\paragraph{LA-0970 [hard]} Explain BentoML from first principles, then show how you would evaluate and operate it in production.
\noindent\textbf{Answer.} Start with the mechanism: BentoML packages model services, dependencies, APIs, and deployment configuration with adaptive batching. Make the tradeoff explicit: Developer workflow is streamlined, while platform integration and runtime tuning still matter. Define workload-specific offline metrics and a latency/cost budget, test important slices, instrument the critical path, and create a rollback criterion. The production review must explicitly guard against this failure: Treating packaging success as load readiness skips concurrency tests.
\par\textit{A strong answer connects mechanism, measurable tradeoffs, failure isolation, observability, and rollback rather than listing product features.}
\paragraph{MCQ-0971 [medium]} Which statement best describes Kubernetes and GPUs?
\begin{enumerate}[label=\Alph*.]
\item Ray Serve composes Python model deployments and scales replicas across a Ray cluster.
\item A registry versions model artifacts, lineage, signatures, evaluation evidence, aliases, and promotion state.
\item TGI is a Hugging Face server for transformer inference with streaming, batching, quantization, and distributed support.
\item Kubernetes schedules containers and device resources while operators add GPU sharing, queues, and autoscaling.
\end{enumerate}
\noindent\textbf{Answer.} Kubernetes schedules containers and device resources while operators add GPU sharing, queues, and autoscaling.
\par\textit{Kubernetes schedules containers and device resources while operators add GPU sharing, queues, and autoscaling. Tradeoff: A common control plane helps governance but default scheduling is not token-aware. Watch for: Requesting whole GPUs for tiny models can strand capacity.}
\paragraph{MCQ-0972 [hard]} What is the central engineering tradeoff of Kubernetes and GPUs?
\begin{enumerate}[label=\Alph*.]
\item Governance and rollback improve but approval gates can slow iteration.
\item A common control plane helps governance but default scheduling is not token-aware.
\item Ecosystem integration is strong, while supported optimization paths vary by model and hardware.
\item Flexible distributed Python pipelines fit multi-model systems but require cluster and object-store discipline.
\end{enumerate}
\noindent\textbf{Answer.} A common control plane helps governance but default scheduling is not token-aware.
\par\textit{Kubernetes schedules containers and device resources while operators add GPU sharing, queues, and autoscaling. Tradeoff: A common control plane helps governance but default scheduling is not token-aware. Watch for: Requesting whole GPUs for tiny models can strand capacity.}
\paragraph{MCQ-0973 [hard]} Which failure mode should an interviewer expect you to identify for Kubernetes and GPUs?
\begin{enumerate}[label=\Alph*.]
\item Enabling quantization without quality and kernel tests can reduce both accuracy and speed.
\item Using mutable stage names without immutable digests makes rollbacks uncertain.
\item Requesting whole GPUs for tiny models can strand capacity.
\item Passing large tensors repeatedly through serialization paths wastes memory and latency.
\end{enumerate}
\noindent\textbf{Answer.} Requesting whole GPUs for tiny models can strand capacity.
\par\textit{Kubernetes schedules containers and device resources while operators add GPU sharing, queues, and autoscaling. Tradeoff: A common control plane helps governance but default scheduling is not token-aware. Watch for: Requesting whole GPUs for tiny models can strand capacity.}
\paragraph{MCQ-0974 [medium]} A design review is specifically concerned with this risk: Requesting whole GPUs for tiny models can strand capacity. Which topic should be reviewed first?
\begin{enumerate}[label=\Alph*.]
\item Text Generation Inference
\item Kubernetes and GPUs
\item Ray Serve
\item Model registry
\end{enumerate}
\noindent\textbf{Answer.} Kubernetes and GPUs
\par\textit{Kubernetes schedules containers and device resources while operators add GPU sharing, queues, and autoscaling. Tradeoff: A common control plane helps governance but default scheduling is not token-aware. Watch for: Requesting whole GPUs for tiny models can strand capacity.}
\paragraph{MCQ-0975 [hard]} Which design note most accurately balances the benefits and costs of Kubernetes and GPUs?
\begin{enumerate}[label=\Alph*.]
\item Ecosystem integration is strong, while supported optimization paths vary by model and hardware.
\item A common control plane helps governance but default scheduling is not token-aware.
\item Flexible distributed Python pipelines fit multi-model systems but require cluster and object-store discipline.
\item Governance and rollback improve but approval gates can slow iteration.
\end{enumerate}
\noindent\textbf{Answer.} A common control plane helps governance but default scheduling is not token-aware.
\par\textit{Kubernetes schedules containers and device resources while operators add GPU sharing, queues, and autoscaling. Tradeoff: A common control plane helps governance but default scheduling is not token-aware. Watch for: Requesting whole GPUs for tiny models can strand capacity.}
\paragraph{MCQ-0976 [medium]} Which explanation would be strongest in a system-design interview about Kubernetes and GPUs?
\begin{enumerate}[label=\Alph*.]
\item TGI is a Hugging Face server for transformer inference with streaming, batching, quantization, and distributed support.
\item A registry versions model artifacts, lineage, signatures, evaluation evidence, aliases, and promotion state.
\item Kubernetes schedules containers and device resources while operators add GPU sharing, queues, and autoscaling.
\item Ray Serve composes Python model deployments and scales replicas across a Ray cluster.
\end{enumerate}
\noindent\textbf{Answer.} Kubernetes schedules containers and device resources while operators add GPU sharing, queues, and autoscaling.
\par\textit{Kubernetes schedules containers and device resources while operators add GPU sharing, queues, and autoscaling. Tradeoff: A common control plane helps governance but default scheduling is not token-aware. Watch for: Requesting whole GPUs for tiny models can strand capacity.}
\paragraph{MCQ-0977 [hard]} Which production warning is most directly associated with Kubernetes and GPUs?
\begin{enumerate}[label=\Alph*.]
\item Passing large tensors repeatedly through serialization paths wastes memory and latency.
\item Requesting whole GPUs for tiny models can strand capacity.
\item Enabling quantization without quality and kernel tests can reduce both accuracy and speed.
\item Using mutable stage names without immutable digests makes rollbacks uncertain.
\end{enumerate}
\noindent\textbf{Answer.} Requesting whole GPUs for tiny models can strand capacity.
\par\textit{Kubernetes schedules containers and device resources while operators add GPU sharing, queues, and autoscaling. Tradeoff: A common control plane helps governance but default scheduling is not token-aware. Watch for: Requesting whole GPUs for tiny models can strand capacity.}
\paragraph{FC-0978 [medium]} Define Kubernetes and GPUs in interview-ready language.
\noindent\textbf{Answer.} Kubernetes schedules containers and device resources while operators add GPU sharing, queues, and autoscaling.
\par\textit{Kubernetes schedules containers and device resources while operators add GPU sharing, queues, and autoscaling.}
\paragraph{FC-0979 [hard]} State the main tradeoff and production pitfall for Kubernetes and GPUs.
\noindent\textbf{Answer.} Tradeoff: A common control plane helps governance but default scheduling is not token-aware. Pitfall: Requesting whole GPUs for tiny models can strand capacity.
\par\textit{Tradeoff: A common control plane helps governance but default scheduling is not token-aware. Pitfall: Requesting whole GPUs for tiny models can strand capacity.}
\paragraph{LA-0980 [hard]} Explain Kubernetes and GPUs from first principles, then show how you would evaluate and operate it in production.
\noindent\textbf{Answer.} Start with the mechanism: Kubernetes schedules containers and device resources while operators add GPU sharing, queues, and autoscaling. Make the tradeoff explicit: A common control plane helps governance but default scheduling is not token-aware. Define workload-specific offline metrics and a latency/cost budget, test important slices, instrument the critical path, and create a rollback criterion. The production review must explicitly guard against this failure: Requesting whole GPUs for tiny models can strand capacity.
\par\textit{A strong answer connects mechanism, measurable tradeoffs, failure isolation, observability, and rollback rather than listing product features.}
\paragraph{MCQ-0981 [medium]} Which statement best describes LLM gateways?
\begin{enumerate}[label=\Alph*.]
\item TensorRT-LLM compiles and serves NVIDIA-optimized LLM graphs and kernels across GPUs.
\item Gateways normalize provider APIs and add routing, budgets, rate limits, fallbacks, caching, and audit policy.
\item Kubernetes schedules containers and device resources while operators add GPU sharing, queues, and autoscaling.
\item SGLang combines a structured generation runtime with optimized model serving, prefix reuse, batching, and parallelism.
\end{enumerate}
\noindent\textbf{Answer.} Gateways normalize provider APIs and add routing, budgets, rate limits, fallbacks, caching, and audit policy.
\par\textit{Gateways normalize provider APIs and add routing, budgets, rate limits, fallbacks, caching, and audit policy. Tradeoff: Central control improves governance but becomes a high-impact availability and security dependency. Watch for: Automatic fallback across non-equivalent models can silently change behavior.}
\paragraph{MCQ-0982 [hard]} What is the central engineering tradeoff of LLM gateways?
\begin{enumerate}[label=\Alph*.]
\item A common control plane helps governance but default scheduling is not token-aware.
\item Central control improves governance but becomes a high-impact availability and security dependency.
\item Tight runtime-serving integration can accelerate agentic workloads but changes application abstractions.
\item Maximum NVIDIA performance comes with build complexity and hardware coupling.
\end{enumerate}
\noindent\textbf{Answer.} Central control improves governance but becomes a high-impact availability and security dependency.
\par\textit{Gateways normalize provider APIs and add routing, budgets, rate limits, fallbacks, caching, and audit policy. Tradeoff: Central control improves governance but becomes a high-impact availability and security dependency. Watch for: Automatic fallback across non-equivalent models can silently change behavior.}
\paragraph{MCQ-0983 [hard]} Which failure mode should an interviewer expect you to identify for LLM gateways?
\begin{enumerate}[label=\Alph*.]
\item Requesting whole GPUs for tiny models can strand capacity.
\item Rebuilding engines for every shape without a deployment cache slows releases.
\item Automatic fallback across non-equivalent models can silently change behavior.
\item Assuming every structured program benefits equally from cache reuse leads to poor capacity models.
\end{enumerate}
\noindent\textbf{Answer.} Automatic fallback across non-equivalent models can silently change behavior.
\par\textit{Gateways normalize provider APIs and add routing, budgets, rate limits, fallbacks, caching, and audit policy. Tradeoff: Central control improves governance but becomes a high-impact availability and security dependency. Watch for: Automatic fallback across non-equivalent models can silently change behavior.}
\paragraph{MCQ-0984 [medium]} A design review is specifically concerned with this risk: Automatic fallback across non-equivalent models can silently change behavior. Which topic should be reviewed first?
\begin{enumerate}[label=\Alph*.]
\item SGLang
\item TensorRT-LLM
\item LLM gateways
\item Kubernetes and GPUs
\end{enumerate}
\noindent\textbf{Answer.} LLM gateways
\par\textit{Gateways normalize provider APIs and add routing, budgets, rate limits, fallbacks, caching, and audit policy. Tradeoff: Central control improves governance but becomes a high-impact availability and security dependency. Watch for: Automatic fallback across non-equivalent models can silently change behavior.}
\paragraph{MCQ-0985 [hard]} Which design note most accurately balances the benefits and costs of LLM gateways?
\begin{enumerate}[label=\Alph*.]
\item Maximum NVIDIA performance comes with build complexity and hardware coupling.
\item A common control plane helps governance but default scheduling is not token-aware.
\item Central control improves governance but becomes a high-impact availability and security dependency.
\item Tight runtime-serving integration can accelerate agentic workloads but changes application abstractions.
\end{enumerate}
\noindent\textbf{Answer.} Central control improves governance but becomes a high-impact availability and security dependency.
\par\textit{Gateways normalize provider APIs and add routing, budgets, rate limits, fallbacks, caching, and audit policy. Tradeoff: Central control improves governance but becomes a high-impact availability and security dependency. Watch for: Automatic fallback across non-equivalent models can silently change behavior.}
\paragraph{MCQ-0986 [medium]} Which explanation would be strongest in a system-design interview about LLM gateways?
\begin{enumerate}[label=\Alph*.]
\item Gateways normalize provider APIs and add routing, budgets, rate limits, fallbacks, caching, and audit policy.
\item SGLang combines a structured generation runtime with optimized model serving, prefix reuse, batching, and parallelism.
\item Kubernetes schedules containers and device resources while operators add GPU sharing, queues, and autoscaling.
\item TensorRT-LLM compiles and serves NVIDIA-optimized LLM graphs and kernels across GPUs.
\end{enumerate}
\noindent\textbf{Answer.} Gateways normalize provider APIs and add routing, budgets, rate limits, fallbacks, caching, and audit policy.
\par\textit{Gateways normalize provider APIs and add routing, budgets, rate limits, fallbacks, caching, and audit policy. Tradeoff: Central control improves governance but becomes a high-impact availability and security dependency. Watch for: Automatic fallback across non-equivalent models can silently change behavior.}
\paragraph{MCQ-0987 [hard]} Which production warning is most directly associated with LLM gateways?
\begin{enumerate}[label=\Alph*.]
\item Automatic fallback across non-equivalent models can silently change behavior.
\item Assuming every structured program benefits equally from cache reuse leads to poor capacity models.
\item Requesting whole GPUs for tiny models can strand capacity.
\item Rebuilding engines for every shape without a deployment cache slows releases.
\end{enumerate}
\noindent\textbf{Answer.} Automatic fallback across non-equivalent models can silently change behavior.
\par\textit{Gateways normalize provider APIs and add routing, budgets, rate limits, fallbacks, caching, and audit policy. Tradeoff: Central control improves governance but becomes a high-impact availability and security dependency. Watch for: Automatic fallback across non-equivalent models can silently change behavior.}
\paragraph{FC-0988 [medium]} Define LLM gateways in interview-ready language.
\noindent\textbf{Answer.} Gateways normalize provider APIs and add routing, budgets, rate limits, fallbacks, caching, and audit policy.
\par\textit{Gateways normalize provider APIs and add routing, budgets, rate limits, fallbacks, caching, and audit policy.}
\paragraph{FC-0989 [hard]} State the main tradeoff and production pitfall for LLM gateways.
\noindent\textbf{Answer.} Tradeoff: Central control improves governance but becomes a high-impact availability and security dependency. Pitfall: Automatic fallback across non-equivalent models can silently change behavior.
\par\textit{Tradeoff: Central control improves governance but becomes a high-impact availability and security dependency. Pitfall: Automatic fallback across non-equivalent models can silently change behavior.}
\paragraph{LA-0990 [hard]} Explain LLM gateways from first principles, then show how you would evaluate and operate it in production.
\noindent\textbf{Answer.} Start with the mechanism: Gateways normalize provider APIs and add routing, budgets, rate limits, fallbacks, caching, and audit policy. Make the tradeoff explicit: Central control improves governance but becomes a high-impact availability and security dependency. Define workload-specific offline metrics and a latency/cost budget, test important slices, instrument the critical path, and create a rollback criterion. The production review must explicitly guard against this failure: Automatic fallback across non-equivalent models can silently change behavior.
\par\textit{A strong answer connects mechanism, measurable tradeoffs, failure isolation, observability, and rollback rather than listing product features.}
\paragraph{MCQ-0991 [medium]} Which statement best describes Canary and shadow releases?
\begin{enumerate}[label=\Alph*.]
\item BentoML packages model services, dependencies, APIs, and deployment configuration with adaptive batching.
\item llama.cpp runs quantized models efficiently on CPUs and varied local accelerators through a portable C++ stack.
\item Canary traffic exposes a new version to a small cohort; shadowing duplicates requests without serving the new output.
\item Ray Serve composes Python model deployments and scales replicas across a Ray cluster.
\end{enumerate}
\noindent\textbf{Answer.} Canary traffic exposes a new version to a small cohort; shadowing duplicates requests without serving the new output.
\par\textit{Canary traffic exposes a new version to a small cohort; shadowing duplicates requests without serving the new output. Tradeoff: Risk falls, while cost, comparison logic, and privacy obligations grow. Watch for: Comparing canaries with different traffic mixes confounds the result.}
\paragraph{MCQ-0992 [hard]} What is the central engineering tradeoff of Canary and shadow releases?
\begin{enumerate}[label=\Alph*.]
\item Flexible distributed Python pipelines fit multi-model systems but require cluster and object-store discipline.
\item Edge privacy and low dependency count trade off against datacenter throughput and model size.
\item Risk falls, while cost, comparison logic, and privacy obligations grow.
\item Developer workflow is streamlined, while platform integration and runtime tuning still matter.
\end{enumerate}
\noindent\textbf{Answer.} Risk falls, while cost, comparison logic, and privacy obligations grow.
\par\textit{Canary traffic exposes a new version to a small cohort; shadowing duplicates requests without serving the new output. Tradeoff: Risk falls, while cost, comparison logic, and privacy obligations grow. Watch for: Comparing canaries with different traffic mixes confounds the result.}
\paragraph{MCQ-0993 [hard]} Which failure mode should an interviewer expect you to identify for Canary and shadow releases?
\begin{enumerate}[label=\Alph*.]
\item Passing large tensors repeatedly through serialization paths wastes memory and latency.
\item Comparing canaries with different traffic mixes confounds the result.
\item Selecting a quantization by file size alone ignores accuracy and hardware kernels.
\item Treating packaging success as load readiness skips concurrency tests.
\end{enumerate}
\noindent\textbf{Answer.} Comparing canaries with different traffic mixes confounds the result.
\par\textit{Canary traffic exposes a new version to a small cohort; shadowing duplicates requests without serving the new output. Tradeoff: Risk falls, while cost, comparison logic, and privacy obligations grow. Watch for: Comparing canaries with different traffic mixes confounds the result.}
\paragraph{MCQ-0994 [medium]} A design review is specifically concerned with this risk: Comparing canaries with different traffic mixes confounds the result. Which topic should be reviewed first?
\begin{enumerate}[label=\Alph*.]
\item Ray Serve
\item Canary and shadow releases
\item BentoML
\item llama.cpp
\end{enumerate}
\noindent\textbf{Answer.} Canary and shadow releases
\par\textit{Canary traffic exposes a new version to a small cohort; shadowing duplicates requests without serving the new output. Tradeoff: Risk falls, while cost, comparison logic, and privacy obligations grow. Watch for: Comparing canaries with different traffic mixes confounds the result.}
\paragraph{MCQ-0995 [hard]} Which design note most accurately balances the benefits and costs of Canary and shadow releases?
\begin{enumerate}[label=\Alph*.]
\item Edge privacy and low dependency count trade off against datacenter throughput and model size.
\item Developer workflow is streamlined, while platform integration and runtime tuning still matter.
\item Flexible distributed Python pipelines fit multi-model systems but require cluster and object-store discipline.
\item Risk falls, while cost, comparison logic, and privacy obligations grow.
\end{enumerate}
\noindent\textbf{Answer.} Risk falls, while cost, comparison logic, and privacy obligations grow.
\par\textit{Canary traffic exposes a new version to a small cohort; shadowing duplicates requests without serving the new output. Tradeoff: Risk falls, while cost, comparison logic, and privacy obligations grow. Watch for: Comparing canaries with different traffic mixes confounds the result.}
\paragraph{MCQ-0996 [medium]} Which explanation would be strongest in a system-design interview about Canary and shadow releases?
\begin{enumerate}[label=\Alph*.]
\item Ray Serve composes Python model deployments and scales replicas across a Ray cluster.
\item llama.cpp runs quantized models efficiently on CPUs and varied local accelerators through a portable C++ stack.
\item Canary traffic exposes a new version to a small cohort; shadowing duplicates requests without serving the new output.
\item BentoML packages model services, dependencies, APIs, and deployment configuration with adaptive batching.
\end{enumerate}
\noindent\textbf{Answer.} Canary traffic exposes a new version to a small cohort; shadowing duplicates requests without serving the new output.
\par\textit{Canary traffic exposes a new version to a small cohort; shadowing duplicates requests without serving the new output. Tradeoff: Risk falls, while cost, comparison logic, and privacy obligations grow. Watch for: Comparing canaries with different traffic mixes confounds the result.}
\paragraph{MCQ-0997 [hard]} Which production warning is most directly associated with Canary and shadow releases?
\begin{enumerate}[label=\Alph*.]
\item Selecting a quantization by file size alone ignores accuracy and hardware kernels.
\item Passing large tensors repeatedly through serialization paths wastes memory and latency.
\item Comparing canaries with different traffic mixes confounds the result.
\item Treating packaging success as load readiness skips concurrency tests.
\end{enumerate}
\noindent\textbf{Answer.} Comparing canaries with different traffic mixes confounds the result.
\par\textit{Canary traffic exposes a new version to a small cohort; shadowing duplicates requests without serving the new output. Tradeoff: Risk falls, while cost, comparison logic, and privacy obligations grow. Watch for: Comparing canaries with different traffic mixes confounds the result.}
\paragraph{FC-0998 [medium]} Define Canary and shadow releases in interview-ready language.
\noindent\textbf{Answer.} Canary traffic exposes a new version to a small cohort; shadowing duplicates requests without serving the new output.
\par\textit{Canary traffic exposes a new version to a small cohort; shadowing duplicates requests without serving the new output.}
\paragraph{FC-0999 [hard]} State the main tradeoff and production pitfall for Canary and shadow releases.
\noindent\textbf{Answer.} Tradeoff: Risk falls, while cost, comparison logic, and privacy obligations grow. Pitfall: Comparing canaries with different traffic mixes confounds the result.
\par\textit{Tradeoff: Risk falls, while cost, comparison logic, and privacy obligations grow. Pitfall: Comparing canaries with different traffic mixes confounds the result.}
\paragraph{LA-1000 [hard]} Explain Canary and shadow releases from first principles, then show how you would evaluate and operate it in production.
\noindent\textbf{Answer.} Start with the mechanism: Canary traffic exposes a new version to a small cohort; shadowing duplicates requests without serving the new output. Make the tradeoff explicit: Risk falls, while cost, comparison logic, and privacy obligations grow. Define workload-specific offline metrics and a latency/cost budget, test important slices, instrument the critical path, and create a rollback criterion. The production review must explicitly guard against this failure: Comparing canaries with different traffic mixes confounds the result.
\par\textit{A strong answer connects mechanism, measurable tradeoffs, failure isolation, observability, and rollback rather than listing product features.}
\paragraph{MCQ-1001 [medium]} Which statement best describes Model registry?
\begin{enumerate}[label=\Alph*.]
\item A registry versions model artifacts, lineage, signatures, evaluation evidence, aliases, and promotion state.
\item Ray Serve composes Python model deployments and scales replicas across a Ray cluster.
\item Kubernetes schedules containers and device resources while operators add GPU sharing, queues, and autoscaling.
\item TensorRT-LLM compiles and serves NVIDIA-optimized LLM graphs and kernels across GPUs.
\end{enumerate}
\noindent\textbf{Answer.} A registry versions model artifacts, lineage, signatures, evaluation evidence, aliases, and promotion state.
\par\textit{A registry versions model artifacts, lineage, signatures, evaluation evidence, aliases, and promotion state. Tradeoff: Governance and rollback improve but approval gates can slow iteration. Watch for: Using mutable stage names without immutable digests makes rollbacks uncertain.}
\paragraph{MCQ-1002 [hard]} What is the central engineering tradeoff of Model registry?
\begin{enumerate}[label=\Alph*.]
\item Maximum NVIDIA performance comes with build complexity and hardware coupling.
\item Governance and rollback improve but approval gates can slow iteration.
\item A common control plane helps governance but default scheduling is not token-aware.
\item Flexible distributed Python pipelines fit multi-model systems but require cluster and object-store discipline.
\end{enumerate}
\noindent\textbf{Answer.} Governance and rollback improve but approval gates can slow iteration.
\par\textit{A registry versions model artifacts, lineage, signatures, evaluation evidence, aliases, and promotion state. Tradeoff: Governance and rollback improve but approval gates can slow iteration. Watch for: Using mutable stage names without immutable digests makes rollbacks uncertain.}
\paragraph{MCQ-1003 [hard]} Which failure mode should an interviewer expect you to identify for Model registry?
\begin{enumerate}[label=\Alph*.]
\item Requesting whole GPUs for tiny models can strand capacity.
\item Passing large tensors repeatedly through serialization paths wastes memory and latency.
\item Using mutable stage names without immutable digests makes rollbacks uncertain.
\item Rebuilding engines for every shape without a deployment cache slows releases.
\end{enumerate}
\noindent\textbf{Answer.} Using mutable stage names without immutable digests makes rollbacks uncertain.
\par\textit{A registry versions model artifacts, lineage, signatures, evaluation evidence, aliases, and promotion state. Tradeoff: Governance and rollback improve but approval gates can slow iteration. Watch for: Using mutable stage names without immutable digests makes rollbacks uncertain.}
\paragraph{MCQ-1004 [medium]} A design review is specifically concerned with this risk: Using mutable stage names without immutable digests makes rollbacks uncertain. Which topic should be reviewed first?
\begin{enumerate}[label=\Alph*.]
\item Kubernetes and GPUs
\item TensorRT-LLM
\item Model registry
\item Ray Serve
\end{enumerate}
\noindent\textbf{Answer.} Model registry
\par\textit{A registry versions model artifacts, lineage, signatures, evaluation evidence, aliases, and promotion state. Tradeoff: Governance and rollback improve but approval gates can slow iteration. Watch for: Using mutable stage names without immutable digests makes rollbacks uncertain.}
\paragraph{MCQ-1005 [hard]} Which design note most accurately balances the benefits and costs of Model registry?
\begin{enumerate}[label=\Alph*.]
\item Flexible distributed Python pipelines fit multi-model systems but require cluster and object-store discipline.
\item Maximum NVIDIA performance comes with build complexity and hardware coupling.
\item A common control plane helps governance but default scheduling is not token-aware.
\item Governance and rollback improve but approval gates can slow iteration.
\end{enumerate}
\noindent\textbf{Answer.} Governance and rollback improve but approval gates can slow iteration.
\par\textit{A registry versions model artifacts, lineage, signatures, evaluation evidence, aliases, and promotion state. Tradeoff: Governance and rollback improve but approval gates can slow iteration. Watch for: Using mutable stage names without immutable digests makes rollbacks uncertain.}
\paragraph{MCQ-1006 [medium]} Which explanation would be strongest in a system-design interview about Model registry?
\begin{enumerate}[label=\Alph*.]
\item TensorRT-LLM compiles and serves NVIDIA-optimized LLM graphs and kernels across GPUs.
\item Kubernetes schedules containers and device resources while operators add GPU sharing, queues, and autoscaling.
\item Ray Serve composes Python model deployments and scales replicas across a Ray cluster.
\item A registry versions model artifacts, lineage, signatures, evaluation evidence, aliases, and promotion state.
\end{enumerate}
\noindent\textbf{Answer.} A registry versions model artifacts, lineage, signatures, evaluation evidence, aliases, and promotion state.
\par\textit{A registry versions model artifacts, lineage, signatures, evaluation evidence, aliases, and promotion state. Tradeoff: Governance and rollback improve but approval gates can slow iteration. Watch for: Using mutable stage names without immutable digests makes rollbacks uncertain.}
\paragraph{MCQ-1007 [hard]} Which production warning is most directly associated with Model registry?
\begin{enumerate}[label=\Alph*.]
\item Requesting whole GPUs for tiny models can strand capacity.
\item Rebuilding engines for every shape without a deployment cache slows releases.
\item Using mutable stage names without immutable digests makes rollbacks uncertain.
\item Passing large tensors repeatedly through serialization paths wastes memory and latency.
\end{enumerate}
\noindent\textbf{Answer.} Using mutable stage names without immutable digests makes rollbacks uncertain.
\par\textit{A registry versions model artifacts, lineage, signatures, evaluation evidence, aliases, and promotion state. Tradeoff: Governance and rollback improve but approval gates can slow iteration. Watch for: Using mutable stage names without immutable digests makes rollbacks uncertain.}
\paragraph{FC-1008 [medium]} Define Model registry in interview-ready language.
\noindent\textbf{Answer.} A registry versions model artifacts, lineage, signatures, evaluation evidence, aliases, and promotion state.
\par\textit{A registry versions model artifacts, lineage, signatures, evaluation evidence, aliases, and promotion state.}
\paragraph{FC-1009 [hard]} State the main tradeoff and production pitfall for Model registry.
\noindent\textbf{Answer.} Tradeoff: Governance and rollback improve but approval gates can slow iteration. Pitfall: Using mutable stage names without immutable digests makes rollbacks uncertain.
\par\textit{Tradeoff: Governance and rollback improve but approval gates can slow iteration. Pitfall: Using mutable stage names without immutable digests makes rollbacks uncertain.}
\paragraph{LA-1010 [hard]} Explain Model registry from first principles, then show how you would evaluate and operate it in production.
\noindent\textbf{Answer.} Start with the mechanism: A registry versions model artifacts, lineage, signatures, evaluation evidence, aliases, and promotion state. Make the tradeoff explicit: Governance and rollback improve but approval gates can slow iteration. Define workload-specific offline metrics and a latency/cost budget, test important slices, instrument the critical path, and create a rollback criterion. The production review must explicitly guard against this failure: Using mutable stage names without immutable digests makes rollbacks uncertain.
\par\textit{A strong answer connects mechanism, measurable tradeoffs, failure isolation, observability, and rollback rather than listing product features.}
\paragraph{MCQ-1011 [medium]} Which statement best describes Autoscaling for LLMs?
\begin{enumerate}[label=\Alph*.]
\item TensorRT-LLM compiles and serves NVIDIA-optimized LLM graphs and kernels across GPUs.
\item LLM autoscaling should consider queue depth, time-to-first-token, tokens per second, memory, and startup time.
\item vLLM is a high-throughput inference server using PagedAttention, continuous batching, parallelism, and OpenAI-compatible APIs.
\item BentoML packages model services, dependencies, APIs, and deployment configuration with adaptive batching.
\end{enumerate}
\noindent\textbf{Answer.} LLM autoscaling should consider queue depth, time-to-first-token, tokens per second, memory, and startup time.
\par\textit{LLM autoscaling should consider queue depth, time-to-first-token, tokens per second, memory, and startup time. Tradeoff: Extra replicas reduce latency but GPUs are costly and slow to warm. Watch for: CPU-based autoscaling reacts late or incorrectly to GPU saturation.}
\paragraph{MCQ-1012 [hard]} What is the central engineering tradeoff of Autoscaling for LLMs?
\begin{enumerate}[label=\Alph*.]
\item Extra replicas reduce latency but GPUs are costly and slow to warm.
\item Throughput and memory utilization improve, while model and kernel compatibility must be verified.
\item Maximum NVIDIA performance comes with build complexity and hardware coupling.
\item Developer workflow is streamlined, while platform integration and runtime tuning still matter.
\end{enumerate}
\noindent\textbf{Answer.} Extra replicas reduce latency but GPUs are costly and slow to warm.
\par\textit{LLM autoscaling should consider queue depth, time-to-first-token, tokens per second, memory, and startup time. Tradeoff: Extra replicas reduce latency but GPUs are costly and slow to warm. Watch for: CPU-based autoscaling reacts late or incorrectly to GPU saturation.}
\paragraph{MCQ-1013 [hard]} Which failure mode should an interviewer expect you to identify for Autoscaling for LLMs?
\begin{enumerate}[label=\Alph*.]
\item Rebuilding engines for every shape without a deployment cache slows releases.
\item Benchmarking only one concurrency level hides latency-throughput collapse points.
\item Treating packaging success as load readiness skips concurrency tests.
\item CPU-based autoscaling reacts late or incorrectly to GPU saturation.
\end{enumerate}
\noindent\textbf{Answer.} CPU-based autoscaling reacts late or incorrectly to GPU saturation.
\par\textit{LLM autoscaling should consider queue depth, time-to-first-token, tokens per second, memory, and startup time. Tradeoff: Extra replicas reduce latency but GPUs are costly and slow to warm. Watch for: CPU-based autoscaling reacts late or incorrectly to GPU saturation.}
\paragraph{MCQ-1014 [medium]} A design review is specifically concerned with this risk: CPU-based autoscaling reacts late or incorrectly to GPU saturation. Which topic should be reviewed first?
\begin{enumerate}[label=\Alph*.]
\item TensorRT-LLM
\item vLLM
\item BentoML
\item Autoscaling for LLMs
\end{enumerate}
\noindent\textbf{Answer.} Autoscaling for LLMs
\par\textit{LLM autoscaling should consider queue depth, time-to-first-token, tokens per second, memory, and startup time. Tradeoff: Extra replicas reduce latency but GPUs are costly and slow to warm. Watch for: CPU-based autoscaling reacts late or incorrectly to GPU saturation.}
\paragraph{MCQ-1015 [hard]} Which design note most accurately balances the benefits and costs of Autoscaling for LLMs?
\begin{enumerate}[label=\Alph*.]
\item Throughput and memory utilization improve, while model and kernel compatibility must be verified.
\item Developer workflow is streamlined, while platform integration and runtime tuning still matter.
\item Extra replicas reduce latency but GPUs are costly and slow to warm.
\item Maximum NVIDIA performance comes with build complexity and hardware coupling.
\end{enumerate}
\noindent\textbf{Answer.} Extra replicas reduce latency but GPUs are costly and slow to warm.
\par\textit{LLM autoscaling should consider queue depth, time-to-first-token, tokens per second, memory, and startup time. Tradeoff: Extra replicas reduce latency but GPUs are costly and slow to warm. Watch for: CPU-based autoscaling reacts late or incorrectly to GPU saturation.}
\paragraph{MCQ-1016 [medium]} Which explanation would be strongest in a system-design interview about Autoscaling for LLMs?
\begin{enumerate}[label=\Alph*.]
\item vLLM is a high-throughput inference server using PagedAttention, continuous batching, parallelism, and OpenAI-compatible APIs.
\item TensorRT-LLM compiles and serves NVIDIA-optimized LLM graphs and kernels across GPUs.
\item LLM autoscaling should consider queue depth, time-to-first-token, tokens per second, memory, and startup time.
\item BentoML packages model services, dependencies, APIs, and deployment configuration with adaptive batching.
\end{enumerate}
\noindent\textbf{Answer.} LLM autoscaling should consider queue depth, time-to-first-token, tokens per second, memory, and startup time.
\par\textit{LLM autoscaling should consider queue depth, time-to-first-token, tokens per second, memory, and startup time. Tradeoff: Extra replicas reduce latency but GPUs are costly and slow to warm. Watch for: CPU-based autoscaling reacts late or incorrectly to GPU saturation.}
\paragraph{MCQ-1017 [hard]} Which production warning is most directly associated with Autoscaling for LLMs?
\begin{enumerate}[label=\Alph*.]
\item Benchmarking only one concurrency level hides latency-throughput collapse points.
\item Rebuilding engines for every shape without a deployment cache slows releases.
\item CPU-based autoscaling reacts late or incorrectly to GPU saturation.
\item Treating packaging success as load readiness skips concurrency tests.
\end{enumerate}
\noindent\textbf{Answer.} CPU-based autoscaling reacts late or incorrectly to GPU saturation.
\par\textit{LLM autoscaling should consider queue depth, time-to-first-token, tokens per second, memory, and startup time. Tradeoff: Extra replicas reduce latency but GPUs are costly and slow to warm. Watch for: CPU-based autoscaling reacts late or incorrectly to GPU saturation.}
\paragraph{FC-1018 [medium]} Define Autoscaling for LLMs in interview-ready language.
\noindent\textbf{Answer.} LLM autoscaling should consider queue depth, time-to-first-token, tokens per second, memory, and startup time.
\par\textit{LLM autoscaling should consider queue depth, time-to-first-token, tokens per second, memory, and startup time.}
\paragraph{FC-1019 [hard]} State the main tradeoff and production pitfall for Autoscaling for LLMs.
\noindent\textbf{Answer.} Tradeoff: Extra replicas reduce latency but GPUs are costly and slow to warm. Pitfall: CPU-based autoscaling reacts late or incorrectly to GPU saturation.
\par\textit{Tradeoff: Extra replicas reduce latency but GPUs are costly and slow to warm. Pitfall: CPU-based autoscaling reacts late or incorrectly to GPU saturation.}
\paragraph{LA-1020 [hard]} Explain Autoscaling for LLMs from first principles, then show how you would evaluate and operate it in production.
\noindent\textbf{Answer.} Start with the mechanism: LLM autoscaling should consider queue depth, time-to-first-token, tokens per second, memory, and startup time. Make the tradeoff explicit: Extra replicas reduce latency but GPUs are costly and slow to warm. Define workload-specific offline metrics and a latency/cost budget, test important slices, instrument the critical path, and create a rollback criterion. The production review must explicitly guard against this failure: CPU-based autoscaling reacts late or incorrectly to GPU saturation.
\par\textit{A strong answer connects mechanism, measurable tradeoffs, failure isolation, observability, and rollback rather than listing product features.}
\subsection{Data and ML Lifecycle}
\paragraph{MCQ-1021 [medium]} Which statement best describes DVC?
\begin{enumerate}[label=\Alph*.]
\item DVC versions large data and model pointers with Git and defines reproducible pipeline stages and experiments.
\item Tests cover schema, nulls, ranges, uniqueness, distribution, freshness, and cross-table invariants.
\item Lockfiles, containers, hardware metadata, seeds, and deterministic settings constrain run variability.
\item Runs capture code, parameters, data, metrics, artifacts, environment, and parent relationships.
\end{enumerate}
\noindent\textbf{Answer.} DVC versions large data and model pointers with Git and defines reproducible pipeline stages and experiments.
\par\textit{DVC versions large data and model pointers with Git and defines reproducible pipeline stages and experiments. Tradeoff: Familiar Git workflows improve reproducibility, while remote storage and cache discipline are required. Watch for: Committing only code while silently replacing remote data breaks lineage.}
\paragraph{MCQ-1022 [hard]} What is the central engineering tradeoff of DVC?
\begin{enumerate}[label=\Alph*.]
\item Early detection prevents downstream damage, while brittle thresholds create alert fatigue.
\item Comparison and reproducibility improve but inconsistent naming creates a junk drawer.
\item Reproducibility improves at the expense of build maintenance and sometimes performance.
\item Familiar Git workflows improve reproducibility, while remote storage and cache discipline are required.
\end{enumerate}
\noindent\textbf{Answer.} Familiar Git workflows improve reproducibility, while remote storage and cache discipline are required.
\par\textit{DVC versions large data and model pointers with Git and defines reproducible pipeline stages and experiments. Tradeoff: Familiar Git workflows improve reproducibility, while remote storage and cache discipline are required. Watch for: Committing only code while silently replacing remote data breaks lineage.}
\paragraph{MCQ-1023 [hard]} Which failure mode should an interviewer expect you to identify for DVC?
\begin{enumerate}[label=\Alph*.]
\item A random seed alone cannot guarantee deterministic GPU execution.
\item Committing only code while silently replacing remote data breaks lineage.
\item Logging a metric without dataset and code versions makes it uninterpretable.
\item Checking only aggregate distributions misses failures in important slices.
\end{enumerate}
\noindent\textbf{Answer.} Committing only code while silently replacing remote data breaks lineage.
\par\textit{DVC versions large data and model pointers with Git and defines reproducible pipeline stages and experiments. Tradeoff: Familiar Git workflows improve reproducibility, while remote storage and cache discipline are required. Watch for: Committing only code while silently replacing remote data breaks lineage.}
\paragraph{MCQ-1024 [medium]} A design review is specifically concerned with this risk: Committing only code while silently replacing remote data breaks lineage. Which topic should be reviewed first?
\begin{enumerate}[label=\Alph*.]
\item Experiment tracking
\item DVC
\item Data quality testing
\item Reproducible environments
\end{enumerate}
\noindent\textbf{Answer.} DVC
\par\textit{DVC versions large data and model pointers with Git and defines reproducible pipeline stages and experiments. Tradeoff: Familiar Git workflows improve reproducibility, while remote storage and cache discipline are required. Watch for: Committing only code while silently replacing remote data breaks lineage.}
\paragraph{MCQ-1025 [hard]} Which design note most accurately balances the benefits and costs of DVC?
\begin{enumerate}[label=\Alph*.]
\item Comparison and reproducibility improve but inconsistent naming creates a junk drawer.
\item Early detection prevents downstream damage, while brittle thresholds create alert fatigue.
\item Reproducibility improves at the expense of build maintenance and sometimes performance.
\item Familiar Git workflows improve reproducibility, while remote storage and cache discipline are required.
\end{enumerate}
\noindent\textbf{Answer.} Familiar Git workflows improve reproducibility, while remote storage and cache discipline are required.
\par\textit{DVC versions large data and model pointers with Git and defines reproducible pipeline stages and experiments. Tradeoff: Familiar Git workflows improve reproducibility, while remote storage and cache discipline are required. Watch for: Committing only code while silently replacing remote data breaks lineage.}
\paragraph{MCQ-1026 [medium]} Which explanation would be strongest in a system-design interview about DVC?
\begin{enumerate}[label=\Alph*.]
\item DVC versions large data and model pointers with Git and defines reproducible pipeline stages and experiments.
\item Runs capture code, parameters, data, metrics, artifacts, environment, and parent relationships.
\item Tests cover schema, nulls, ranges, uniqueness, distribution, freshness, and cross-table invariants.
\item Lockfiles, containers, hardware metadata, seeds, and deterministic settings constrain run variability.
\end{enumerate}
\noindent\textbf{Answer.} DVC versions large data and model pointers with Git and defines reproducible pipeline stages and experiments.
\par\textit{DVC versions large data and model pointers with Git and defines reproducible pipeline stages and experiments. Tradeoff: Familiar Git workflows improve reproducibility, while remote storage and cache discipline are required. Watch for: Committing only code while silently replacing remote data breaks lineage.}
\paragraph{MCQ-1027 [hard]} Which production warning is most directly associated with DVC?
\begin{enumerate}[label=\Alph*.]
\item Checking only aggregate distributions misses failures in important slices.
\item A random seed alone cannot guarantee deterministic GPU execution.
\item Logging a metric without dataset and code versions makes it uninterpretable.
\item Committing only code while silently replacing remote data breaks lineage.
\end{enumerate}
\noindent\textbf{Answer.} Committing only code while silently replacing remote data breaks lineage.
\par\textit{DVC versions large data and model pointers with Git and defines reproducible pipeline stages and experiments. Tradeoff: Familiar Git workflows improve reproducibility, while remote storage and cache discipline are required. Watch for: Committing only code while silently replacing remote data breaks lineage.}
\paragraph{FC-1028 [medium]} Define DVC in interview-ready language.
\noindent\textbf{Answer.} DVC versions large data and model pointers with Git and defines reproducible pipeline stages and experiments.
\par\textit{DVC versions large data and model pointers with Git and defines reproducible pipeline stages and experiments.}
\paragraph{FC-1029 [hard]} State the main tradeoff and production pitfall for DVC.
\noindent\textbf{Answer.} Tradeoff: Familiar Git workflows improve reproducibility, while remote storage and cache discipline are required. Pitfall: Committing only code while silently replacing remote data breaks lineage.
\par\textit{Tradeoff: Familiar Git workflows improve reproducibility, while remote storage and cache discipline are required. Pitfall: Committing only code while silently replacing remote data breaks lineage.}
\paragraph{LA-1030 [hard]} Explain DVC from first principles, then show how you would evaluate and operate it in production.
\noindent\textbf{Answer.} Start with the mechanism: DVC versions large data and model pointers with Git and defines reproducible pipeline stages and experiments. Make the tradeoff explicit: Familiar Git workflows improve reproducibility, while remote storage and cache discipline are required. Define workload-specific offline metrics and a latency/cost budget, test important slices, instrument the critical path, and create a rollback criterion. The production review must explicitly guard against this failure: Committing only code while silently replacing remote data breaks lineage.
\par\textit{A strong answer connects mechanism, measurable tradeoffs, failure isolation, observability, and rollback rather than listing product features.}
\paragraph{MCQ-1031 [medium]} Which statement best describes lakeFS?
\begin{enumerate}[label=\Alph*.]
\item lakeFS adds Git-like branches, commits, merges, and hooks over object-storage data lakes.
\item Feature stores define reusable features and connect historical offline data with low-latency online serving.
\item Open table formats add schemas, snapshots, partition evolution, and transactional metadata over object storage.
\item Historical joins select only feature values known at each label timestamp.
\end{enumerate}
\noindent\textbf{Answer.} lakeFS adds Git-like branches, commits, merges, and hooks over object-storage data lakes.
\par\textit{lakeFS adds Git-like branches, commits, merges, and hooks over object-storage data lakes. Tradeoff: Atomic data-lake changes enable testing, with server and metadata operations to manage. Watch for: Treating object versioning alone as branch-level data semantics is insufficient.}
\paragraph{MCQ-1032 [hard]} What is the central engineering tradeoff of lakeFS?
\begin{enumerate}[label=\Alph*.]
\item Training-serving consistency improves but adds infrastructure and ownership requirements.
\item Leakage is prevented at the cost of temporal keys and more expensive joins.
\item Reliable analytics improve, while compaction and catalog operations become essential.
\item Atomic data-lake changes enable testing, with server and metadata operations to manage.
\end{enumerate}
\noindent\textbf{Answer.} Atomic data-lake changes enable testing, with server and metadata operations to manage.
\par\textit{lakeFS adds Git-like branches, commits, merges, and hooks over object-storage data lakes. Tradeoff: Atomic data-lake changes enable testing, with server and metadata operations to manage. Watch for: Treating object versioning alone as branch-level data semantics is insufficient.}
\paragraph{MCQ-1033 [hard]} Which failure mode should an interviewer expect you to identify for lakeFS?
\begin{enumerate}[label=\Alph*.]
\item Excess small files destroy scan performance despite correct table semantics.
\item Joining the latest feature value into old examples creates unrealistically strong training results.
\item Treating object versioning alone as branch-level data semantics is insufficient.
\item Recomputing features differently online creates skew.
\end{enumerate}
\noindent\textbf{Answer.} Treating object versioning alone as branch-level data semantics is insufficient.
\par\textit{lakeFS adds Git-like branches, commits, merges, and hooks over object-storage data lakes. Tradeoff: Atomic data-lake changes enable testing, with server and metadata operations to manage. Watch for: Treating object versioning alone as branch-level data semantics is insufficient.}
\paragraph{MCQ-1034 [medium]} A design review is specifically concerned with this risk: Treating object versioning alone as branch-level data semantics is insufficient. Which topic should be reviewed first?
\begin{enumerate}[label=\Alph*.]
\item Feature stores
\item Point-in-time correctness
\item lakeFS
\item Delta Lake and Iceberg
\end{enumerate}
\noindent\textbf{Answer.} lakeFS
\par\textit{lakeFS adds Git-like branches, commits, merges, and hooks over object-storage data lakes. Tradeoff: Atomic data-lake changes enable testing, with server and metadata operations to manage. Watch for: Treating object versioning alone as branch-level data semantics is insufficient.}
\paragraph{MCQ-1035 [hard]} Which design note most accurately balances the benefits and costs of lakeFS?
\begin{enumerate}[label=\Alph*.]
\item Leakage is prevented at the cost of temporal keys and more expensive joins.
\item Atomic data-lake changes enable testing, with server and metadata operations to manage.
\item Training-serving consistency improves but adds infrastructure and ownership requirements.
\item Reliable analytics improve, while compaction and catalog operations become essential.
\end{enumerate}
\noindent\textbf{Answer.} Atomic data-lake changes enable testing, with server and metadata operations to manage.
\par\textit{lakeFS adds Git-like branches, commits, merges, and hooks over object-storage data lakes. Tradeoff: Atomic data-lake changes enable testing, with server and metadata operations to manage. Watch for: Treating object versioning alone as branch-level data semantics is insufficient.}
\paragraph{MCQ-1036 [medium]} Which explanation would be strongest in a system-design interview about lakeFS?
\begin{enumerate}[label=\Alph*.]
\item Open table formats add schemas, snapshots, partition evolution, and transactional metadata over object storage.
\item lakeFS adds Git-like branches, commits, merges, and hooks over object-storage data lakes.
\item Feature stores define reusable features and connect historical offline data with low-latency online serving.
\item Historical joins select only feature values known at each label timestamp.
\end{enumerate}
\noindent\textbf{Answer.} lakeFS adds Git-like branches, commits, merges, and hooks over object-storage data lakes.
\par\textit{lakeFS adds Git-like branches, commits, merges, and hooks over object-storage data lakes. Tradeoff: Atomic data-lake changes enable testing, with server and metadata operations to manage. Watch for: Treating object versioning alone as branch-level data semantics is insufficient.}
\paragraph{MCQ-1037 [hard]} Which production warning is most directly associated with lakeFS?
\begin{enumerate}[label=\Alph*.]
\item Recomputing features differently online creates skew.
\item Treating object versioning alone as branch-level data semantics is insufficient.
\item Excess small files destroy scan performance despite correct table semantics.
\item Joining the latest feature value into old examples creates unrealistically strong training results.
\end{enumerate}
\noindent\textbf{Answer.} Treating object versioning alone as branch-level data semantics is insufficient.
\par\textit{lakeFS adds Git-like branches, commits, merges, and hooks over object-storage data lakes. Tradeoff: Atomic data-lake changes enable testing, with server and metadata operations to manage. Watch for: Treating object versioning alone as branch-level data semantics is insufficient.}
\paragraph{FC-1038 [medium]} Define lakeFS in interview-ready language.
\noindent\textbf{Answer.} lakeFS adds Git-like branches, commits, merges, and hooks over object-storage data lakes.
\par\textit{lakeFS adds Git-like branches, commits, merges, and hooks over object-storage data lakes.}
\paragraph{FC-1039 [hard]} State the main tradeoff and production pitfall for lakeFS.
\noindent\textbf{Answer.} Tradeoff: Atomic data-lake changes enable testing, with server and metadata operations to manage. Pitfall: Treating object versioning alone as branch-level data semantics is insufficient.
\par\textit{Tradeoff: Atomic data-lake changes enable testing, with server and metadata operations to manage. Pitfall: Treating object versioning alone as branch-level data semantics is insufficient.}
\paragraph{LA-1040 [hard]} Explain lakeFS from first principles, then show how you would evaluate and operate it in production.
\noindent\textbf{Answer.} Start with the mechanism: lakeFS adds Git-like branches, commits, merges, and hooks over object-storage data lakes. Make the tradeoff explicit: Atomic data-lake changes enable testing, with server and metadata operations to manage. Define workload-specific offline metrics and a latency/cost budget, test important slices, instrument the critical path, and create a rollback criterion. The production review must explicitly guard against this failure: Treating object versioning alone as branch-level data semantics is insufficient.
\par\textit{A strong answer connects mechanism, measurable tradeoffs, failure isolation, observability, and rollback rather than listing product features.}
\paragraph{MCQ-1041 [medium]} Which statement best describes Delta Lake and Iceberg?
\begin{enumerate}[label=\Alph*.]
\item Feature stores define reusable features and connect historical offline data with low-latency online serving.
\item Open table formats add schemas, snapshots, partition evolution, and transactional metadata over object storage.
\item lakeFS adds Git-like branches, commits, merges, and hooks over object-storage data lakes.
\item Tests cover schema, nulls, ranges, uniqueness, distribution, freshness, and cross-table invariants.
\end{enumerate}
\noindent\textbf{Answer.} Open table formats add schemas, snapshots, partition evolution, and transactional metadata over object storage.
\par\textit{Open table formats add schemas, snapshots, partition evolution, and transactional metadata over object storage. Tradeoff: Reliable analytics improve, while compaction and catalog operations become essential. Watch for: Excess small files destroy scan performance despite correct table semantics.}
\paragraph{MCQ-1042 [hard]} What is the central engineering tradeoff of Delta Lake and Iceberg?
\begin{enumerate}[label=\Alph*.]
\item Reliable analytics improve, while compaction and catalog operations become essential.
\item Early detection prevents downstream damage, while brittle thresholds create alert fatigue.
\item Atomic data-lake changes enable testing, with server and metadata operations to manage.
\item Training-serving consistency improves but adds infrastructure and ownership requirements.
\end{enumerate}
\noindent\textbf{Answer.} Reliable analytics improve, while compaction and catalog operations become essential.
\par\textit{Open table formats add schemas, snapshots, partition evolution, and transactional metadata over object storage. Tradeoff: Reliable analytics improve, while compaction and catalog operations become essential. Watch for: Excess small files destroy scan performance despite correct table semantics.}
\paragraph{MCQ-1043 [hard]} Which failure mode should an interviewer expect you to identify for Delta Lake and Iceberg?
\begin{enumerate}[label=\Alph*.]
\item Checking only aggregate distributions misses failures in important slices.
\item Excess small files destroy scan performance despite correct table semantics.
\item Treating object versioning alone as branch-level data semantics is insufficient.
\item Recomputing features differently online creates skew.
\end{enumerate}
\noindent\textbf{Answer.} Excess small files destroy scan performance despite correct table semantics.
\par\textit{Open table formats add schemas, snapshots, partition evolution, and transactional metadata over object storage. Tradeoff: Reliable analytics improve, while compaction and catalog operations become essential. Watch for: Excess small files destroy scan performance despite correct table semantics.}
\paragraph{MCQ-1044 [medium]} A design review is specifically concerned with this risk: Excess small files destroy scan performance despite correct table semantics. Which topic should be reviewed first?
\begin{enumerate}[label=\Alph*.]
\item Feature stores
\item Data quality testing
\item lakeFS
\item Delta Lake and Iceberg
\end{enumerate}
\noindent\textbf{Answer.} Delta Lake and Iceberg
\par\textit{Open table formats add schemas, snapshots, partition evolution, and transactional metadata over object storage. Tradeoff: Reliable analytics improve, while compaction and catalog operations become essential. Watch for: Excess small files destroy scan performance despite correct table semantics.}
\paragraph{MCQ-1045 [hard]} Which design note most accurately balances the benefits and costs of Delta Lake and Iceberg?
\begin{enumerate}[label=\Alph*.]
\item Atomic data-lake changes enable testing, with server and metadata operations to manage.
\item Training-serving consistency improves but adds infrastructure and ownership requirements.
\item Early detection prevents downstream damage, while brittle thresholds create alert fatigue.
\item Reliable analytics improve, while compaction and catalog operations become essential.
\end{enumerate}
\noindent\textbf{Answer.} Reliable analytics improve, while compaction and catalog operations become essential.
\par\textit{Open table formats add schemas, snapshots, partition evolution, and transactional metadata over object storage. Tradeoff: Reliable analytics improve, while compaction and catalog operations become essential. Watch for: Excess small files destroy scan performance despite correct table semantics.}
\paragraph{MCQ-1046 [medium]} Which explanation would be strongest in a system-design interview about Delta Lake and Iceberg?
\begin{enumerate}[label=\Alph*.]
\item lakeFS adds Git-like branches, commits, merges, and hooks over object-storage data lakes.
\item Tests cover schema, nulls, ranges, uniqueness, distribution, freshness, and cross-table invariants.
\item Open table formats add schemas, snapshots, partition evolution, and transactional metadata over object storage.
\item Feature stores define reusable features and connect historical offline data with low-latency online serving.
\end{enumerate}
\noindent\textbf{Answer.} Open table formats add schemas, snapshots, partition evolution, and transactional metadata over object storage.
\par\textit{Open table formats add schemas, snapshots, partition evolution, and transactional metadata over object storage. Tradeoff: Reliable analytics improve, while compaction and catalog operations become essential. Watch for: Excess small files destroy scan performance despite correct table semantics.}
\paragraph{MCQ-1047 [hard]} Which production warning is most directly associated with Delta Lake and Iceberg?
\begin{enumerate}[label=\Alph*.]
\item Treating object versioning alone as branch-level data semantics is insufficient.
\item Excess small files destroy scan performance despite correct table semantics.
\item Recomputing features differently online creates skew.
\item Checking only aggregate distributions misses failures in important slices.
\end{enumerate}
\noindent\textbf{Answer.} Excess small files destroy scan performance despite correct table semantics.
\par\textit{Open table formats add schemas, snapshots, partition evolution, and transactional metadata over object storage. Tradeoff: Reliable analytics improve, while compaction and catalog operations become essential. Watch for: Excess small files destroy scan performance despite correct table semantics.}
\paragraph{FC-1048 [medium]} Define Delta Lake and Iceberg in interview-ready language.
\noindent\textbf{Answer.} Open table formats add schemas, snapshots, partition evolution, and transactional metadata over object storage.
\par\textit{Open table formats add schemas, snapshots, partition evolution, and transactional metadata over object storage.}
\paragraph{FC-1049 [hard]} State the main tradeoff and production pitfall for Delta Lake and Iceberg.
\noindent\textbf{Answer.} Tradeoff: Reliable analytics improve, while compaction and catalog operations become essential. Pitfall: Excess small files destroy scan performance despite correct table semantics.
\par\textit{Tradeoff: Reliable analytics improve, while compaction and catalog operations become essential. Pitfall: Excess small files destroy scan performance despite correct table semantics.}
\paragraph{LA-1050 [hard]} Explain Delta Lake and Iceberg from first principles, then show how you would evaluate and operate it in production.
\noindent\textbf{Answer.} Start with the mechanism: Open table formats add schemas, snapshots, partition evolution, and transactional metadata over object storage. Make the tradeoff explicit: Reliable analytics improve, while compaction and catalog operations become essential. Define workload-specific offline metrics and a latency/cost budget, test important slices, instrument the critical path, and create a rollback criterion. The production review must explicitly guard against this failure: Excess small files destroy scan performance despite correct table semantics.
\par\textit{A strong answer connects mechanism, measurable tradeoffs, failure isolation, observability, and rollback rather than listing product features.}
\paragraph{MCQ-1051 [medium]} Which statement best describes Feature stores?
\begin{enumerate}[label=\Alph*.]
\item Feature stores define reusable features and connect historical offline data with low-latency online serving.
\item Lineage links sources, transformations, datasets, features, prompts, models, evaluations, and deployments.
\item Historical joins select only feature values known at each label timestamp.
\item Runs capture code, parameters, data, metrics, artifacts, environment, and parent relationships.
\end{enumerate}
\noindent\textbf{Answer.} Feature stores define reusable features and connect historical offline data with low-latency online serving.
\par\textit{Feature stores define reusable features and connect historical offline data with low-latency online serving. Tradeoff: Training-serving consistency improves but adds infrastructure and ownership requirements. Watch for: Recomputing features differently online creates skew.}
\paragraph{MCQ-1052 [hard]} What is the central engineering tradeoff of Feature stores?
\begin{enumerate}[label=\Alph*.]
\item Impact analysis and auditability improve, while automated lineage can miss dynamic behavior.
\item Comparison and reproducibility improve but inconsistent naming creates a junk drawer.
\item Training-serving consistency improves but adds infrastructure and ownership requirements.
\item Leakage is prevented at the cost of temporal keys and more expensive joins.
\end{enumerate}
\noindent\textbf{Answer.} Training-serving consistency improves but adds infrastructure and ownership requirements.
\par\textit{Feature stores define reusable features and connect historical offline data with low-latency online serving. Tradeoff: Training-serving consistency improves but adds infrastructure and ownership requirements. Watch for: Recomputing features differently online creates skew.}
\paragraph{MCQ-1053 [hard]} Which failure mode should an interviewer expect you to identify for Feature stores?
\begin{enumerate}[label=\Alph*.]
\item Joining the latest feature value into old examples creates unrealistically strong training results.
\item Logging a metric without dataset and code versions makes it uninterpretable.
\item Lineage without immutable identifiers gives a false sense of reproducibility.
\item Recomputing features differently online creates skew.
\end{enumerate}
\noindent\textbf{Answer.} Recomputing features differently online creates skew.
\par\textit{Feature stores define reusable features and connect historical offline data with low-latency online serving. Tradeoff: Training-serving consistency improves but adds infrastructure and ownership requirements. Watch for: Recomputing features differently online creates skew.}
\paragraph{MCQ-1054 [medium]} A design review is specifically concerned with this risk: Recomputing features differently online creates skew. Which topic should be reviewed first?
\begin{enumerate}[label=\Alph*.]
\item Point-in-time correctness
\item Feature stores
\item Data lineage
\item Experiment tracking
\end{enumerate}
\noindent\textbf{Answer.} Feature stores
\par\textit{Feature stores define reusable features and connect historical offline data with low-latency online serving. Tradeoff: Training-serving consistency improves but adds infrastructure and ownership requirements. Watch for: Recomputing features differently online creates skew.}
\paragraph{MCQ-1055 [hard]} Which design note most accurately balances the benefits and costs of Feature stores?
\begin{enumerate}[label=\Alph*.]
\item Impact analysis and auditability improve, while automated lineage can miss dynamic behavior.
\item Comparison and reproducibility improve but inconsistent naming creates a junk drawer.
\item Leakage is prevented at the cost of temporal keys and more expensive joins.
\item Training-serving consistency improves but adds infrastructure and ownership requirements.
\end{enumerate}
\noindent\textbf{Answer.} Training-serving consistency improves but adds infrastructure and ownership requirements.
\par\textit{Feature stores define reusable features and connect historical offline data with low-latency online serving. Tradeoff: Training-serving consistency improves but adds infrastructure and ownership requirements. Watch for: Recomputing features differently online creates skew.}
\paragraph{MCQ-1056 [medium]} Which explanation would be strongest in a system-design interview about Feature stores?
\begin{enumerate}[label=\Alph*.]
\item Historical joins select only feature values known at each label timestamp.
\item Feature stores define reusable features and connect historical offline data with low-latency online serving.
\item Lineage links sources, transformations, datasets, features, prompts, models, evaluations, and deployments.
\item Runs capture code, parameters, data, metrics, artifacts, environment, and parent relationships.
\end{enumerate}
\noindent\textbf{Answer.} Feature stores define reusable features and connect historical offline data with low-latency online serving.
\par\textit{Feature stores define reusable features and connect historical offline data with low-latency online serving. Tradeoff: Training-serving consistency improves but adds infrastructure and ownership requirements. Watch for: Recomputing features differently online creates skew.}
\paragraph{MCQ-1057 [hard]} Which production warning is most directly associated with Feature stores?
\begin{enumerate}[label=\Alph*.]
\item Recomputing features differently online creates skew.
\item Lineage without immutable identifiers gives a false sense of reproducibility.
\item Joining the latest feature value into old examples creates unrealistically strong training results.
\item Logging a metric without dataset and code versions makes it uninterpretable.
\end{enumerate}
\noindent\textbf{Answer.} Recomputing features differently online creates skew.
\par\textit{Feature stores define reusable features and connect historical offline data with low-latency online serving. Tradeoff: Training-serving consistency improves but adds infrastructure and ownership requirements. Watch for: Recomputing features differently online creates skew.}
\paragraph{FC-1058 [medium]} Define Feature stores in interview-ready language.
\noindent\textbf{Answer.} Feature stores define reusable features and connect historical offline data with low-latency online serving.
\par\textit{Feature stores define reusable features and connect historical offline data with low-latency online serving.}
\paragraph{FC-1059 [hard]} State the main tradeoff and production pitfall for Feature stores.
\noindent\textbf{Answer.} Tradeoff: Training-serving consistency improves but adds infrastructure and ownership requirements. Pitfall: Recomputing features differently online creates skew.
\par\textit{Tradeoff: Training-serving consistency improves but adds infrastructure and ownership requirements. Pitfall: Recomputing features differently online creates skew.}
\paragraph{LA-1060 [hard]} Explain Feature stores from first principles, then show how you would evaluate and operate it in production.
\noindent\textbf{Answer.} Start with the mechanism: Feature stores define reusable features and connect historical offline data with low-latency online serving. Make the tradeoff explicit: Training-serving consistency improves but adds infrastructure and ownership requirements. Define workload-specific offline metrics and a latency/cost budget, test important slices, instrument the critical path, and create a rollback criterion. The production review must explicitly guard against this failure: Recomputing features differently online creates skew.
\par\textit{A strong answer connects mechanism, measurable tradeoffs, failure isolation, observability, and rollback rather than listing product features.}
\paragraph{MCQ-1061 [medium]} Which statement best describes Point-in-time correctness?
\begin{enumerate}[label=\Alph*.]
\item Historical joins select only feature values known at each label timestamp.
\item DVC versions large data and model pointers with Git and defines reproducible pipeline stages and experiments.
\item lakeFS adds Git-like branches, commits, merges, and hooks over object-storage data lakes.
\item Lockfiles, containers, hardware metadata, seeds, and deterministic settings constrain run variability.
\end{enumerate}
\noindent\textbf{Answer.} Historical joins select only feature values known at each label timestamp.
\par\textit{Historical joins select only feature values known at each label timestamp. Tradeoff: Leakage is prevented at the cost of temporal keys and more expensive joins. Watch for: Joining the latest feature value into old examples creates unrealistically strong training results.}
\paragraph{MCQ-1062 [hard]} What is the central engineering tradeoff of Point-in-time correctness?
\begin{enumerate}[label=\Alph*.]
\item Atomic data-lake changes enable testing, with server and metadata operations to manage.
\item Familiar Git workflows improve reproducibility, while remote storage and cache discipline are required.
\item Leakage is prevented at the cost of temporal keys and more expensive joins.
\item Reproducibility improves at the expense of build maintenance and sometimes performance.
\end{enumerate}
\noindent\textbf{Answer.} Leakage is prevented at the cost of temporal keys and more expensive joins.
\par\textit{Historical joins select only feature values known at each label timestamp. Tradeoff: Leakage is prevented at the cost of temporal keys and more expensive joins. Watch for: Joining the latest feature value into old examples creates unrealistically strong training results.}
\paragraph{MCQ-1063 [hard]} Which failure mode should an interviewer expect you to identify for Point-in-time correctness?
\begin{enumerate}[label=\Alph*.]
\item A random seed alone cannot guarantee deterministic GPU execution.
\item Treating object versioning alone as branch-level data semantics is insufficient.
\item Committing only code while silently replacing remote data breaks lineage.
\item Joining the latest feature value into old examples creates unrealistically strong training results.
\end{enumerate}
\noindent\textbf{Answer.} Joining the latest feature value into old examples creates unrealistically strong training results.
\par\textit{Historical joins select only feature values known at each label timestamp. Tradeoff: Leakage is prevented at the cost of temporal keys and more expensive joins. Watch for: Joining the latest feature value into old examples creates unrealistically strong training results.}
\paragraph{MCQ-1064 [medium]} A design review is specifically concerned with this risk: Joining the latest feature value into old examples creates unrealistically strong training results. Which topic should be reviewed first?
\begin{enumerate}[label=\Alph*.]
\item Point-in-time correctness
\item lakeFS
\item Reproducible environments
\item DVC
\end{enumerate}
\noindent\textbf{Answer.} Point-in-time correctness
\par\textit{Historical joins select only feature values known at each label timestamp. Tradeoff: Leakage is prevented at the cost of temporal keys and more expensive joins. Watch for: Joining the latest feature value into old examples creates unrealistically strong training results.}
\paragraph{MCQ-1065 [hard]} Which design note most accurately balances the benefits and costs of Point-in-time correctness?
\begin{enumerate}[label=\Alph*.]
\item Familiar Git workflows improve reproducibility, while remote storage and cache discipline are required.
\item Reproducibility improves at the expense of build maintenance and sometimes performance.
\item Atomic data-lake changes enable testing, with server and metadata operations to manage.
\item Leakage is prevented at the cost of temporal keys and more expensive joins.
\end{enumerate}
\noindent\textbf{Answer.} Leakage is prevented at the cost of temporal keys and more expensive joins.
\par\textit{Historical joins select only feature values known at each label timestamp. Tradeoff: Leakage is prevented at the cost of temporal keys and more expensive joins. Watch for: Joining the latest feature value into old examples creates unrealistically strong training results.}
\paragraph{MCQ-1066 [medium]} Which explanation would be strongest in a system-design interview about Point-in-time correctness?
\begin{enumerate}[label=\Alph*.]
\item DVC versions large data and model pointers with Git and defines reproducible pipeline stages and experiments.
\item Historical joins select only feature values known at each label timestamp.
\item lakeFS adds Git-like branches, commits, merges, and hooks over object-storage data lakes.
\item Lockfiles, containers, hardware metadata, seeds, and deterministic settings constrain run variability.
\end{enumerate}
\noindent\textbf{Answer.} Historical joins select only feature values known at each label timestamp.
\par\textit{Historical joins select only feature values known at each label timestamp. Tradeoff: Leakage is prevented at the cost of temporal keys and more expensive joins. Watch for: Joining the latest feature value into old examples creates unrealistically strong training results.}
\paragraph{MCQ-1067 [hard]} Which production warning is most directly associated with Point-in-time correctness?
\begin{enumerate}[label=\Alph*.]
\item Treating object versioning alone as branch-level data semantics is insufficient.
\item Joining the latest feature value into old examples creates unrealistically strong training results.
\item Committing only code while silently replacing remote data breaks lineage.
\item A random seed alone cannot guarantee deterministic GPU execution.
\end{enumerate}
\noindent\textbf{Answer.} Joining the latest feature value into old examples creates unrealistically strong training results.
\par\textit{Historical joins select only feature values known at each label timestamp. Tradeoff: Leakage is prevented at the cost of temporal keys and more expensive joins. Watch for: Joining the latest feature value into old examples creates unrealistically strong training results.}
\paragraph{FC-1068 [medium]} Define Point-in-time correctness in interview-ready language.
\noindent\textbf{Answer.} Historical joins select only feature values known at each label timestamp.
\par\textit{Historical joins select only feature values known at each label timestamp.}
\paragraph{FC-1069 [hard]} State the main tradeoff and production pitfall for Point-in-time correctness.
\noindent\textbf{Answer.} Tradeoff: Leakage is prevented at the cost of temporal keys and more expensive joins. Pitfall: Joining the latest feature value into old examples creates unrealistically strong training results.
\par\textit{Tradeoff: Leakage is prevented at the cost of temporal keys and more expensive joins. Pitfall: Joining the latest feature value into old examples creates unrealistically strong training results.}
\paragraph{LA-1070 [hard]} Explain Point-in-time correctness from first principles, then show how you would evaluate and operate it in production.
\noindent\textbf{Answer.} Start with the mechanism: Historical joins select only feature values known at each label timestamp. Make the tradeoff explicit: Leakage is prevented at the cost of temporal keys and more expensive joins. Define workload-specific offline metrics and a latency/cost budget, test important slices, instrument the critical path, and create a rollback criterion. The production review must explicitly guard against this failure: Joining the latest feature value into old examples creates unrealistically strong training results.
\par\textit{A strong answer connects mechanism, measurable tradeoffs, failure isolation, observability, and rollback rather than listing product features.}
\paragraph{MCQ-1071 [medium]} Which statement best describes Experiment tracking?
\begin{enumerate}[label=\Alph*.]
\item Historical joins select only feature values known at each label timestamp.
\item Feature stores define reusable features and connect historical offline data with low-latency online serving.
\item Tests cover schema, nulls, ranges, uniqueness, distribution, freshness, and cross-table invariants.
\item Runs capture code, parameters, data, metrics, artifacts, environment, and parent relationships.
\end{enumerate}
\noindent\textbf{Answer.} Runs capture code, parameters, data, metrics, artifacts, environment, and parent relationships.
\par\textit{Runs capture code, parameters, data, metrics, artifacts, environment, and parent relationships. Tradeoff: Comparison and reproducibility improve but inconsistent naming creates a junk drawer. Watch for: Logging a metric without dataset and code versions makes it uninterpretable.}
\paragraph{MCQ-1072 [hard]} What is the central engineering tradeoff of Experiment tracking?
\begin{enumerate}[label=\Alph*.]
\item Comparison and reproducibility improve but inconsistent naming creates a junk drawer.
\item Training-serving consistency improves but adds infrastructure and ownership requirements.
\item Early detection prevents downstream damage, while brittle thresholds create alert fatigue.
\item Leakage is prevented at the cost of temporal keys and more expensive joins.
\end{enumerate}
\noindent\textbf{Answer.} Comparison and reproducibility improve but inconsistent naming creates a junk drawer.
\par\textit{Runs capture code, parameters, data, metrics, artifacts, environment, and parent relationships. Tradeoff: Comparison and reproducibility improve but inconsistent naming creates a junk drawer. Watch for: Logging a metric without dataset and code versions makes it uninterpretable.}
\paragraph{MCQ-1073 [hard]} Which failure mode should an interviewer expect you to identify for Experiment tracking?
\begin{enumerate}[label=\Alph*.]
\item Checking only aggregate distributions misses failures in important slices.
\item Recomputing features differently online creates skew.
\item Logging a metric without dataset and code versions makes it uninterpretable.
\item Joining the latest feature value into old examples creates unrealistically strong training results.
\end{enumerate}
\noindent\textbf{Answer.} Logging a metric without dataset and code versions makes it uninterpretable.
\par\textit{Runs capture code, parameters, data, metrics, artifacts, environment, and parent relationships. Tradeoff: Comparison and reproducibility improve but inconsistent naming creates a junk drawer. Watch for: Logging a metric without dataset and code versions makes it uninterpretable.}
\paragraph{MCQ-1074 [medium]} A design review is specifically concerned with this risk: Logging a metric without dataset and code versions makes it uninterpretable. Which topic should be reviewed first?
\begin{enumerate}[label=\Alph*.]
\item Experiment tracking
\item Feature stores
\item Point-in-time correctness
\item Data quality testing
\end{enumerate}
\noindent\textbf{Answer.} Experiment tracking
\par\textit{Runs capture code, parameters, data, metrics, artifacts, environment, and parent relationships. Tradeoff: Comparison and reproducibility improve but inconsistent naming creates a junk drawer. Watch for: Logging a metric without dataset and code versions makes it uninterpretable.}
\paragraph{MCQ-1075 [hard]} Which design note most accurately balances the benefits and costs of Experiment tracking?
\begin{enumerate}[label=\Alph*.]
\item Leakage is prevented at the cost of temporal keys and more expensive joins.
\item Early detection prevents downstream damage, while brittle thresholds create alert fatigue.
\item Comparison and reproducibility improve but inconsistent naming creates a junk drawer.
\item Training-serving consistency improves but adds infrastructure and ownership requirements.
\end{enumerate}
\noindent\textbf{Answer.} Comparison and reproducibility improve but inconsistent naming creates a junk drawer.
\par\textit{Runs capture code, parameters, data, metrics, artifacts, environment, and parent relationships. Tradeoff: Comparison and reproducibility improve but inconsistent naming creates a junk drawer. Watch for: Logging a metric without dataset and code versions makes it uninterpretable.}
\paragraph{MCQ-1076 [medium]} Which explanation would be strongest in a system-design interview about Experiment tracking?
\begin{enumerate}[label=\Alph*.]
\item Runs capture code, parameters, data, metrics, artifacts, environment, and parent relationships.
\item Tests cover schema, nulls, ranges, uniqueness, distribution, freshness, and cross-table invariants.
\item Feature stores define reusable features and connect historical offline data with low-latency online serving.
\item Historical joins select only feature values known at each label timestamp.
\end{enumerate}
\noindent\textbf{Answer.} Runs capture code, parameters, data, metrics, artifacts, environment, and parent relationships.
\par\textit{Runs capture code, parameters, data, metrics, artifacts, environment, and parent relationships. Tradeoff: Comparison and reproducibility improve but inconsistent naming creates a junk drawer. Watch for: Logging a metric without dataset and code versions makes it uninterpretable.}
\paragraph{MCQ-1077 [hard]} Which production warning is most directly associated with Experiment tracking?
\begin{enumerate}[label=\Alph*.]
\item Joining the latest feature value into old examples creates unrealistically strong training results.
\item Logging a metric without dataset and code versions makes it uninterpretable.
\item Recomputing features differently online creates skew.
\item Checking only aggregate distributions misses failures in important slices.
\end{enumerate}
\noindent\textbf{Answer.} Logging a metric without dataset and code versions makes it uninterpretable.
\par\textit{Runs capture code, parameters, data, metrics, artifacts, environment, and parent relationships. Tradeoff: Comparison and reproducibility improve but inconsistent naming creates a junk drawer. Watch for: Logging a metric without dataset and code versions makes it uninterpretable.}
\paragraph{FC-1078 [medium]} Define Experiment tracking in interview-ready language.
\noindent\textbf{Answer.} Runs capture code, parameters, data, metrics, artifacts, environment, and parent relationships.
\par\textit{Runs capture code, parameters, data, metrics, artifacts, environment, and parent relationships.}
\paragraph{FC-1079 [hard]} State the main tradeoff and production pitfall for Experiment tracking.
\noindent\textbf{Answer.} Tradeoff: Comparison and reproducibility improve but inconsistent naming creates a junk drawer. Pitfall: Logging a metric without dataset and code versions makes it uninterpretable.
\par\textit{Tradeoff: Comparison and reproducibility improve but inconsistent naming creates a junk drawer. Pitfall: Logging a metric without dataset and code versions makes it uninterpretable.}
\paragraph{LA-1080 [hard]} Explain Experiment tracking from first principles, then show how you would evaluate and operate it in production.
\noindent\textbf{Answer.} Start with the mechanism: Runs capture code, parameters, data, metrics, artifacts, environment, and parent relationships. Make the tradeoff explicit: Comparison and reproducibility improve but inconsistent naming creates a junk drawer. Define workload-specific offline metrics and a latency/cost budget, test important slices, instrument the critical path, and create a rollback criterion. The production review must explicitly guard against this failure: Logging a metric without dataset and code versions makes it uninterpretable.
\par\textit{A strong answer connects mechanism, measurable tradeoffs, failure isolation, observability, and rollback rather than listing product features.}
\paragraph{MCQ-1081 [medium]} Which statement best describes Data contracts?
\begin{enumerate}[label=\Alph*.]
\item Runs capture code, parameters, data, metrics, artifacts, environment, and parent relationships.
\item Contracts specify schema, semantics, freshness, ownership, quality, and compatibility expectations between producers and consumers.
\item Historical joins select only feature values known at each label timestamp.
\item lakeFS adds Git-like branches, commits, merges, and hooks over object-storage data lakes.
\end{enumerate}
\noindent\textbf{Answer.} Contracts specify schema, semantics, freshness, ownership, quality, and compatibility expectations between producers and consumers.
\par\textit{Contracts specify schema, semantics, freshness, ownership, quality, and compatibility expectations between producers and consumers. Tradeoff: Clear boundaries reduce silent breakage but require governance and enforcement. Watch for: Validating shape without semantic definitions misses unit and meaning changes.}
\paragraph{MCQ-1082 [hard]} What is the central engineering tradeoff of Data contracts?
\begin{enumerate}[label=\Alph*.]
\item Comparison and reproducibility improve but inconsistent naming creates a junk drawer.
\item Leakage is prevented at the cost of temporal keys and more expensive joins.
\item Clear boundaries reduce silent breakage but require governance and enforcement.
\item Atomic data-lake changes enable testing, with server and metadata operations to manage.
\end{enumerate}
\noindent\textbf{Answer.} Clear boundaries reduce silent breakage but require governance and enforcement.
\par\textit{Contracts specify schema, semantics, freshness, ownership, quality, and compatibility expectations between producers and consumers. Tradeoff: Clear boundaries reduce silent breakage but require governance and enforcement. Watch for: Validating shape without semantic definitions misses unit and meaning changes.}
\paragraph{MCQ-1083 [hard]} Which failure mode should an interviewer expect you to identify for Data contracts?
\begin{enumerate}[label=\Alph*.]
\item Treating object versioning alone as branch-level data semantics is insufficient.
\item Logging a metric without dataset and code versions makes it uninterpretable.
\item Validating shape without semantic definitions misses unit and meaning changes.
\item Joining the latest feature value into old examples creates unrealistically strong training results.
\end{enumerate}
\noindent\textbf{Answer.} Validating shape without semantic definitions misses unit and meaning changes.
\par\textit{Contracts specify schema, semantics, freshness, ownership, quality, and compatibility expectations between producers and consumers. Tradeoff: Clear boundaries reduce silent breakage but require governance and enforcement. Watch for: Validating shape without semantic definitions misses unit and meaning changes.}
\paragraph{MCQ-1084 [medium]} A design review is specifically concerned with this risk: Validating shape without semantic definitions misses unit and meaning changes. Which topic should be reviewed first?
\begin{enumerate}[label=\Alph*.]
\item Data contracts
\item Experiment tracking
\item lakeFS
\item Point-in-time correctness
\end{enumerate}
\noindent\textbf{Answer.} Data contracts
\par\textit{Contracts specify schema, semantics, freshness, ownership, quality, and compatibility expectations between producers and consumers. Tradeoff: Clear boundaries reduce silent breakage but require governance and enforcement. Watch for: Validating shape without semantic definitions misses unit and meaning changes.}
\paragraph{MCQ-1085 [hard]} Which design note most accurately balances the benefits and costs of Data contracts?
\begin{enumerate}[label=\Alph*.]
\item Clear boundaries reduce silent breakage but require governance and enforcement.
\item Comparison and reproducibility improve but inconsistent naming creates a junk drawer.
\item Atomic data-lake changes enable testing, with server and metadata operations to manage.
\item Leakage is prevented at the cost of temporal keys and more expensive joins.
\end{enumerate}
\noindent\textbf{Answer.} Clear boundaries reduce silent breakage but require governance and enforcement.
\par\textit{Contracts specify schema, semantics, freshness, ownership, quality, and compatibility expectations between producers and consumers. Tradeoff: Clear boundaries reduce silent breakage but require governance and enforcement. Watch for: Validating shape without semantic definitions misses unit and meaning changes.}
\paragraph{MCQ-1086 [medium]} Which explanation would be strongest in a system-design interview about Data contracts?
\begin{enumerate}[label=\Alph*.]
\item lakeFS adds Git-like branches, commits, merges, and hooks over object-storage data lakes.
\item Contracts specify schema, semantics, freshness, ownership, quality, and compatibility expectations between producers and consumers.
\item Historical joins select only feature values known at each label timestamp.
\item Runs capture code, parameters, data, metrics, artifacts, environment, and parent relationships.
\end{enumerate}
\noindent\textbf{Answer.} Contracts specify schema, semantics, freshness, ownership, quality, and compatibility expectations between producers and consumers.
\par\textit{Contracts specify schema, semantics, freshness, ownership, quality, and compatibility expectations between producers and consumers. Tradeoff: Clear boundaries reduce silent breakage but require governance and enforcement. Watch for: Validating shape without semantic definitions misses unit and meaning changes.}
\paragraph{MCQ-1087 [hard]} Which production warning is most directly associated with Data contracts?
\begin{enumerate}[label=\Alph*.]
\item Logging a metric without dataset and code versions makes it uninterpretable.
\item Validating shape without semantic definitions misses unit and meaning changes.
\item Treating object versioning alone as branch-level data semantics is insufficient.
\item Joining the latest feature value into old examples creates unrealistically strong training results.
\end{enumerate}
\noindent\textbf{Answer.} Validating shape without semantic definitions misses unit and meaning changes.
\par\textit{Contracts specify schema, semantics, freshness, ownership, quality, and compatibility expectations between producers and consumers. Tradeoff: Clear boundaries reduce silent breakage but require governance and enforcement. Watch for: Validating shape without semantic definitions misses unit and meaning changes.}
\paragraph{FC-1088 [medium]} Define Data contracts in interview-ready language.
\noindent\textbf{Answer.} Contracts specify schema, semantics, freshness, ownership, quality, and compatibility expectations between producers and consumers.
\par\textit{Contracts specify schema, semantics, freshness, ownership, quality, and compatibility expectations between producers and consumers.}
\paragraph{FC-1089 [hard]} State the main tradeoff and production pitfall for Data contracts.
\noindent\textbf{Answer.} Tradeoff: Clear boundaries reduce silent breakage but require governance and enforcement. Pitfall: Validating shape without semantic definitions misses unit and meaning changes.
\par\textit{Tradeoff: Clear boundaries reduce silent breakage but require governance and enforcement. Pitfall: Validating shape without semantic definitions misses unit and meaning changes.}
\paragraph{LA-1090 [hard]} Explain Data contracts from first principles, then show how you would evaluate and operate it in production.
\noindent\textbf{Answer.} Start with the mechanism: Contracts specify schema, semantics, freshness, ownership, quality, and compatibility expectations between producers and consumers. Make the tradeoff explicit: Clear boundaries reduce silent breakage but require governance and enforcement. Define workload-specific offline metrics and a latency/cost budget, test important slices, instrument the critical path, and create a rollback criterion. The production review must explicitly guard against this failure: Validating shape without semantic definitions misses unit and meaning changes.
\par\textit{A strong answer connects mechanism, measurable tradeoffs, failure isolation, observability, and rollback rather than listing product features.}
\paragraph{MCQ-1091 [medium]} Which statement best describes Data quality testing?
\begin{enumerate}[label=\Alph*.]
\item Tests cover schema, nulls, ranges, uniqueness, distribution, freshness, and cross-table invariants.
\item Runs capture code, parameters, data, metrics, artifacts, environment, and parent relationships.
\item Lineage links sources, transformations, datasets, features, prompts, models, evaluations, and deployments.
\item lakeFS adds Git-like branches, commits, merges, and hooks over object-storage data lakes.
\end{enumerate}
\noindent\textbf{Answer.} Tests cover schema, nulls, ranges, uniqueness, distribution, freshness, and cross-table invariants.
\par\textit{Tests cover schema, nulls, ranges, uniqueness, distribution, freshness, and cross-table invariants. Tradeoff: Early detection prevents downstream damage, while brittle thresholds create alert fatigue. Watch for: Checking only aggregate distributions misses failures in important slices.}
\paragraph{MCQ-1092 [hard]} What is the central engineering tradeoff of Data quality testing?
\begin{enumerate}[label=\Alph*.]
\item Impact analysis and auditability improve, while automated lineage can miss dynamic behavior.
\item Early detection prevents downstream damage, while brittle thresholds create alert fatigue.
\item Comparison and reproducibility improve but inconsistent naming creates a junk drawer.
\item Atomic data-lake changes enable testing, with server and metadata operations to manage.
\end{enumerate}
\noindent\textbf{Answer.} Early detection prevents downstream damage, while brittle thresholds create alert fatigue.
\par\textit{Tests cover schema, nulls, ranges, uniqueness, distribution, freshness, and cross-table invariants. Tradeoff: Early detection prevents downstream damage, while brittle thresholds create alert fatigue. Watch for: Checking only aggregate distributions misses failures in important slices.}
\paragraph{MCQ-1093 [hard]} Which failure mode should an interviewer expect you to identify for Data quality testing?
\begin{enumerate}[label=\Alph*.]
\item Checking only aggregate distributions misses failures in important slices.
\item Treating object versioning alone as branch-level data semantics is insufficient.
\item Lineage without immutable identifiers gives a false sense of reproducibility.
\item Logging a metric without dataset and code versions makes it uninterpretable.
\end{enumerate}
\noindent\textbf{Answer.} Checking only aggregate distributions misses failures in important slices.
\par\textit{Tests cover schema, nulls, ranges, uniqueness, distribution, freshness, and cross-table invariants. Tradeoff: Early detection prevents downstream damage, while brittle thresholds create alert fatigue. Watch for: Checking only aggregate distributions misses failures in important slices.}
\paragraph{MCQ-1094 [medium]} A design review is specifically concerned with this risk: Checking only aggregate distributions misses failures in important slices. Which topic should be reviewed first?
\begin{enumerate}[label=\Alph*.]
\item Experiment tracking
\item Data quality testing
\item Data lineage
\item lakeFS
\end{enumerate}
\noindent\textbf{Answer.} Data quality testing
\par\textit{Tests cover schema, nulls, ranges, uniqueness, distribution, freshness, and cross-table invariants. Tradeoff: Early detection prevents downstream damage, while brittle thresholds create alert fatigue. Watch for: Checking only aggregate distributions misses failures in important slices.}
\paragraph{MCQ-1095 [hard]} Which design note most accurately balances the benefits and costs of Data quality testing?
\begin{enumerate}[label=\Alph*.]
\item Impact analysis and auditability improve, while automated lineage can miss dynamic behavior.
\item Early detection prevents downstream damage, while brittle thresholds create alert fatigue.
\item Comparison and reproducibility improve but inconsistent naming creates a junk drawer.
\item Atomic data-lake changes enable testing, with server and metadata operations to manage.
\end{enumerate}
\noindent\textbf{Answer.} Early detection prevents downstream damage, while brittle thresholds create alert fatigue.
\par\textit{Tests cover schema, nulls, ranges, uniqueness, distribution, freshness, and cross-table invariants. Tradeoff: Early detection prevents downstream damage, while brittle thresholds create alert fatigue. Watch for: Checking only aggregate distributions misses failures in important slices.}
\paragraph{MCQ-1096 [medium]} Which explanation would be strongest in a system-design interview about Data quality testing?
\begin{enumerate}[label=\Alph*.]
\item Lineage links sources, transformations, datasets, features, prompts, models, evaluations, and deployments.
\item lakeFS adds Git-like branches, commits, merges, and hooks over object-storage data lakes.
\item Runs capture code, parameters, data, metrics, artifacts, environment, and parent relationships.
\item Tests cover schema, nulls, ranges, uniqueness, distribution, freshness, and cross-table invariants.
\end{enumerate}
\noindent\textbf{Answer.} Tests cover schema, nulls, ranges, uniqueness, distribution, freshness, and cross-table invariants.
\par\textit{Tests cover schema, nulls, ranges, uniqueness, distribution, freshness, and cross-table invariants. Tradeoff: Early detection prevents downstream damage, while brittle thresholds create alert fatigue. Watch for: Checking only aggregate distributions misses failures in important slices.}
\paragraph{MCQ-1097 [hard]} Which production warning is most directly associated with Data quality testing?
\begin{enumerate}[label=\Alph*.]
\item Checking only aggregate distributions misses failures in important slices.
\item Logging a metric without dataset and code versions makes it uninterpretable.
\item Lineage without immutable identifiers gives a false sense of reproducibility.
\item Treating object versioning alone as branch-level data semantics is insufficient.
\end{enumerate}
\noindent\textbf{Answer.} Checking only aggregate distributions misses failures in important slices.
\par\textit{Tests cover schema, nulls, ranges, uniqueness, distribution, freshness, and cross-table invariants. Tradeoff: Early detection prevents downstream damage, while brittle thresholds create alert fatigue. Watch for: Checking only aggregate distributions misses failures in important slices.}
\paragraph{FC-1098 [medium]} Define Data quality testing in interview-ready language.
\noindent\textbf{Answer.} Tests cover schema, nulls, ranges, uniqueness, distribution, freshness, and cross-table invariants.
\par\textit{Tests cover schema, nulls, ranges, uniqueness, distribution, freshness, and cross-table invariants.}
\paragraph{FC-1099 [hard]} State the main tradeoff and production pitfall for Data quality testing.
\noindent\textbf{Answer.} Tradeoff: Early detection prevents downstream damage, while brittle thresholds create alert fatigue. Pitfall: Checking only aggregate distributions misses failures in important slices.
\par\textit{Tradeoff: Early detection prevents downstream damage, while brittle thresholds create alert fatigue. Pitfall: Checking only aggregate distributions misses failures in important slices.}
\paragraph{LA-1100 [hard]} Explain Data quality testing from first principles, then show how you would evaluate and operate it in production.
\noindent\textbf{Answer.} Start with the mechanism: Tests cover schema, nulls, ranges, uniqueness, distribution, freshness, and cross-table invariants. Make the tradeoff explicit: Early detection prevents downstream damage, while brittle thresholds create alert fatigue. Define workload-specific offline metrics and a latency/cost budget, test important slices, instrument the critical path, and create a rollback criterion. The production review must explicitly guard against this failure: Checking only aggregate distributions misses failures in important slices.
\par\textit{A strong answer connects mechanism, measurable tradeoffs, failure isolation, observability, and rollback rather than listing product features.}
\paragraph{MCQ-1101 [medium]} Which statement best describes Data lineage?
\begin{enumerate}[label=\Alph*.]
\item Contracts specify schema, semantics, freshness, ownership, quality, and compatibility expectations between producers and consumers.
\item lakeFS adds Git-like branches, commits, merges, and hooks over object-storage data lakes.
\item Lineage links sources, transformations, datasets, features, prompts, models, evaluations, and deployments.
\item Runs capture code, parameters, data, metrics, artifacts, environment, and parent relationships.
\end{enumerate}
\noindent\textbf{Answer.} Lineage links sources, transformations, datasets, features, prompts, models, evaluations, and deployments.
\par\textit{Lineage links sources, transformations, datasets, features, prompts, models, evaluations, and deployments. Tradeoff: Impact analysis and auditability improve, while automated lineage can miss dynamic behavior. Watch for: Lineage without immutable identifiers gives a false sense of reproducibility.}
\paragraph{MCQ-1102 [hard]} What is the central engineering tradeoff of Data lineage?
\begin{enumerate}[label=\Alph*.]
\item Clear boundaries reduce silent breakage but require governance and enforcement.
\item Impact analysis and auditability improve, while automated lineage can miss dynamic behavior.
\item Comparison and reproducibility improve but inconsistent naming creates a junk drawer.
\item Atomic data-lake changes enable testing, with server and metadata operations to manage.
\end{enumerate}
\noindent\textbf{Answer.} Impact analysis and auditability improve, while automated lineage can miss dynamic behavior.
\par\textit{Lineage links sources, transformations, datasets, features, prompts, models, evaluations, and deployments. Tradeoff: Impact analysis and auditability improve, while automated lineage can miss dynamic behavior. Watch for: Lineage without immutable identifiers gives a false sense of reproducibility.}
\paragraph{MCQ-1103 [hard]} Which failure mode should an interviewer expect you to identify for Data lineage?
\begin{enumerate}[label=\Alph*.]
\item Validating shape without semantic definitions misses unit and meaning changes.
\item Lineage without immutable identifiers gives a false sense of reproducibility.
\item Treating object versioning alone as branch-level data semantics is insufficient.
\item Logging a metric without dataset and code versions makes it uninterpretable.
\end{enumerate}
\noindent\textbf{Answer.} Lineage without immutable identifiers gives a false sense of reproducibility.
\par\textit{Lineage links sources, transformations, datasets, features, prompts, models, evaluations, and deployments. Tradeoff: Impact analysis and auditability improve, while automated lineage can miss dynamic behavior. Watch for: Lineage without immutable identifiers gives a false sense of reproducibility.}
\paragraph{MCQ-1104 [medium]} A design review is specifically concerned with this risk: Lineage without immutable identifiers gives a false sense of reproducibility. Which topic should be reviewed first?
\begin{enumerate}[label=\Alph*.]
\item Data lineage
\item lakeFS
\item Data contracts
\item Experiment tracking
\end{enumerate}
\noindent\textbf{Answer.} Data lineage
\par\textit{Lineage links sources, transformations, datasets, features, prompts, models, evaluations, and deployments. Tradeoff: Impact analysis and auditability improve, while automated lineage can miss dynamic behavior. Watch for: Lineage without immutable identifiers gives a false sense of reproducibility.}
\paragraph{MCQ-1105 [hard]} Which design note most accurately balances the benefits and costs of Data lineage?
\begin{enumerate}[label=\Alph*.]
\item Impact analysis and auditability improve, while automated lineage can miss dynamic behavior.
\item Comparison and reproducibility improve but inconsistent naming creates a junk drawer.
\item Clear boundaries reduce silent breakage but require governance and enforcement.
\item Atomic data-lake changes enable testing, with server and metadata operations to manage.
\end{enumerate}
\noindent\textbf{Answer.} Impact analysis and auditability improve, while automated lineage can miss dynamic behavior.
\par\textit{Lineage links sources, transformations, datasets, features, prompts, models, evaluations, and deployments. Tradeoff: Impact analysis and auditability improve, while automated lineage can miss dynamic behavior. Watch for: Lineage without immutable identifiers gives a false sense of reproducibility.}
\paragraph{MCQ-1106 [medium]} Which explanation would be strongest in a system-design interview about Data lineage?
\begin{enumerate}[label=\Alph*.]
\item Contracts specify schema, semantics, freshness, ownership, quality, and compatibility expectations between producers and consumers.
\item lakeFS adds Git-like branches, commits, merges, and hooks over object-storage data lakes.
\item Runs capture code, parameters, data, metrics, artifacts, environment, and parent relationships.
\item Lineage links sources, transformations, datasets, features, prompts, models, evaluations, and deployments.
\end{enumerate}
\noindent\textbf{Answer.} Lineage links sources, transformations, datasets, features, prompts, models, evaluations, and deployments.
\par\textit{Lineage links sources, transformations, datasets, features, prompts, models, evaluations, and deployments. Tradeoff: Impact analysis and auditability improve, while automated lineage can miss dynamic behavior. Watch for: Lineage without immutable identifiers gives a false sense of reproducibility.}
\paragraph{MCQ-1107 [hard]} Which production warning is most directly associated with Data lineage?
\begin{enumerate}[label=\Alph*.]
\item Validating shape without semantic definitions misses unit and meaning changes.
\item Logging a metric without dataset and code versions makes it uninterpretable.
\item Treating object versioning alone as branch-level data semantics is insufficient.
\item Lineage without immutable identifiers gives a false sense of reproducibility.
\end{enumerate}
\noindent\textbf{Answer.} Lineage without immutable identifiers gives a false sense of reproducibility.
\par\textit{Lineage links sources, transformations, datasets, features, prompts, models, evaluations, and deployments. Tradeoff: Impact analysis and auditability improve, while automated lineage can miss dynamic behavior. Watch for: Lineage without immutable identifiers gives a false sense of reproducibility.}
\paragraph{FC-1108 [medium]} Define Data lineage in interview-ready language.
\noindent\textbf{Answer.} Lineage links sources, transformations, datasets, features, prompts, models, evaluations, and deployments.
\par\textit{Lineage links sources, transformations, datasets, features, prompts, models, evaluations, and deployments.}
\paragraph{FC-1109 [hard]} State the main tradeoff and production pitfall for Data lineage.
\noindent\textbf{Answer.} Tradeoff: Impact analysis and auditability improve, while automated lineage can miss dynamic behavior. Pitfall: Lineage without immutable identifiers gives a false sense of reproducibility.
\par\textit{Tradeoff: Impact analysis and auditability improve, while automated lineage can miss dynamic behavior. Pitfall: Lineage without immutable identifiers gives a false sense of reproducibility.}
\paragraph{LA-1110 [hard]} Explain Data lineage from first principles, then show how you would evaluate and operate it in production.
\noindent\textbf{Answer.} Start with the mechanism: Lineage links sources, transformations, datasets, features, prompts, models, evaluations, and deployments. Make the tradeoff explicit: Impact analysis and auditability improve, while automated lineage can miss dynamic behavior. Define workload-specific offline metrics and a latency/cost budget, test important slices, instrument the critical path, and create a rollback criterion. The production review must explicitly guard against this failure: Lineage without immutable identifiers gives a false sense of reproducibility.
\par\textit{A strong answer connects mechanism, measurable tradeoffs, failure isolation, observability, and rollback rather than listing product features.}
\paragraph{MCQ-1111 [medium]} Which statement best describes Reproducible environments?
\begin{enumerate}[label=\Alph*.]
\item Lockfiles, containers, hardware metadata, seeds, and deterministic settings constrain run variability.
\item DVC versions large data and model pointers with Git and defines reproducible pipeline stages and experiments.
\item Open table formats add schemas, snapshots, partition evolution, and transactional metadata over object storage.
\item Historical joins select only feature values known at each label timestamp.
\end{enumerate}
\noindent\textbf{Answer.} Lockfiles, containers, hardware metadata, seeds, and deterministic settings constrain run variability.
\par\textit{Lockfiles, containers, hardware metadata, seeds, and deterministic settings constrain run variability. Tradeoff: Reproducibility improves at the expense of build maintenance and sometimes performance. Watch for: A random seed alone cannot guarantee deterministic GPU execution.}
\paragraph{MCQ-1112 [hard]} What is the central engineering tradeoff of Reproducible environments?
\begin{enumerate}[label=\Alph*.]
\item Reliable analytics improve, while compaction and catalog operations become essential.
\item Familiar Git workflows improve reproducibility, while remote storage and cache discipline are required.
\item Leakage is prevented at the cost of temporal keys and more expensive joins.
\item Reproducibility improves at the expense of build maintenance and sometimes performance.
\end{enumerate}
\noindent\textbf{Answer.} Reproducibility improves at the expense of build maintenance and sometimes performance.
\par\textit{Lockfiles, containers, hardware metadata, seeds, and deterministic settings constrain run variability. Tradeoff: Reproducibility improves at the expense of build maintenance and sometimes performance. Watch for: A random seed alone cannot guarantee deterministic GPU execution.}
\paragraph{MCQ-1113 [hard]} Which failure mode should an interviewer expect you to identify for Reproducible environments?
\begin{enumerate}[label=\Alph*.]
\item Excess small files destroy scan performance despite correct table semantics.
\item Joining the latest feature value into old examples creates unrealistically strong training results.
\item A random seed alone cannot guarantee deterministic GPU execution.
\item Committing only code while silently replacing remote data breaks lineage.
\end{enumerate}
\noindent\textbf{Answer.} A random seed alone cannot guarantee deterministic GPU execution.
\par\textit{Lockfiles, containers, hardware metadata, seeds, and deterministic settings constrain run variability. Tradeoff: Reproducibility improves at the expense of build maintenance and sometimes performance. Watch for: A random seed alone cannot guarantee deterministic GPU execution.}
\paragraph{MCQ-1114 [medium]} A design review is specifically concerned with this risk: A random seed alone cannot guarantee deterministic GPU execution. Which topic should be reviewed first?
\begin{enumerate}[label=\Alph*.]
\item Reproducible environments
\item DVC
\item Delta Lake and Iceberg
\item Point-in-time correctness
\end{enumerate}
\noindent\textbf{Answer.} Reproducible environments
\par\textit{Lockfiles, containers, hardware metadata, seeds, and deterministic settings constrain run variability. Tradeoff: Reproducibility improves at the expense of build maintenance and sometimes performance. Watch for: A random seed alone cannot guarantee deterministic GPU execution.}
\paragraph{MCQ-1115 [hard]} Which design note most accurately balances the benefits and costs of Reproducible environments?
\begin{enumerate}[label=\Alph*.]
\item Leakage is prevented at the cost of temporal keys and more expensive joins.
\item Reproducibility improves at the expense of build maintenance and sometimes performance.
\item Familiar Git workflows improve reproducibility, while remote storage and cache discipline are required.
\item Reliable analytics improve, while compaction and catalog operations become essential.
\end{enumerate}
\noindent\textbf{Answer.} Reproducibility improves at the expense of build maintenance and sometimes performance.
\par\textit{Lockfiles, containers, hardware metadata, seeds, and deterministic settings constrain run variability. Tradeoff: Reproducibility improves at the expense of build maintenance and sometimes performance. Watch for: A random seed alone cannot guarantee deterministic GPU execution.}
\paragraph{MCQ-1116 [medium]} Which explanation would be strongest in a system-design interview about Reproducible environments?
\begin{enumerate}[label=\Alph*.]
\item Open table formats add schemas, snapshots, partition evolution, and transactional metadata over object storage.
\item Historical joins select only feature values known at each label timestamp.
\item Lockfiles, containers, hardware metadata, seeds, and deterministic settings constrain run variability.
\item DVC versions large data and model pointers with Git and defines reproducible pipeline stages and experiments.
\end{enumerate}
\noindent\textbf{Answer.} Lockfiles, containers, hardware metadata, seeds, and deterministic settings constrain run variability.
\par\textit{Lockfiles, containers, hardware metadata, seeds, and deterministic settings constrain run variability. Tradeoff: Reproducibility improves at the expense of build maintenance and sometimes performance. Watch for: A random seed alone cannot guarantee deterministic GPU execution.}
\paragraph{MCQ-1117 [hard]} Which production warning is most directly associated with Reproducible environments?
\begin{enumerate}[label=\Alph*.]
\item Excess small files destroy scan performance despite correct table semantics.
\item Committing only code while silently replacing remote data breaks lineage.
\item Joining the latest feature value into old examples creates unrealistically strong training results.
\item A random seed alone cannot guarantee deterministic GPU execution.
\end{enumerate}
\noindent\textbf{Answer.} A random seed alone cannot guarantee deterministic GPU execution.
\par\textit{Lockfiles, containers, hardware metadata, seeds, and deterministic settings constrain run variability. Tradeoff: Reproducibility improves at the expense of build maintenance and sometimes performance. Watch for: A random seed alone cannot guarantee deterministic GPU execution.}
\paragraph{FC-1118 [medium]} Define Reproducible environments in interview-ready language.
\noindent\textbf{Answer.} Lockfiles, containers, hardware metadata, seeds, and deterministic settings constrain run variability.
\par\textit{Lockfiles, containers, hardware metadata, seeds, and deterministic settings constrain run variability.}
\paragraph{FC-1119 [hard]} State the main tradeoff and production pitfall for Reproducible environments.
\noindent\textbf{Answer.} Tradeoff: Reproducibility improves at the expense of build maintenance and sometimes performance. Pitfall: A random seed alone cannot guarantee deterministic GPU execution.
\par\textit{Tradeoff: Reproducibility improves at the expense of build maintenance and sometimes performance. Pitfall: A random seed alone cannot guarantee deterministic GPU execution.}
\paragraph{LA-1120 [hard]} Explain Reproducible environments from first principles, then show how you would evaluate and operate it in production.
\noindent\textbf{Answer.} Start with the mechanism: Lockfiles, containers, hardware metadata, seeds, and deterministic settings constrain run variability. Make the tradeoff explicit: Reproducibility improves at the expense of build maintenance and sometimes performance. Define workload-specific offline metrics and a latency/cost budget, test important slices, instrument the critical path, and create a rollback criterion. The production review must explicitly guard against this failure: A random seed alone cannot guarantee deterministic GPU execution.
\par\textit{A strong answer connects mechanism, measurable tradeoffs, failure isolation, observability, and rollback rather than listing product features.}
\subsection{Transformer Theory and Training}
\paragraph{MCQ-1121 [medium]} Which statement best describes Tokenization?
\begin{enumerate}[label=\Alph*.]
\item Tokenizers map text or bytes into model vocabulary IDs using algorithms such as BPE, WordPiece, or unigram models.
\item Position signals break permutation invariance using absolute, relative, rotary, or bias-based schemes.
\item Direct Preference Optimization fits preferred over rejected responses through a classification-style objective relative to a reference policy.
\item Multiple heads project attention into separate subspaces before concatenation and output projection.
\end{enumerate}
\noindent\textbf{Answer.} Tokenizers map text or bytes into model vocabulary IDs using algorithms such as BPE, WordPiece, or unigram models.
\par\textit{Tokenizers map text or bytes into model vocabulary IDs using algorithms such as BPE, WordPiece, or unigram models. Tradeoff: Larger vocabularies shorten sequences but increase embedding parameters and sparse coverage. Watch for: Comparing context lengths across tokenizers as if tokens were equal misstates capacity and cost.}
\paragraph{MCQ-1122 [hard]} What is the central engineering tradeoff of Tokenization?
\begin{enumerate}[label=\Alph*.]
\item It simplifies preference training but still depends on pair quality and distribution coverage.
\item Larger vocabularies shorten sequences but increase embedding parameters and sparse coverage.
\item Diverse interactions improve capacity but add projections and KV memory.
\item Relative methods extrapolate differently, while implementation and scaling choices affect long context.
\end{enumerate}
\noindent\textbf{Answer.} Larger vocabularies shorten sequences but increase embedding parameters and sparse coverage.
\par\textit{Tokenizers map text or bytes into model vocabulary IDs using algorithms such as BPE, WordPiece, or unigram models. Tradeoff: Larger vocabularies shorten sequences but increase embedding parameters and sparse coverage. Watch for: Comparing context lengths across tokenizers as if tokens were equal misstates capacity and cost.}
\paragraph{MCQ-1123 [hard]} Which failure mode should an interviewer expect you to identify for Tokenization?
\begin{enumerate}[label=\Alph*.]
\item Comparing context lengths across tokenizers as if tokens were equal misstates capacity and cost.
\item Extending context without validating the position scheme can collapse quality.
\item Assuming heads are independently interpretable can overstate mechanistic conclusions.
\item Treating implicit user choices as clean preferences introduces confounding.
\end{enumerate}
\noindent\textbf{Answer.} Comparing context lengths across tokenizers as if tokens were equal misstates capacity and cost.
\par\textit{Tokenizers map text or bytes into model vocabulary IDs using algorithms such as BPE, WordPiece, or unigram models. Tradeoff: Larger vocabularies shorten sequences but increase embedding parameters and sparse coverage. Watch for: Comparing context lengths across tokenizers as if tokens were equal misstates capacity and cost.}
\paragraph{MCQ-1124 [medium]} A design review is specifically concerned with this risk: Comparing context lengths across tokenizers as if tokens were equal misstates capacity and cost. Which topic should be reviewed first?
\begin{enumerate}[label=\Alph*.]
\item DPO
\item Positional encoding
\item Multi-head attention
\item Tokenization
\end{enumerate}
\noindent\textbf{Answer.} Tokenization
\par\textit{Tokenizers map text or bytes into model vocabulary IDs using algorithms such as BPE, WordPiece, or unigram models. Tradeoff: Larger vocabularies shorten sequences but increase embedding parameters and sparse coverage. Watch for: Comparing context lengths across tokenizers as if tokens were equal misstates capacity and cost.}
\paragraph{MCQ-1125 [hard]} Which design note most accurately balances the benefits and costs of Tokenization?
\begin{enumerate}[label=\Alph*.]
\item Diverse interactions improve capacity but add projections and KV memory.
\item It simplifies preference training but still depends on pair quality and distribution coverage.
\item Larger vocabularies shorten sequences but increase embedding parameters and sparse coverage.
\item Relative methods extrapolate differently, while implementation and scaling choices affect long context.
\end{enumerate}
\noindent\textbf{Answer.} Larger vocabularies shorten sequences but increase embedding parameters and sparse coverage.
\par\textit{Tokenizers map text or bytes into model vocabulary IDs using algorithms such as BPE, WordPiece, or unigram models. Tradeoff: Larger vocabularies shorten sequences but increase embedding parameters and sparse coverage. Watch for: Comparing context lengths across tokenizers as if tokens were equal misstates capacity and cost.}
\paragraph{MCQ-1126 [medium]} Which explanation would be strongest in a system-design interview about Tokenization?
\begin{enumerate}[label=\Alph*.]
\item Direct Preference Optimization fits preferred over rejected responses through a classification-style objective relative to a reference policy.
\item Multiple heads project attention into separate subspaces before concatenation and output projection.
\item Position signals break permutation invariance using absolute, relative, rotary, or bias-based schemes.
\item Tokenizers map text or bytes into model vocabulary IDs using algorithms such as BPE, WordPiece, or unigram models.
\end{enumerate}
\noindent\textbf{Answer.} Tokenizers map text or bytes into model vocabulary IDs using algorithms such as BPE, WordPiece, or unigram models.
\par\textit{Tokenizers map text or bytes into model vocabulary IDs using algorithms such as BPE, WordPiece, or unigram models. Tradeoff: Larger vocabularies shorten sequences but increase embedding parameters and sparse coverage. Watch for: Comparing context lengths across tokenizers as if tokens were equal misstates capacity and cost.}
\paragraph{MCQ-1127 [hard]} Which production warning is most directly associated with Tokenization?
\begin{enumerate}[label=\Alph*.]
\item Extending context without validating the position scheme can collapse quality.
\item Assuming heads are independently interpretable can overstate mechanistic conclusions.
\item Treating implicit user choices as clean preferences introduces confounding.
\item Comparing context lengths across tokenizers as if tokens were equal misstates capacity and cost.
\end{enumerate}
\noindent\textbf{Answer.} Comparing context lengths across tokenizers as if tokens were equal misstates capacity and cost.
\par\textit{Tokenizers map text or bytes into model vocabulary IDs using algorithms such as BPE, WordPiece, or unigram models. Tradeoff: Larger vocabularies shorten sequences but increase embedding parameters and sparse coverage. Watch for: Comparing context lengths across tokenizers as if tokens were equal misstates capacity and cost.}
\paragraph{FC-1128 [medium]} Define Tokenization in interview-ready language.
\noindent\textbf{Answer.} Tokenizers map text or bytes into model vocabulary IDs using algorithms such as BPE, WordPiece, or unigram models.
\par\textit{Tokenizers map text or bytes into model vocabulary IDs using algorithms such as BPE, WordPiece, or unigram models.}
\paragraph{FC-1129 [hard]} State the main tradeoff and production pitfall for Tokenization.
\noindent\textbf{Answer.} Tradeoff: Larger vocabularies shorten sequences but increase embedding parameters and sparse coverage. Pitfall: Comparing context lengths across tokenizers as if tokens were equal misstates capacity and cost.
\par\textit{Tradeoff: Larger vocabularies shorten sequences but increase embedding parameters and sparse coverage. Pitfall: Comparing context lengths across tokenizers as if tokens were equal misstates capacity and cost.}
\paragraph{LA-1130 [hard]} Explain Tokenization from first principles, then show how you would evaluate and operate it in production.
\noindent\textbf{Answer.} Start with the mechanism: Tokenizers map text or bytes into model vocabulary IDs using algorithms such as BPE, WordPiece, or unigram models. Make the tradeoff explicit: Larger vocabularies shorten sequences but increase embedding parameters and sparse coverage. Define workload-specific offline metrics and a latency/cost budget, test important slices, instrument the critical path, and create a rollback criterion. The production review must explicitly guard against this failure: Comparing context lengths across tokenizers as if tokens were equal misstates capacity and cost.
\par\textit{A strong answer connects mechanism, measurable tradeoffs, failure isolation, observability, and rollback rather than listing product features.}
\paragraph{MCQ-1131 [medium]} Which statement best describes Embedding layer?
\begin{enumerate}[label=\Alph*.]
\item MoE routes each token to a subset of expert networks, increasing parameters without proportional per-token compute.
\item Learned token vectors map discrete IDs into continuous model space and are often tied to output weights.
\item Residual paths add transformed activations back to the stream to improve gradient flow and feature reuse.
\item Transformer MLP blocks expand hidden dimensions, apply nonlinear gates, and project back per token.
\end{enumerate}
\noindent\textbf{Answer.} Learned token vectors map discrete IDs into continuous model space and are often tied to output weights.
\par\textit{Learned token vectors map discrete IDs into continuous model space and are often tied to output weights. Tradeoff: Weight tying saves parameters but constrains input-output geometry. Watch for: Forgetting vocabulary changes invalidates checkpoint embedding shapes.}
\paragraph{MCQ-1132 [hard]} What is the central engineering tradeoff of Embedding layer?
\begin{enumerate}[label=\Alph*.]
\item Weight tying saves parameters but constrains input-output geometry.
\item They hold substantial parameters and compute while offering parallel token processing.
\item Capacity scales efficiently but routing balance, communication, and memory are hard.
\item Deep optimization becomes feasible, while residual scale can grow and require controls.
\end{enumerate}
\noindent\textbf{Answer.} Weight tying saves parameters but constrains input-output geometry.
\par\textit{Learned token vectors map discrete IDs into continuous model space and are often tied to output weights. Tradeoff: Weight tying saves parameters but constrains input-output geometry. Watch for: Forgetting vocabulary changes invalidates checkpoint embedding shapes.}
\paragraph{MCQ-1133 [hard]} Which failure mode should an interviewer expect you to identify for Embedding layer?
\begin{enumerate}[label=\Alph*.]
\item In-place mutation of residual tensors can break autograd or checkpointing.
\item Ignoring activation memory underestimates training capacity.
\item Forgetting vocabulary changes invalidates checkpoint embedding shapes.
\item Expert collapse and uneven routing create stragglers and weak specialization.
\end{enumerate}
\noindent\textbf{Answer.} Forgetting vocabulary changes invalidates checkpoint embedding shapes.
\par\textit{Learned token vectors map discrete IDs into continuous model space and are often tied to output weights. Tradeoff: Weight tying saves parameters but constrains input-output geometry. Watch for: Forgetting vocabulary changes invalidates checkpoint embedding shapes.}
\paragraph{MCQ-1134 [medium]} A design review is specifically concerned with this risk: Forgetting vocabulary changes invalidates checkpoint embedding shapes. Which topic should be reviewed first?
\begin{enumerate}[label=\Alph*.]
\item Residual connections
\item Mixture of Experts
\item Embedding layer
\item Feed-forward networks
\end{enumerate}
\noindent\textbf{Answer.} Embedding layer
\par\textit{Learned token vectors map discrete IDs into continuous model space and are often tied to output weights. Tradeoff: Weight tying saves parameters but constrains input-output geometry. Watch for: Forgetting vocabulary changes invalidates checkpoint embedding shapes.}
\paragraph{MCQ-1135 [hard]} Which design note most accurately balances the benefits and costs of Embedding layer?
\begin{enumerate}[label=\Alph*.]
\item Capacity scales efficiently but routing balance, communication, and memory are hard.
\item Deep optimization becomes feasible, while residual scale can grow and require controls.
\item Weight tying saves parameters but constrains input-output geometry.
\item They hold substantial parameters and compute while offering parallel token processing.
\end{enumerate}
\noindent\textbf{Answer.} Weight tying saves parameters but constrains input-output geometry.
\par\textit{Learned token vectors map discrete IDs into continuous model space and are often tied to output weights. Tradeoff: Weight tying saves parameters but constrains input-output geometry. Watch for: Forgetting vocabulary changes invalidates checkpoint embedding shapes.}
\paragraph{MCQ-1136 [medium]} Which explanation would be strongest in a system-design interview about Embedding layer?
\begin{enumerate}[label=\Alph*.]
\item Learned token vectors map discrete IDs into continuous model space and are often tied to output weights.
\item MoE routes each token to a subset of expert networks, increasing parameters without proportional per-token compute.
\item Residual paths add transformed activations back to the stream to improve gradient flow and feature reuse.
\item Transformer MLP blocks expand hidden dimensions, apply nonlinear gates, and project back per token.
\end{enumerate}
\noindent\textbf{Answer.} Learned token vectors map discrete IDs into continuous model space and are often tied to output weights.
\par\textit{Learned token vectors map discrete IDs into continuous model space and are often tied to output weights. Tradeoff: Weight tying saves parameters but constrains input-output geometry. Watch for: Forgetting vocabulary changes invalidates checkpoint embedding shapes.}
\paragraph{MCQ-1137 [hard]} Which production warning is most directly associated with Embedding layer?
\begin{enumerate}[label=\Alph*.]
\item Expert collapse and uneven routing create stragglers and weak specialization.
\item In-place mutation of residual tensors can break autograd or checkpointing.
\item Forgetting vocabulary changes invalidates checkpoint embedding shapes.
\item Ignoring activation memory underestimates training capacity.
\end{enumerate}
\noindent\textbf{Answer.} Forgetting vocabulary changes invalidates checkpoint embedding shapes.
\par\textit{Learned token vectors map discrete IDs into continuous model space and are often tied to output weights. Tradeoff: Weight tying saves parameters but constrains input-output geometry. Watch for: Forgetting vocabulary changes invalidates checkpoint embedding shapes.}
\paragraph{FC-1138 [medium]} Define Embedding layer in interview-ready language.
\noindent\textbf{Answer.} Learned token vectors map discrete IDs into continuous model space and are often tied to output weights.
\par\textit{Learned token vectors map discrete IDs into continuous model space and are often tied to output weights.}
\paragraph{FC-1139 [hard]} State the main tradeoff and production pitfall for Embedding layer.
\noindent\textbf{Answer.} Tradeoff: Weight tying saves parameters but constrains input-output geometry. Pitfall: Forgetting vocabulary changes invalidates checkpoint embedding shapes.
\par\textit{Tradeoff: Weight tying saves parameters but constrains input-output geometry. Pitfall: Forgetting vocabulary changes invalidates checkpoint embedding shapes.}
\paragraph{LA-1140 [hard]} Explain Embedding layer from first principles, then show how you would evaluate and operate it in production.
\noindent\textbf{Answer.} Start with the mechanism: Learned token vectors map discrete IDs into continuous model space and are often tied to output weights. Make the tradeoff explicit: Weight tying saves parameters but constrains input-output geometry. Define workload-specific offline metrics and a latency/cost budget, test important slices, instrument the critical path, and create a rollback criterion. The production review must explicitly guard against this failure: Forgetting vocabulary changes invalidates checkpoint embedding shapes.
\par\textit{A strong answer connects mechanism, measurable tradeoffs, failure isolation, observability, and rollback rather than listing product features.}
\paragraph{MCQ-1141 [medium]} Which statement best describes Self-attention math?
\begin{enumerate}[label=\Alph*.]
\item Normalization stabilizes activations; RMSNorm removes mean centering and normalizes by root-mean-square magnitude.
\item A preference model supplies rewards while PPO constrains policy updates relative to a reference model.
\item MoE routes each token to a subset of expert networks, increasing parameters without proportional per-token compute.
\item Attention computes scaled query-key scores, applies masking and softmax, then mixes value vectors.
\end{enumerate}
\noindent\textbf{Answer.} Attention computes scaled query-key scores, applies masking and softmax, then mixes value vectors.
\par\textit{Attention computes scaled query-key scores, applies masking and softmax, then mixes value vectors. Tradeoff: Every token can interact directly, while standard sequence cost grows quadratically. Watch for: Omitting the square-root head-dimension scale saturates softmax gradients.}
\paragraph{MCQ-1142 [hard]} What is the central engineering tradeoff of Self-attention math?
\begin{enumerate}[label=\Alph*.]
\item Every token can interact directly, while standard sequence cost grows quadratically.
\item Capacity scales efficiently but routing balance, communication, and memory are hard.
\item It can optimize human preferences but adds reward hacking and training complexity.
\item Simpler normalization can be faster but architecture placement changes optimization dynamics.
\end{enumerate}
\noindent\textbf{Answer.} Every token can interact directly, while standard sequence cost grows quadratically.
\par\textit{Attention computes scaled query-key scores, applies masking and softmax, then mixes value vectors. Tradeoff: Every token can interact directly, while standard sequence cost grows quadratically. Watch for: Omitting the square-root head-dimension scale saturates softmax gradients.}
\paragraph{MCQ-1143 [hard]} Which failure mode should an interviewer expect you to identify for Self-attention math?
\begin{enumerate}[label=\Alph*.]
\item Expert collapse and uneven routing create stragglers and weak specialization.
\item Omitting the square-root head-dimension scale saturates softmax gradients.
\item A biased reward model can produce confidently optimized failure modes.
\item Moving pre-norm to post-norm without retuning can destabilize deep training.
\end{enumerate}
\noindent\textbf{Answer.} Omitting the square-root head-dimension scale saturates softmax gradients.
\par\textit{Attention computes scaled query-key scores, applies masking and softmax, then mixes value vectors. Tradeoff: Every token can interact directly, while standard sequence cost grows quadratically. Watch for: Omitting the square-root head-dimension scale saturates softmax gradients.}
\paragraph{MCQ-1144 [medium]} A design review is specifically concerned with this risk: Omitting the square-root head-dimension scale saturates softmax gradients. Which topic should be reviewed first?
\begin{enumerate}[label=\Alph*.]
\item Self-attention math
\item Layer normalization and RMSNorm
\item RLHF and PPO
\item Mixture of Experts
\end{enumerate}
\noindent\textbf{Answer.} Self-attention math
\par\textit{Attention computes scaled query-key scores, applies masking and softmax, then mixes value vectors. Tradeoff: Every token can interact directly, while standard sequence cost grows quadratically. Watch for: Omitting the square-root head-dimension scale saturates softmax gradients.}
\paragraph{MCQ-1145 [hard]} Which design note most accurately balances the benefits and costs of Self-attention math?
\begin{enumerate}[label=\Alph*.]
\item Every token can interact directly, while standard sequence cost grows quadratically.
\item It can optimize human preferences but adds reward hacking and training complexity.
\item Simpler normalization can be faster but architecture placement changes optimization dynamics.
\item Capacity scales efficiently but routing balance, communication, and memory are hard.
\end{enumerate}
\noindent\textbf{Answer.} Every token can interact directly, while standard sequence cost grows quadratically.
\par\textit{Attention computes scaled query-key scores, applies masking and softmax, then mixes value vectors. Tradeoff: Every token can interact directly, while standard sequence cost grows quadratically. Watch for: Omitting the square-root head-dimension scale saturates softmax gradients.}
\paragraph{MCQ-1146 [medium]} Which explanation would be strongest in a system-design interview about Self-attention math?
\begin{enumerate}[label=\Alph*.]
\item A preference model supplies rewards while PPO constrains policy updates relative to a reference model.
\item Attention computes scaled query-key scores, applies masking and softmax, then mixes value vectors.
\item MoE routes each token to a subset of expert networks, increasing parameters without proportional per-token compute.
\item Normalization stabilizes activations; RMSNorm removes mean centering and normalizes by root-mean-square magnitude.
\end{enumerate}
\noindent\textbf{Answer.} Attention computes scaled query-key scores, applies masking and softmax, then mixes value vectors.
\par\textit{Attention computes scaled query-key scores, applies masking and softmax, then mixes value vectors. Tradeoff: Every token can interact directly, while standard sequence cost grows quadratically. Watch for: Omitting the square-root head-dimension scale saturates softmax gradients.}
\paragraph{MCQ-1147 [hard]} Which production warning is most directly associated with Self-attention math?
\begin{enumerate}[label=\Alph*.]
\item Omitting the square-root head-dimension scale saturates softmax gradients.
\item Moving pre-norm to post-norm without retuning can destabilize deep training.
\item A biased reward model can produce confidently optimized failure modes.
\item Expert collapse and uneven routing create stragglers and weak specialization.
\end{enumerate}
\noindent\textbf{Answer.} Omitting the square-root head-dimension scale saturates softmax gradients.
\par\textit{Attention computes scaled query-key scores, applies masking and softmax, then mixes value vectors. Tradeoff: Every token can interact directly, while standard sequence cost grows quadratically. Watch for: Omitting the square-root head-dimension scale saturates softmax gradients.}
\paragraph{FC-1148 [medium]} Define Self-attention math in interview-ready language.
\noindent\textbf{Answer.} Attention computes scaled query-key scores, applies masking and softmax, then mixes value vectors.
\par\textit{Attention computes scaled query-key scores, applies masking and softmax, then mixes value vectors.}
\paragraph{FC-1149 [hard]} State the main tradeoff and production pitfall for Self-attention math.
\noindent\textbf{Answer.} Tradeoff: Every token can interact directly, while standard sequence cost grows quadratically. Pitfall: Omitting the square-root head-dimension scale saturates softmax gradients.
\par\textit{Tradeoff: Every token can interact directly, while standard sequence cost grows quadratically. Pitfall: Omitting the square-root head-dimension scale saturates softmax gradients.}
\paragraph{LA-1150 [hard]} Explain Self-attention math from first principles, then show how you would evaluate and operate it in production.
\noindent\textbf{Answer.} Start with the mechanism: Attention computes scaled query-key scores, applies masking and softmax, then mixes value vectors. Make the tradeoff explicit: Every token can interact directly, while standard sequence cost grows quadratically. Define workload-specific offline metrics and a latency/cost budget, test important slices, instrument the critical path, and create a rollback criterion. The production review must explicitly guard against this failure: Omitting the square-root head-dimension scale saturates softmax gradients.
\par\textit{A strong answer connects mechanism, measurable tradeoffs, failure isolation, observability, and rollback rather than listing product features.}
\paragraph{MCQ-1151 [medium]} Which statement best describes Multi-head attention?
\begin{enumerate}[label=\Alph*.]
\item Transformer MLP blocks expand hidden dimensions, apply nonlinear gates, and project back per token.
\item LoRA trains low-rank update matrices while freezing base weights for parameter-efficient adaptation.
\item Next-token training factorizes sequence likelihood from left to right under a causal mask.
\item Multiple heads project attention into separate subspaces before concatenation and output projection.
\end{enumerate}
\noindent\textbf{Answer.} Multiple heads project attention into separate subspaces before concatenation and output projection.
\par\textit{Multiple heads project attention into separate subspaces before concatenation and output projection. Tradeoff: Diverse interactions improve capacity but add projections and KV memory. Watch for: Assuming heads are independently interpretable can overstate mechanistic conclusions.}
\paragraph{MCQ-1152 [hard]} What is the central engineering tradeoff of Multi-head attention?
\begin{enumerate}[label=\Alph*.]
\item They hold substantial parameters and compute while offering parallel token processing.
\item Memory and storage fall, while rank and target-module choices limit expressiveness.
\item Diverse interactions improve capacity but add projections and KV memory.
\item It provides a simple scalable objective but does not directly optimize instruction usefulness or truth.
\end{enumerate}
\noindent\textbf{Answer.} Diverse interactions improve capacity but add projections and KV memory.
\par\textit{Multiple heads project attention into separate subspaces before concatenation and output projection. Tradeoff: Diverse interactions improve capacity but add projections and KV memory. Watch for: Assuming heads are independently interpretable can overstate mechanistic conclusions.}
\paragraph{MCQ-1153 [hard]} Which failure mode should an interviewer expect you to identify for Multi-head attention?
\begin{enumerate}[label=\Alph*.]
\item Ignoring activation memory underestimates training capacity.
\item Merging incompatible adapters without evaluation creates interference.
\item Assuming heads are independently interpretable can overstate mechanistic conclusions.
\item Interpreting low perplexity as safe task performance is invalid.
\end{enumerate}
\noindent\textbf{Answer.} Assuming heads are independently interpretable can overstate mechanistic conclusions.
\par\textit{Multiple heads project attention into separate subspaces before concatenation and output projection. Tradeoff: Diverse interactions improve capacity but add projections and KV memory. Watch for: Assuming heads are independently interpretable can overstate mechanistic conclusions.}
\paragraph{MCQ-1154 [medium]} A design review is specifically concerned with this risk: Assuming heads are independently interpretable can overstate mechanistic conclusions. Which topic should be reviewed first?
\begin{enumerate}[label=\Alph*.]
\item Feed-forward networks
\item Multi-head attention
\item LoRA
\item Causal language modeling
\end{enumerate}
\noindent\textbf{Answer.} Multi-head attention
\par\textit{Multiple heads project attention into separate subspaces before concatenation and output projection. Tradeoff: Diverse interactions improve capacity but add projections and KV memory. Watch for: Assuming heads are independently interpretable can overstate mechanistic conclusions.}
\paragraph{MCQ-1155 [hard]} Which design note most accurately balances the benefits and costs of Multi-head attention?
\begin{enumerate}[label=\Alph*.]
\item It provides a simple scalable objective but does not directly optimize instruction usefulness or truth.
\item Diverse interactions improve capacity but add projections and KV memory.
\item Memory and storage fall, while rank and target-module choices limit expressiveness.
\item They hold substantial parameters and compute while offering parallel token processing.
\end{enumerate}
\noindent\textbf{Answer.} Diverse interactions improve capacity but add projections and KV memory.
\par\textit{Multiple heads project attention into separate subspaces before concatenation and output projection. Tradeoff: Diverse interactions improve capacity but add projections and KV memory. Watch for: Assuming heads are independently interpretable can overstate mechanistic conclusions.}
\paragraph{MCQ-1156 [medium]} Which explanation would be strongest in a system-design interview about Multi-head attention?
\begin{enumerate}[label=\Alph*.]
\item Next-token training factorizes sequence likelihood from left to right under a causal mask.
\item Multiple heads project attention into separate subspaces before concatenation and output projection.
\item LoRA trains low-rank update matrices while freezing base weights for parameter-efficient adaptation.
\item Transformer MLP blocks expand hidden dimensions, apply nonlinear gates, and project back per token.
\end{enumerate}
\noindent\textbf{Answer.} Multiple heads project attention into separate subspaces before concatenation and output projection.
\par\textit{Multiple heads project attention into separate subspaces before concatenation and output projection. Tradeoff: Diverse interactions improve capacity but add projections and KV memory. Watch for: Assuming heads are independently interpretable can overstate mechanistic conclusions.}
\paragraph{MCQ-1157 [hard]} Which production warning is most directly associated with Multi-head attention?
\begin{enumerate}[label=\Alph*.]
\item Merging incompatible adapters without evaluation creates interference.
\item Assuming heads are independently interpretable can overstate mechanistic conclusions.
\item Interpreting low perplexity as safe task performance is invalid.
\item Ignoring activation memory underestimates training capacity.
\end{enumerate}
\noindent\textbf{Answer.} Assuming heads are independently interpretable can overstate mechanistic conclusions.
\par\textit{Multiple heads project attention into separate subspaces before concatenation and output projection. Tradeoff: Diverse interactions improve capacity but add projections and KV memory. Watch for: Assuming heads are independently interpretable can overstate mechanistic conclusions.}
\paragraph{FC-1158 [medium]} Define Multi-head attention in interview-ready language.
\noindent\textbf{Answer.} Multiple heads project attention into separate subspaces before concatenation and output projection.
\par\textit{Multiple heads project attention into separate subspaces before concatenation and output projection.}
\paragraph{FC-1159 [hard]} State the main tradeoff and production pitfall for Multi-head attention.
\noindent\textbf{Answer.} Tradeoff: Diverse interactions improve capacity but add projections and KV memory. Pitfall: Assuming heads are independently interpretable can overstate mechanistic conclusions.
\par\textit{Tradeoff: Diverse interactions improve capacity but add projections and KV memory. Pitfall: Assuming heads are independently interpretable can overstate mechanistic conclusions.}
\paragraph{LA-1160 [hard]} Explain Multi-head attention from first principles, then show how you would evaluate and operate it in production.
\noindent\textbf{Answer.} Start with the mechanism: Multiple heads project attention into separate subspaces before concatenation and output projection. Make the tradeoff explicit: Diverse interactions improve capacity but add projections and KV memory. Define workload-specific offline metrics and a latency/cost budget, test important slices, instrument the critical path, and create a rollback criterion. The production review must explicitly guard against this failure: Assuming heads are independently interpretable can overstate mechanistic conclusions.
\par\textit{A strong answer connects mechanism, measurable tradeoffs, failure isolation, observability, and rollback rather than listing product features.}
\paragraph{MCQ-1161 [medium]} Which statement best describes Positional encoding?
\begin{enumerate}[label=\Alph*.]
\item QLoRA backpropagates through a frozen quantized base model into LoRA adapters using memory-saving quantization techniques.
\item Rotary position embeddings rotate query and key coordinates so dot products encode relative offsets.
\item Position signals break permutation invariance using absolute, relative, rotary, or bias-based schemes.
\item Transformer MLP blocks expand hidden dimensions, apply nonlinear gates, and project back per token.
\end{enumerate}
\noindent\textbf{Answer.} Position signals break permutation invariance using absolute, relative, rotary, or bias-based schemes.
\par\textit{Position signals break permutation invariance using absolute, relative, rotary, or bias-based schemes. Tradeoff: Relative methods extrapolate differently, while implementation and scaling choices affect long context. Watch for: Extending context without validating the position scheme can collapse quality.}
\paragraph{MCQ-1162 [hard]} What is the central engineering tradeoff of Positional encoding?
\begin{enumerate}[label=\Alph*.]
\item They hold substantial parameters and compute while offering parallel token processing.
\item Relative methods extrapolate differently, while implementation and scaling choices affect long context.
\item Fine-tuning large models becomes cheaper, while kernels and quantization settings affect stability.
\item RoPE integrates cleanly with attention but long-context scaling needs careful frequency treatment.
\end{enumerate}
\noindent\textbf{Answer.} Relative methods extrapolate differently, while implementation and scaling choices affect long context.
\par\textit{Position signals break permutation invariance using absolute, relative, rotary, or bias-based schemes. Tradeoff: Relative methods extrapolate differently, while implementation and scaling choices affect long context. Watch for: Extending context without validating the position scheme can collapse quality.}
\paragraph{MCQ-1163 [hard]} Which failure mode should an interviewer expect you to identify for Positional encoding?
\begin{enumerate}[label=\Alph*.]
\item Naively increasing the maximum length does not preserve trained frequency behavior.
\item Ignoring activation memory underestimates training capacity.
\item Extending context without validating the position scheme can collapse quality.
\item Confusing training quantization with final serving quantization produces wrong expectations.
\end{enumerate}
\noindent\textbf{Answer.} Extending context without validating the position scheme can collapse quality.
\par\textit{Position signals break permutation invariance using absolute, relative, rotary, or bias-based schemes. Tradeoff: Relative methods extrapolate differently, while implementation and scaling choices affect long context. Watch for: Extending context without validating the position scheme can collapse quality.}
\paragraph{MCQ-1164 [medium]} A design review is specifically concerned with this risk: Extending context without validating the position scheme can collapse quality. Which topic should be reviewed first?
\begin{enumerate}[label=\Alph*.]
\item QLoRA
\item RoPE
\item Positional encoding
\item Feed-forward networks
\end{enumerate}
\noindent\textbf{Answer.} Positional encoding
\par\textit{Position signals break permutation invariance using absolute, relative, rotary, or bias-based schemes. Tradeoff: Relative methods extrapolate differently, while implementation and scaling choices affect long context. Watch for: Extending context without validating the position scheme can collapse quality.}
\paragraph{MCQ-1165 [hard]} Which design note most accurately balances the benefits and costs of Positional encoding?
\begin{enumerate}[label=\Alph*.]
\item Relative methods extrapolate differently, while implementation and scaling choices affect long context.
\item Fine-tuning large models becomes cheaper, while kernels and quantization settings affect stability.
\item RoPE integrates cleanly with attention but long-context scaling needs careful frequency treatment.
\item They hold substantial parameters and compute while offering parallel token processing.
\end{enumerate}
\noindent\textbf{Answer.} Relative methods extrapolate differently, while implementation and scaling choices affect long context.
\par\textit{Position signals break permutation invariance using absolute, relative, rotary, or bias-based schemes. Tradeoff: Relative methods extrapolate differently, while implementation and scaling choices affect long context. Watch for: Extending context without validating the position scheme can collapse quality.}
\paragraph{MCQ-1166 [medium]} Which explanation would be strongest in a system-design interview about Positional encoding?
\begin{enumerate}[label=\Alph*.]
\item Position signals break permutation invariance using absolute, relative, rotary, or bias-based schemes.
\item Rotary position embeddings rotate query and key coordinates so dot products encode relative offsets.
\item QLoRA backpropagates through a frozen quantized base model into LoRA adapters using memory-saving quantization techniques.
\item Transformer MLP blocks expand hidden dimensions, apply nonlinear gates, and project back per token.
\end{enumerate}
\noindent\textbf{Answer.} Position signals break permutation invariance using absolute, relative, rotary, or bias-based schemes.
\par\textit{Position signals break permutation invariance using absolute, relative, rotary, or bias-based schemes. Tradeoff: Relative methods extrapolate differently, while implementation and scaling choices affect long context. Watch for: Extending context without validating the position scheme can collapse quality.}
\paragraph{MCQ-1167 [hard]} Which production warning is most directly associated with Positional encoding?
\begin{enumerate}[label=\Alph*.]
\item Naively increasing the maximum length does not preserve trained frequency behavior.
\item Confusing training quantization with final serving quantization produces wrong expectations.
\item Extending context without validating the position scheme can collapse quality.
\item Ignoring activation memory underestimates training capacity.
\end{enumerate}
\noindent\textbf{Answer.} Extending context without validating the position scheme can collapse quality.
\par\textit{Position signals break permutation invariance using absolute, relative, rotary, or bias-based schemes. Tradeoff: Relative methods extrapolate differently, while implementation and scaling choices affect long context. Watch for: Extending context without validating the position scheme can collapse quality.}
\paragraph{FC-1168 [medium]} Define Positional encoding in interview-ready language.
\noindent\textbf{Answer.} Position signals break permutation invariance using absolute, relative, rotary, or bias-based schemes.
\par\textit{Position signals break permutation invariance using absolute, relative, rotary, or bias-based schemes.}
\paragraph{FC-1169 [hard]} State the main tradeoff and production pitfall for Positional encoding.
\noindent\textbf{Answer.} Tradeoff: Relative methods extrapolate differently, while implementation and scaling choices affect long context. Pitfall: Extending context without validating the position scheme can collapse quality.
\par\textit{Tradeoff: Relative methods extrapolate differently, while implementation and scaling choices affect long context. Pitfall: Extending context without validating the position scheme can collapse quality.}
\paragraph{LA-1170 [hard]} Explain Positional encoding from first principles, then show how you would evaluate and operate it in production.
\noindent\textbf{Answer.} Start with the mechanism: Position signals break permutation invariance using absolute, relative, rotary, or bias-based schemes. Make the tradeoff explicit: Relative methods extrapolate differently, while implementation and scaling choices affect long context. Define workload-specific offline metrics and a latency/cost budget, test important slices, instrument the critical path, and create a rollback criterion. The production review must explicitly guard against this failure: Extending context without validating the position scheme can collapse quality.
\par\textit{A strong answer connects mechanism, measurable tradeoffs, failure isolation, observability, and rollback rather than listing product features.}
\paragraph{MCQ-1171 [medium]} Which statement best describes RoPE?
\begin{enumerate}[label=\Alph*.]
\item Residual paths add transformed activations back to the stream to improve gradient flow and feature reuse.
\item Position signals break permutation invariance using absolute, relative, rotary, or bias-based schemes.
\item Tokenizers map text or bytes into model vocabulary IDs using algorithms such as BPE, WordPiece, or unigram models.
\item Rotary position embeddings rotate query and key coordinates so dot products encode relative offsets.
\end{enumerate}
\noindent\textbf{Answer.} Rotary position embeddings rotate query and key coordinates so dot products encode relative offsets.
\par\textit{Rotary position embeddings rotate query and key coordinates so dot products encode relative offsets. Tradeoff: RoPE integrates cleanly with attention but long-context scaling needs careful frequency treatment. Watch for: Naively increasing the maximum length does not preserve trained frequency behavior.}
\paragraph{MCQ-1172 [hard]} What is the central engineering tradeoff of RoPE?
\begin{enumerate}[label=\Alph*.]
\item Relative methods extrapolate differently, while implementation and scaling choices affect long context.
\item Larger vocabularies shorten sequences but increase embedding parameters and sparse coverage.
\item RoPE integrates cleanly with attention but long-context scaling needs careful frequency treatment.
\item Deep optimization becomes feasible, while residual scale can grow and require controls.
\end{enumerate}
\noindent\textbf{Answer.} RoPE integrates cleanly with attention but long-context scaling needs careful frequency treatment.
\par\textit{Rotary position embeddings rotate query and key coordinates so dot products encode relative offsets. Tradeoff: RoPE integrates cleanly with attention but long-context scaling needs careful frequency treatment. Watch for: Naively increasing the maximum length does not preserve trained frequency behavior.}
\paragraph{MCQ-1173 [hard]} Which failure mode should an interviewer expect you to identify for RoPE?
\begin{enumerate}[label=\Alph*.]
\item Comparing context lengths across tokenizers as if tokens were equal misstates capacity and cost.
\item Naively increasing the maximum length does not preserve trained frequency behavior.
\item In-place mutation of residual tensors can break autograd or checkpointing.
\item Extending context without validating the position scheme can collapse quality.
\end{enumerate}
\noindent\textbf{Answer.} Naively increasing the maximum length does not preserve trained frequency behavior.
\par\textit{Rotary position embeddings rotate query and key coordinates so dot products encode relative offsets. Tradeoff: RoPE integrates cleanly with attention but long-context scaling needs careful frequency treatment. Watch for: Naively increasing the maximum length does not preserve trained frequency behavior.}
\paragraph{MCQ-1174 [medium]} A design review is specifically concerned with this risk: Naively increasing the maximum length does not preserve trained frequency behavior. Which topic should be reviewed first?
\begin{enumerate}[label=\Alph*.]
\item Positional encoding
\item Residual connections
\item Tokenization
\item RoPE
\end{enumerate}
\noindent\textbf{Answer.} RoPE
\par\textit{Rotary position embeddings rotate query and key coordinates so dot products encode relative offsets. Tradeoff: RoPE integrates cleanly with attention but long-context scaling needs careful frequency treatment. Watch for: Naively increasing the maximum length does not preserve trained frequency behavior.}
\paragraph{MCQ-1175 [hard]} Which design note most accurately balances the benefits and costs of RoPE?
\begin{enumerate}[label=\Alph*.]
\item Larger vocabularies shorten sequences but increase embedding parameters and sparse coverage.
\item Relative methods extrapolate differently, while implementation and scaling choices affect long context.
\item Deep optimization becomes feasible, while residual scale can grow and require controls.
\item RoPE integrates cleanly with attention but long-context scaling needs careful frequency treatment.
\end{enumerate}
\noindent\textbf{Answer.} RoPE integrates cleanly with attention but long-context scaling needs careful frequency treatment.
\par\textit{Rotary position embeddings rotate query and key coordinates so dot products encode relative offsets. Tradeoff: RoPE integrates cleanly with attention but long-context scaling needs careful frequency treatment. Watch for: Naively increasing the maximum length does not preserve trained frequency behavior.}
\paragraph{MCQ-1176 [medium]} Which explanation would be strongest in a system-design interview about RoPE?
\begin{enumerate}[label=\Alph*.]
\item Residual paths add transformed activations back to the stream to improve gradient flow and feature reuse.
\item Rotary position embeddings rotate query and key coordinates so dot products encode relative offsets.
\item Tokenizers map text or bytes into model vocabulary IDs using algorithms such as BPE, WordPiece, or unigram models.
\item Position signals break permutation invariance using absolute, relative, rotary, or bias-based schemes.
\end{enumerate}
\noindent\textbf{Answer.} Rotary position embeddings rotate query and key coordinates so dot products encode relative offsets.
\par\textit{Rotary position embeddings rotate query and key coordinates so dot products encode relative offsets. Tradeoff: RoPE integrates cleanly with attention but long-context scaling needs careful frequency treatment. Watch for: Naively increasing the maximum length does not preserve trained frequency behavior.}
\paragraph{MCQ-1177 [hard]} Which production warning is most directly associated with RoPE?
\begin{enumerate}[label=\Alph*.]
\item In-place mutation of residual tensors can break autograd or checkpointing.
\item Naively increasing the maximum length does not preserve trained frequency behavior.
\item Comparing context lengths across tokenizers as if tokens were equal misstates capacity and cost.
\item Extending context without validating the position scheme can collapse quality.
\end{enumerate}
\noindent\textbf{Answer.} Naively increasing the maximum length does not preserve trained frequency behavior.
\par\textit{Rotary position embeddings rotate query and key coordinates so dot products encode relative offsets. Tradeoff: RoPE integrates cleanly with attention but long-context scaling needs careful frequency treatment. Watch for: Naively increasing the maximum length does not preserve trained frequency behavior.}
\paragraph{FC-1178 [medium]} Define RoPE in interview-ready language.
\noindent\textbf{Answer.} Rotary position embeddings rotate query and key coordinates so dot products encode relative offsets.
\par\textit{Rotary position embeddings rotate query and key coordinates so dot products encode relative offsets.}
\paragraph{FC-1179 [hard]} State the main tradeoff and production pitfall for RoPE.
\noindent\textbf{Answer.} Tradeoff: RoPE integrates cleanly with attention but long-context scaling needs careful frequency treatment. Pitfall: Naively increasing the maximum length does not preserve trained frequency behavior.
\par\textit{Tradeoff: RoPE integrates cleanly with attention but long-context scaling needs careful frequency treatment. Pitfall: Naively increasing the maximum length does not preserve trained frequency behavior.}
\paragraph{LA-1180 [hard]} Explain RoPE from first principles, then show how you would evaluate and operate it in production.
\noindent\textbf{Answer.} Start with the mechanism: Rotary position embeddings rotate query and key coordinates so dot products encode relative offsets. Make the tradeoff explicit: RoPE integrates cleanly with attention but long-context scaling needs careful frequency treatment. Define workload-specific offline metrics and a latency/cost budget, test important slices, instrument the critical path, and create a rollback criterion. The production review must explicitly guard against this failure: Naively increasing the maximum length does not preserve trained frequency behavior.
\par\textit{A strong answer connects mechanism, measurable tradeoffs, failure isolation, observability, and rollback rather than listing product features.}
\paragraph{MCQ-1181 [medium]} Which statement best describes Feed-forward networks?
\begin{enumerate}[label=\Alph*.]
\item MoE routes each token to a subset of expert networks, increasing parameters without proportional per-token compute.
\item Rotary position embeddings rotate query and key coordinates so dot products encode relative offsets.
\item LoRA trains low-rank update matrices while freezing base weights for parameter-efficient adaptation.
\item Transformer MLP blocks expand hidden dimensions, apply nonlinear gates, and project back per token.
\end{enumerate}
\noindent\textbf{Answer.} Transformer MLP blocks expand hidden dimensions, apply nonlinear gates, and project back per token.
\par\textit{Transformer MLP blocks expand hidden dimensions, apply nonlinear gates, and project back per token. Tradeoff: They hold substantial parameters and compute while offering parallel token processing. Watch for: Ignoring activation memory underestimates training capacity.}
\paragraph{MCQ-1182 [hard]} What is the central engineering tradeoff of Feed-forward networks?
\begin{enumerate}[label=\Alph*.]
\item Capacity scales efficiently but routing balance, communication, and memory are hard.
\item RoPE integrates cleanly with attention but long-context scaling needs careful frequency treatment.
\item They hold substantial parameters and compute while offering parallel token processing.
\item Memory and storage fall, while rank and target-module choices limit expressiveness.
\end{enumerate}
\noindent\textbf{Answer.} They hold substantial parameters and compute while offering parallel token processing.
\par\textit{Transformer MLP blocks expand hidden dimensions, apply nonlinear gates, and project back per token. Tradeoff: They hold substantial parameters and compute while offering parallel token processing. Watch for: Ignoring activation memory underestimates training capacity.}
\paragraph{MCQ-1183 [hard]} Which failure mode should an interviewer expect you to identify for Feed-forward networks?
\begin{enumerate}[label=\Alph*.]
\item Merging incompatible adapters without evaluation creates interference.
\item Expert collapse and uneven routing create stragglers and weak specialization.
\item Ignoring activation memory underestimates training capacity.
\item Naively increasing the maximum length does not preserve trained frequency behavior.
\end{enumerate}
\noindent\textbf{Answer.} Ignoring activation memory underestimates training capacity.
\par\textit{Transformer MLP blocks expand hidden dimensions, apply nonlinear gates, and project back per token. Tradeoff: They hold substantial parameters and compute while offering parallel token processing. Watch for: Ignoring activation memory underestimates training capacity.}
\paragraph{MCQ-1184 [medium]} A design review is specifically concerned with this risk: Ignoring activation memory underestimates training capacity. Which topic should be reviewed first?
\begin{enumerate}[label=\Alph*.]
\item RoPE
\item Mixture of Experts
\item Feed-forward networks
\item LoRA
\end{enumerate}
\noindent\textbf{Answer.} Feed-forward networks
\par\textit{Transformer MLP blocks expand hidden dimensions, apply nonlinear gates, and project back per token. Tradeoff: They hold substantial parameters and compute while offering parallel token processing. Watch for: Ignoring activation memory underestimates training capacity.}
\paragraph{MCQ-1185 [hard]} Which design note most accurately balances the benefits and costs of Feed-forward networks?
\begin{enumerate}[label=\Alph*.]
\item They hold substantial parameters and compute while offering parallel token processing.
\item RoPE integrates cleanly with attention but long-context scaling needs careful frequency treatment.
\item Memory and storage fall, while rank and target-module choices limit expressiveness.
\item Capacity scales efficiently but routing balance, communication, and memory are hard.
\end{enumerate}
\noindent\textbf{Answer.} They hold substantial parameters and compute while offering parallel token processing.
\par\textit{Transformer MLP blocks expand hidden dimensions, apply nonlinear gates, and project back per token. Tradeoff: They hold substantial parameters and compute while offering parallel token processing. Watch for: Ignoring activation memory underestimates training capacity.}
\paragraph{MCQ-1186 [medium]} Which explanation would be strongest in a system-design interview about Feed-forward networks?
\begin{enumerate}[label=\Alph*.]
\item Transformer MLP blocks expand hidden dimensions, apply nonlinear gates, and project back per token.
\item MoE routes each token to a subset of expert networks, increasing parameters without proportional per-token compute.
\item LoRA trains low-rank update matrices while freezing base weights for parameter-efficient adaptation.
\item Rotary position embeddings rotate query and key coordinates so dot products encode relative offsets.
\end{enumerate}
\noindent\textbf{Answer.} Transformer MLP blocks expand hidden dimensions, apply nonlinear gates, and project back per token.
\par\textit{Transformer MLP blocks expand hidden dimensions, apply nonlinear gates, and project back per token. Tradeoff: They hold substantial parameters and compute while offering parallel token processing. Watch for: Ignoring activation memory underestimates training capacity.}
\paragraph{MCQ-1187 [hard]} Which production warning is most directly associated with Feed-forward networks?
\begin{enumerate}[label=\Alph*.]
\item Ignoring activation memory underestimates training capacity.
\item Naively increasing the maximum length does not preserve trained frequency behavior.
\item Merging incompatible adapters without evaluation creates interference.
\item Expert collapse and uneven routing create stragglers and weak specialization.
\end{enumerate}
\noindent\textbf{Answer.} Ignoring activation memory underestimates training capacity.
\par\textit{Transformer MLP blocks expand hidden dimensions, apply nonlinear gates, and project back per token. Tradeoff: They hold substantial parameters and compute while offering parallel token processing. Watch for: Ignoring activation memory underestimates training capacity.}
\paragraph{FC-1188 [medium]} Define Feed-forward networks in interview-ready language.
\noindent\textbf{Answer.} Transformer MLP blocks expand hidden dimensions, apply nonlinear gates, and project back per token.
\par\textit{Transformer MLP blocks expand hidden dimensions, apply nonlinear gates, and project back per token.}
\paragraph{FC-1189 [hard]} State the main tradeoff and production pitfall for Feed-forward networks.
\noindent\textbf{Answer.} Tradeoff: They hold substantial parameters and compute while offering parallel token processing. Pitfall: Ignoring activation memory underestimates training capacity.
\par\textit{Tradeoff: They hold substantial parameters and compute while offering parallel token processing. Pitfall: Ignoring activation memory underestimates training capacity.}
\paragraph{LA-1190 [hard]} Explain Feed-forward networks from first principles, then show how you would evaluate and operate it in production.
\noindent\textbf{Answer.} Start with the mechanism: Transformer MLP blocks expand hidden dimensions, apply nonlinear gates, and project back per token. Make the tradeoff explicit: They hold substantial parameters and compute while offering parallel token processing. Define workload-specific offline metrics and a latency/cost budget, test important slices, instrument the critical path, and create a rollback criterion. The production review must explicitly guard against this failure: Ignoring activation memory underestimates training capacity.
\par\textit{A strong answer connects mechanism, measurable tradeoffs, failure isolation, observability, and rollback rather than listing product features.}
\paragraph{MCQ-1191 [medium]} Which statement best describes Layer normalization and RMSNorm?
\begin{enumerate}[label=\Alph*.]
\item Multiple heads project attention into separate subspaces before concatenation and output projection.
\item A student learns from teacher probabilities, generated data, rationales, or intermediate representations.
\item Normalization stabilizes activations; RMSNorm removes mean centering and normalizes by root-mean-square magnitude.
\item Attention computes scaled query-key scores, applies masking and softmax, then mixes value vectors.
\end{enumerate}
\noindent\textbf{Answer.} Normalization stabilizes activations; RMSNorm removes mean centering and normalizes by root-mean-square magnitude.
\par\textit{Normalization stabilizes activations; RMSNorm removes mean centering and normalizes by root-mean-square magnitude. Tradeoff: Simpler normalization can be faster but architecture placement changes optimization dynamics. Watch for: Moving pre-norm to post-norm without retuning can destabilize deep training.}
\paragraph{MCQ-1192 [hard]} What is the central engineering tradeoff of Layer normalization and RMSNorm?
\begin{enumerate}[label=\Alph*.]
\item Every token can interact directly, while standard sequence cost grows quadratically.
\item Simpler normalization can be faster but architecture placement changes optimization dynamics.
\item Diverse interactions improve capacity but add projections and KV memory.
\item Serving cost falls but rare capabilities and calibration may be lost.
\end{enumerate}
\noindent\textbf{Answer.} Simpler normalization can be faster but architecture placement changes optimization dynamics.
\par\textit{Normalization stabilizes activations; RMSNorm removes mean centering and normalizes by root-mean-square magnitude. Tradeoff: Simpler normalization can be faster but architecture placement changes optimization dynamics. Watch for: Moving pre-norm to post-norm without retuning can destabilize deep training.}
\paragraph{MCQ-1193 [hard]} Which failure mode should an interviewer expect you to identify for Layer normalization and RMSNorm?
\begin{enumerate}[label=\Alph*.]
\item Assuming heads are independently interpretable can overstate mechanistic conclusions.
\item Omitting the square-root head-dimension scale saturates softmax gradients.
\item Moving pre-norm to post-norm without retuning can destabilize deep training.
\item Distilling only easy teacher outputs creates a deceptively strong narrow model.
\end{enumerate}
\noindent\textbf{Answer.} Moving pre-norm to post-norm without retuning can destabilize deep training.
\par\textit{Normalization stabilizes activations; RMSNorm removes mean centering and normalizes by root-mean-square magnitude. Tradeoff: Simpler normalization can be faster but architecture placement changes optimization dynamics. Watch for: Moving pre-norm to post-norm without retuning can destabilize deep training.}
\paragraph{MCQ-1194 [medium]} A design review is specifically concerned with this risk: Moving pre-norm to post-norm without retuning can destabilize deep training. Which topic should be reviewed first?
\begin{enumerate}[label=\Alph*.]
\item Layer normalization and RMSNorm
\item Self-attention math
\item Multi-head attention
\item Knowledge distillation
\end{enumerate}
\noindent\textbf{Answer.} Layer normalization and RMSNorm
\par\textit{Normalization stabilizes activations; RMSNorm removes mean centering and normalizes by root-mean-square magnitude. Tradeoff: Simpler normalization can be faster but architecture placement changes optimization dynamics. Watch for: Moving pre-norm to post-norm without retuning can destabilize deep training.}
\paragraph{MCQ-1195 [hard]} Which design note most accurately balances the benefits and costs of Layer normalization and RMSNorm?
\begin{enumerate}[label=\Alph*.]
\item Diverse interactions improve capacity but add projections and KV memory.
\item Every token can interact directly, while standard sequence cost grows quadratically.
\item Simpler normalization can be faster but architecture placement changes optimization dynamics.
\item Serving cost falls but rare capabilities and calibration may be lost.
\end{enumerate}
\noindent\textbf{Answer.} Simpler normalization can be faster but architecture placement changes optimization dynamics.
\par\textit{Normalization stabilizes activations; RMSNorm removes mean centering and normalizes by root-mean-square magnitude. Tradeoff: Simpler normalization can be faster but architecture placement changes optimization dynamics. Watch for: Moving pre-norm to post-norm without retuning can destabilize deep training.}
\paragraph{MCQ-1196 [medium]} Which explanation would be strongest in a system-design interview about Layer normalization and RMSNorm?
\begin{enumerate}[label=\Alph*.]
\item Normalization stabilizes activations; RMSNorm removes mean centering and normalizes by root-mean-square magnitude.
\item Multiple heads project attention into separate subspaces before concatenation and output projection.
\item Attention computes scaled query-key scores, applies masking and softmax, then mixes value vectors.
\item A student learns from teacher probabilities, generated data, rationales, or intermediate representations.
\end{enumerate}
\noindent\textbf{Answer.} Normalization stabilizes activations; RMSNorm removes mean centering and normalizes by root-mean-square magnitude.
\par\textit{Normalization stabilizes activations; RMSNorm removes mean centering and normalizes by root-mean-square magnitude. Tradeoff: Simpler normalization can be faster but architecture placement changes optimization dynamics. Watch for: Moving pre-norm to post-norm without retuning can destabilize deep training.}
\paragraph{MCQ-1197 [hard]} Which production warning is most directly associated with Layer normalization and RMSNorm?
\begin{enumerate}[label=\Alph*.]
\item Omitting the square-root head-dimension scale saturates softmax gradients.
\item Distilling only easy teacher outputs creates a deceptively strong narrow model.
\item Assuming heads are independently interpretable can overstate mechanistic conclusions.
\item Moving pre-norm to post-norm without retuning can destabilize deep training.
\end{enumerate}
\noindent\textbf{Answer.} Moving pre-norm to post-norm without retuning can destabilize deep training.
\par\textit{Normalization stabilizes activations; RMSNorm removes mean centering and normalizes by root-mean-square magnitude. Tradeoff: Simpler normalization can be faster but architecture placement changes optimization dynamics. Watch for: Moving pre-norm to post-norm without retuning can destabilize deep training.}
\paragraph{FC-1198 [medium]} Define Layer normalization and RMSNorm in interview-ready language.
\noindent\textbf{Answer.} Normalization stabilizes activations; RMSNorm removes mean centering and normalizes by root-mean-square magnitude.
\par\textit{Normalization stabilizes activations; RMSNorm removes mean centering and normalizes by root-mean-square magnitude.}
\paragraph{FC-1199 [hard]} State the main tradeoff and production pitfall for Layer normalization and RMSNorm.
\noindent\textbf{Answer.} Tradeoff: Simpler normalization can be faster but architecture placement changes optimization dynamics. Pitfall: Moving pre-norm to post-norm without retuning can destabilize deep training.
\par\textit{Tradeoff: Simpler normalization can be faster but architecture placement changes optimization dynamics. Pitfall: Moving pre-norm to post-norm without retuning can destabilize deep training.}
\paragraph{LA-1200 [hard]} Explain Layer normalization and RMSNorm from first principles, then show how you would evaluate and operate it in production.
\noindent\textbf{Answer.} Start with the mechanism: Normalization stabilizes activations; RMSNorm removes mean centering and normalizes by root-mean-square magnitude. Make the tradeoff explicit: Simpler normalization can be faster but architecture placement changes optimization dynamics. Define workload-specific offline metrics and a latency/cost budget, test important slices, instrument the critical path, and create a rollback criterion. The production review must explicitly guard against this failure: Moving pre-norm to post-norm without retuning can destabilize deep training.
\par\textit{A strong answer connects mechanism, measurable tradeoffs, failure isolation, observability, and rollback rather than listing product features.}
\paragraph{MCQ-1201 [medium]} Which statement best describes Residual connections?
\begin{enumerate}[label=\Alph*.]
\item Residual paths add transformed activations back to the stream to improve gradient flow and feature reuse.
\item Supervised fine-tuning teaches instruction-response behavior from curated demonstrations.
\item A preference model supplies rewards while PPO constrains policy updates relative to a reference model.
\item MoE routes each token to a subset of expert networks, increasing parameters without proportional per-token compute.
\end{enumerate}
\noindent\textbf{Answer.} Residual paths add transformed activations back to the stream to improve gradient flow and feature reuse.
\par\textit{Residual paths add transformed activations back to the stream to improve gradient flow and feature reuse. Tradeoff: Deep optimization becomes feasible, while residual scale can grow and require controls. Watch for: In-place mutation of residual tensors can break autograd or checkpointing.}
\paragraph{MCQ-1202 [hard]} What is the central engineering tradeoff of Residual connections?
\begin{enumerate}[label=\Alph*.]
\item It can optimize human preferences but adds reward hacking and training complexity.
\item Deep optimization becomes feasible, while residual scale can grow and require controls.
\item It is direct and stable but can overfit style and inherit annotation errors.
\item Capacity scales efficiently but routing balance, communication, and memory are hard.
\end{enumerate}
\noindent\textbf{Answer.} Deep optimization becomes feasible, while residual scale can grow and require controls.
\par\textit{Residual paths add transformed activations back to the stream to improve gradient flow and feature reuse. Tradeoff: Deep optimization becomes feasible, while residual scale can grow and require controls. Watch for: In-place mutation of residual tensors can break autograd or checkpointing.}
\paragraph{MCQ-1203 [hard]} Which failure mode should an interviewer expect you to identify for Residual connections?
\begin{enumerate}[label=\Alph*.]
\item Training on conflicting instructions without provenance makes behavior unstable.
\item A biased reward model can produce confidently optimized failure modes.
\item Expert collapse and uneven routing create stragglers and weak specialization.
\item In-place mutation of residual tensors can break autograd or checkpointing.
\end{enumerate}
\noindent\textbf{Answer.} In-place mutation of residual tensors can break autograd or checkpointing.
\par\textit{Residual paths add transformed activations back to the stream to improve gradient flow and feature reuse. Tradeoff: Deep optimization becomes feasible, while residual scale can grow and require controls. Watch for: In-place mutation of residual tensors can break autograd or checkpointing.}
\paragraph{MCQ-1204 [medium]} A design review is specifically concerned with this risk: In-place mutation of residual tensors can break autograd or checkpointing. Which topic should be reviewed first?
\begin{enumerate}[label=\Alph*.]
\item SFT
\item RLHF and PPO
\item Residual connections
\item Mixture of Experts
\end{enumerate}
\noindent\textbf{Answer.} Residual connections
\par\textit{Residual paths add transformed activations back to the stream to improve gradient flow and feature reuse. Tradeoff: Deep optimization becomes feasible, while residual scale can grow and require controls. Watch for: In-place mutation of residual tensors can break autograd or checkpointing.}
\paragraph{MCQ-1205 [hard]} Which design note most accurately balances the benefits and costs of Residual connections?
\begin{enumerate}[label=\Alph*.]
\item Capacity scales efficiently but routing balance, communication, and memory are hard.
\item Deep optimization becomes feasible, while residual scale can grow and require controls.
\item It can optimize human preferences but adds reward hacking and training complexity.
\item It is direct and stable but can overfit style and inherit annotation errors.
\end{enumerate}
\noindent\textbf{Answer.} Deep optimization becomes feasible, while residual scale can grow and require controls.
\par\textit{Residual paths add transformed activations back to the stream to improve gradient flow and feature reuse. Tradeoff: Deep optimization becomes feasible, while residual scale can grow and require controls. Watch for: In-place mutation of residual tensors can break autograd or checkpointing.}
\paragraph{MCQ-1206 [medium]} Which explanation would be strongest in a system-design interview about Residual connections?
\begin{enumerate}[label=\Alph*.]
\item Residual paths add transformed activations back to the stream to improve gradient flow and feature reuse.
\item MoE routes each token to a subset of expert networks, increasing parameters without proportional per-token compute.
\item Supervised fine-tuning teaches instruction-response behavior from curated demonstrations.
\item A preference model supplies rewards while PPO constrains policy updates relative to a reference model.
\end{enumerate}
\noindent\textbf{Answer.} Residual paths add transformed activations back to the stream to improve gradient flow and feature reuse.
\par\textit{Residual paths add transformed activations back to the stream to improve gradient flow and feature reuse. Tradeoff: Deep optimization becomes feasible, while residual scale can grow and require controls. Watch for: In-place mutation of residual tensors can break autograd or checkpointing.}
\paragraph{MCQ-1207 [hard]} Which production warning is most directly associated with Residual connections?
\begin{enumerate}[label=\Alph*.]
\item A biased reward model can produce confidently optimized failure modes.
\item In-place mutation of residual tensors can break autograd or checkpointing.
\item Training on conflicting instructions without provenance makes behavior unstable.
\item Expert collapse and uneven routing create stragglers and weak specialization.
\end{enumerate}
\noindent\textbf{Answer.} In-place mutation of residual tensors can break autograd or checkpointing.
\par\textit{Residual paths add transformed activations back to the stream to improve gradient flow and feature reuse. Tradeoff: Deep optimization becomes feasible, while residual scale can grow and require controls. Watch for: In-place mutation of residual tensors can break autograd or checkpointing.}
\paragraph{FC-1208 [medium]} Define Residual connections in interview-ready language.
\noindent\textbf{Answer.} Residual paths add transformed activations back to the stream to improve gradient flow and feature reuse.
\par\textit{Residual paths add transformed activations back to the stream to improve gradient flow and feature reuse.}
\paragraph{FC-1209 [hard]} State the main tradeoff and production pitfall for Residual connections.
\noindent\textbf{Answer.} Tradeoff: Deep optimization becomes feasible, while residual scale can grow and require controls. Pitfall: In-place mutation of residual tensors can break autograd or checkpointing.
\par\textit{Tradeoff: Deep optimization becomes feasible, while residual scale can grow and require controls. Pitfall: In-place mutation of residual tensors can break autograd or checkpointing.}
\paragraph{LA-1210 [hard]} Explain Residual connections from first principles, then show how you would evaluate and operate it in production.
\noindent\textbf{Answer.} Start with the mechanism: Residual paths add transformed activations back to the stream to improve gradient flow and feature reuse. Make the tradeoff explicit: Deep optimization becomes feasible, while residual scale can grow and require controls. Define workload-specific offline metrics and a latency/cost budget, test important slices, instrument the critical path, and create a rollback criterion. The production review must explicitly guard against this failure: In-place mutation of residual tensors can break autograd or checkpointing.
\par\textit{A strong answer connects mechanism, measurable tradeoffs, failure isolation, observability, and rollback rather than listing product features.}
\paragraph{MCQ-1211 [medium]} Which statement best describes Mixture of Experts?
\begin{enumerate}[label=\Alph*.]
\item QLoRA backpropagates through a frozen quantized base model into LoRA adapters using memory-saving quantization techniques.
\item Cross-entropy is negative log likelihood; perplexity exponentiates average token loss.
\item LoRA trains low-rank update matrices while freezing base weights for parameter-efficient adaptation.
\item MoE routes each token to a subset of expert networks, increasing parameters without proportional per-token compute.
\end{enumerate}
\noindent\textbf{Answer.} MoE routes each token to a subset of expert networks, increasing parameters without proportional per-token compute.
\par\textit{MoE routes each token to a subset of expert networks, increasing parameters without proportional per-token compute. Tradeoff: Capacity scales efficiently but routing balance, communication, and memory are hard. Watch for: Expert collapse and uneven routing create stragglers and weak specialization.}
\paragraph{MCQ-1212 [hard]} What is the central engineering tradeoff of Mixture of Experts?
\begin{enumerate}[label=\Alph*.]
\item Fine-tuning large models becomes cheaper, while kernels and quantization settings affect stability.
\item It is stable for model comparison on the same tokenization and data, not across arbitrary setups.
\item Capacity scales efficiently but routing balance, communication, and memory are hard.
\item Memory and storage fall, while rank and target-module choices limit expressiveness.
\end{enumerate}
\noindent\textbf{Answer.} Capacity scales efficiently but routing balance, communication, and memory are hard.
\par\textit{MoE routes each token to a subset of expert networks, increasing parameters without proportional per-token compute. Tradeoff: Capacity scales efficiently but routing balance, communication, and memory are hard. Watch for: Expert collapse and uneven routing create stragglers and weak specialization.}
\paragraph{MCQ-1213 [hard]} Which failure mode should an interviewer expect you to identify for Mixture of Experts?
\begin{enumerate}[label=\Alph*.]
\item Expert collapse and uneven routing create stragglers and weak specialization.
\item Comparing perplexity across vocabularies or tokenizers is misleading.
\item Merging incompatible adapters without evaluation creates interference.
\item Confusing training quantization with final serving quantization produces wrong expectations.
\end{enumerate}
\noindent\textbf{Answer.} Expert collapse and uneven routing create stragglers and weak specialization.
\par\textit{MoE routes each token to a subset of expert networks, increasing parameters without proportional per-token compute. Tradeoff: Capacity scales efficiently but routing balance, communication, and memory are hard. Watch for: Expert collapse and uneven routing create stragglers and weak specialization.}
\paragraph{MCQ-1214 [medium]} A design review is specifically concerned with this risk: Expert collapse and uneven routing create stragglers and weak specialization. Which topic should be reviewed first?
\begin{enumerate}[label=\Alph*.]
\item Cross-entropy and perplexity
\item LoRA
\item QLoRA
\item Mixture of Experts
\end{enumerate}
\noindent\textbf{Answer.} Mixture of Experts
\par\textit{MoE routes each token to a subset of expert networks, increasing parameters without proportional per-token compute. Tradeoff: Capacity scales efficiently but routing balance, communication, and memory are hard. Watch for: Expert collapse and uneven routing create stragglers and weak specialization.}
\paragraph{MCQ-1215 [hard]} Which design note most accurately balances the benefits and costs of Mixture of Experts?
\begin{enumerate}[label=\Alph*.]
\item Capacity scales efficiently but routing balance, communication, and memory are hard.
\item Fine-tuning large models becomes cheaper, while kernels and quantization settings affect stability.
\item It is stable for model comparison on the same tokenization and data, not across arbitrary setups.
\item Memory and storage fall, while rank and target-module choices limit expressiveness.
\end{enumerate}
\noindent\textbf{Answer.} Capacity scales efficiently but routing balance, communication, and memory are hard.
\par\textit{MoE routes each token to a subset of expert networks, increasing parameters without proportional per-token compute. Tradeoff: Capacity scales efficiently but routing balance, communication, and memory are hard. Watch for: Expert collapse and uneven routing create stragglers and weak specialization.}
\paragraph{MCQ-1216 [medium]} Which explanation would be strongest in a system-design interview about Mixture of Experts?
\begin{enumerate}[label=\Alph*.]
\item Cross-entropy is negative log likelihood; perplexity exponentiates average token loss.
\item MoE routes each token to a subset of expert networks, increasing parameters without proportional per-token compute.
\item QLoRA backpropagates through a frozen quantized base model into LoRA adapters using memory-saving quantization techniques.
\item LoRA trains low-rank update matrices while freezing base weights for parameter-efficient adaptation.
\end{enumerate}
\noindent\textbf{Answer.} MoE routes each token to a subset of expert networks, increasing parameters without proportional per-token compute.
\par\textit{MoE routes each token to a subset of expert networks, increasing parameters without proportional per-token compute. Tradeoff: Capacity scales efficiently but routing balance, communication, and memory are hard. Watch for: Expert collapse and uneven routing create stragglers and weak specialization.}
\paragraph{MCQ-1217 [hard]} Which production warning is most directly associated with Mixture of Experts?
\begin{enumerate}[label=\Alph*.]
\item Confusing training quantization with final serving quantization produces wrong expectations.
\item Expert collapse and uneven routing create stragglers and weak specialization.
\item Merging incompatible adapters without evaluation creates interference.
\item Comparing perplexity across vocabularies or tokenizers is misleading.
\end{enumerate}
\noindent\textbf{Answer.} Expert collapse and uneven routing create stragglers and weak specialization.
\par\textit{MoE routes each token to a subset of expert networks, increasing parameters without proportional per-token compute. Tradeoff: Capacity scales efficiently but routing balance, communication, and memory are hard. Watch for: Expert collapse and uneven routing create stragglers and weak specialization.}
\paragraph{FC-1218 [medium]} Define Mixture of Experts in interview-ready language.
\noindent\textbf{Answer.} MoE routes each token to a subset of expert networks, increasing parameters without proportional per-token compute.
\par\textit{MoE routes each token to a subset of expert networks, increasing parameters without proportional per-token compute.}
\paragraph{FC-1219 [hard]} State the main tradeoff and production pitfall for Mixture of Experts.
\noindent\textbf{Answer.} Tradeoff: Capacity scales efficiently but routing balance, communication, and memory are hard. Pitfall: Expert collapse and uneven routing create stragglers and weak specialization.
\par\textit{Tradeoff: Capacity scales efficiently but routing balance, communication, and memory are hard. Pitfall: Expert collapse and uneven routing create stragglers and weak specialization.}
\paragraph{LA-1220 [hard]} Explain Mixture of Experts from first principles, then show how you would evaluate and operate it in production.
\noindent\textbf{Answer.} Start with the mechanism: MoE routes each token to a subset of expert networks, increasing parameters without proportional per-token compute. Make the tradeoff explicit: Capacity scales efficiently but routing balance, communication, and memory are hard. Define workload-specific offline metrics and a latency/cost budget, test important slices, instrument the critical path, and create a rollback criterion. The production review must explicitly guard against this failure: Expert collapse and uneven routing create stragglers and weak specialization.
\par\textit{A strong answer connects mechanism, measurable tradeoffs, failure isolation, observability, and rollback rather than listing product features.}
\paragraph{MCQ-1221 [medium]} Which statement best describes Causal language modeling?
\begin{enumerate}[label=\Alph*.]
\item Next-token training factorizes sequence likelihood from left to right under a causal mask.
\item Supervised fine-tuning teaches instruction-response behavior from curated demonstrations.
\item Rotary position embeddings rotate query and key coordinates so dot products encode relative offsets.
\item Tokenizers map text or bytes into model vocabulary IDs using algorithms such as BPE, WordPiece, or unigram models.
\end{enumerate}
\noindent\textbf{Answer.} Next-token training factorizes sequence likelihood from left to right under a causal mask.
\par\textit{Next-token training factorizes sequence likelihood from left to right under a causal mask. Tradeoff: It provides a simple scalable objective but does not directly optimize instruction usefulness or truth. Watch for: Interpreting low perplexity as safe task performance is invalid.}
\paragraph{MCQ-1222 [hard]} What is the central engineering tradeoff of Causal language modeling?
\begin{enumerate}[label=\Alph*.]
\item RoPE integrates cleanly with attention but long-context scaling needs careful frequency treatment.
\item It provides a simple scalable objective but does not directly optimize instruction usefulness or truth.
\item Larger vocabularies shorten sequences but increase embedding parameters and sparse coverage.
\item It is direct and stable but can overfit style and inherit annotation errors.
\end{enumerate}
\noindent\textbf{Answer.} It provides a simple scalable objective but does not directly optimize instruction usefulness or truth.
\par\textit{Next-token training factorizes sequence likelihood from left to right under a causal mask. Tradeoff: It provides a simple scalable objective but does not directly optimize instruction usefulness or truth. Watch for: Interpreting low perplexity as safe task performance is invalid.}
\paragraph{MCQ-1223 [hard]} Which failure mode should an interviewer expect you to identify for Causal language modeling?
\begin{enumerate}[label=\Alph*.]
\item Training on conflicting instructions without provenance makes behavior unstable.
\item Naively increasing the maximum length does not preserve trained frequency behavior.
\item Comparing context lengths across tokenizers as if tokens were equal misstates capacity and cost.
\item Interpreting low perplexity as safe task performance is invalid.
\end{enumerate}
\noindent\textbf{Answer.} Interpreting low perplexity as safe task performance is invalid.
\par\textit{Next-token training factorizes sequence likelihood from left to right under a causal mask. Tradeoff: It provides a simple scalable objective but does not directly optimize instruction usefulness or truth. Watch for: Interpreting low perplexity as safe task performance is invalid.}
\paragraph{MCQ-1224 [medium]} A design review is specifically concerned with this risk: Interpreting low perplexity as safe task performance is invalid. Which topic should be reviewed first?
\begin{enumerate}[label=\Alph*.]
\item SFT
\item RoPE
\item Causal language modeling
\item Tokenization
\end{enumerate}
\noindent\textbf{Answer.} Causal language modeling
\par\textit{Next-token training factorizes sequence likelihood from left to right under a causal mask. Tradeoff: It provides a simple scalable objective but does not directly optimize instruction usefulness or truth. Watch for: Interpreting low perplexity as safe task performance is invalid.}
\paragraph{MCQ-1225 [hard]} Which design note most accurately balances the benefits and costs of Causal language modeling?
\begin{enumerate}[label=\Alph*.]
\item It provides a simple scalable objective but does not directly optimize instruction usefulness or truth.
\item It is direct and stable but can overfit style and inherit annotation errors.
\item RoPE integrates cleanly with attention but long-context scaling needs careful frequency treatment.
\item Larger vocabularies shorten sequences but increase embedding parameters and sparse coverage.
\end{enumerate}
\noindent\textbf{Answer.} It provides a simple scalable objective but does not directly optimize instruction usefulness or truth.
\par\textit{Next-token training factorizes sequence likelihood from left to right under a causal mask. Tradeoff: It provides a simple scalable objective but does not directly optimize instruction usefulness or truth. Watch for: Interpreting low perplexity as safe task performance is invalid.}
\paragraph{MCQ-1226 [medium]} Which explanation would be strongest in a system-design interview about Causal language modeling?
\begin{enumerate}[label=\Alph*.]
\item Rotary position embeddings rotate query and key coordinates so dot products encode relative offsets.
\item Tokenizers map text or bytes into model vocabulary IDs using algorithms such as BPE, WordPiece, or unigram models.
\item Supervised fine-tuning teaches instruction-response behavior from curated demonstrations.
\item Next-token training factorizes sequence likelihood from left to right under a causal mask.
\end{enumerate}
\noindent\textbf{Answer.} Next-token training factorizes sequence likelihood from left to right under a causal mask.
\par\textit{Next-token training factorizes sequence likelihood from left to right under a causal mask. Tradeoff: It provides a simple scalable objective but does not directly optimize instruction usefulness or truth. Watch for: Interpreting low perplexity as safe task performance is invalid.}
\paragraph{MCQ-1227 [hard]} Which production warning is most directly associated with Causal language modeling?
\begin{enumerate}[label=\Alph*.]
\item Naively increasing the maximum length does not preserve trained frequency behavior.
\item Training on conflicting instructions without provenance makes behavior unstable.
\item Comparing context lengths across tokenizers as if tokens were equal misstates capacity and cost.
\item Interpreting low perplexity as safe task performance is invalid.
\end{enumerate}
\noindent\textbf{Answer.} Interpreting low perplexity as safe task performance is invalid.
\par\textit{Next-token training factorizes sequence likelihood from left to right under a causal mask. Tradeoff: It provides a simple scalable objective but does not directly optimize instruction usefulness or truth. Watch for: Interpreting low perplexity as safe task performance is invalid.}
\paragraph{FC-1228 [medium]} Define Causal language modeling in interview-ready language.
\noindent\textbf{Answer.} Next-token training factorizes sequence likelihood from left to right under a causal mask.
\par\textit{Next-token training factorizes sequence likelihood from left to right under a causal mask.}
\paragraph{FC-1229 [hard]} State the main tradeoff and production pitfall for Causal language modeling.
\noindent\textbf{Answer.} Tradeoff: It provides a simple scalable objective but does not directly optimize instruction usefulness or truth. Pitfall: Interpreting low perplexity as safe task performance is invalid.
\par\textit{Tradeoff: It provides a simple scalable objective but does not directly optimize instruction usefulness or truth. Pitfall: Interpreting low perplexity as safe task performance is invalid.}
\paragraph{LA-1230 [hard]} Explain Causal language modeling from first principles, then show how you would evaluate and operate it in production.
\noindent\textbf{Answer.} Start with the mechanism: Next-token training factorizes sequence likelihood from left to right under a causal mask. Make the tradeoff explicit: It provides a simple scalable objective but does not directly optimize instruction usefulness or truth. Define workload-specific offline metrics and a latency/cost budget, test important slices, instrument the critical path, and create a rollback criterion. The production review must explicitly guard against this failure: Interpreting low perplexity as safe task performance is invalid.
\par\textit{A strong answer connects mechanism, measurable tradeoffs, failure isolation, observability, and rollback rather than listing product features.}
\paragraph{MCQ-1231 [medium]} Which statement best describes Cross-entropy and perplexity?
\begin{enumerate}[label=\Alph*.]
\item QLoRA backpropagates through a frozen quantized base model into LoRA adapters using memory-saving quantization techniques.
\item Attention computes scaled query-key scores, applies masking and softmax, then mixes value vectors.
\item Cross-entropy is negative log likelihood; perplexity exponentiates average token loss.
\item Multiple heads project attention into separate subspaces before concatenation and output projection.
\end{enumerate}
\noindent\textbf{Answer.} Cross-entropy is negative log likelihood; perplexity exponentiates average token loss.
\par\textit{Cross-entropy is negative log likelihood; perplexity exponentiates average token loss. Tradeoff: It is stable for model comparison on the same tokenization and data, not across arbitrary setups. Watch for: Comparing perplexity across vocabularies or tokenizers is misleading.}
\paragraph{MCQ-1232 [hard]} What is the central engineering tradeoff of Cross-entropy and perplexity?
\begin{enumerate}[label=\Alph*.]
\item Diverse interactions improve capacity but add projections and KV memory.
\item Every token can interact directly, while standard sequence cost grows quadratically.
\item Fine-tuning large models becomes cheaper, while kernels and quantization settings affect stability.
\item It is stable for model comparison on the same tokenization and data, not across arbitrary setups.
\end{enumerate}
\noindent\textbf{Answer.} It is stable for model comparison on the same tokenization and data, not across arbitrary setups.
\par\textit{Cross-entropy is negative log likelihood; perplexity exponentiates average token loss. Tradeoff: It is stable for model comparison on the same tokenization and data, not across arbitrary setups. Watch for: Comparing perplexity across vocabularies or tokenizers is misleading.}
\paragraph{MCQ-1233 [hard]} Which failure mode should an interviewer expect you to identify for Cross-entropy and perplexity?
\begin{enumerate}[label=\Alph*.]
\item Assuming heads are independently interpretable can overstate mechanistic conclusions.
\item Confusing training quantization with final serving quantization produces wrong expectations.
\item Comparing perplexity across vocabularies or tokenizers is misleading.
\item Omitting the square-root head-dimension scale saturates softmax gradients.
\end{enumerate}
\noindent\textbf{Answer.} Comparing perplexity across vocabularies or tokenizers is misleading.
\par\textit{Cross-entropy is negative log likelihood; perplexity exponentiates average token loss. Tradeoff: It is stable for model comparison on the same tokenization and data, not across arbitrary setups. Watch for: Comparing perplexity across vocabularies or tokenizers is misleading.}
\paragraph{MCQ-1234 [medium]} A design review is specifically concerned with this risk: Comparing perplexity across vocabularies or tokenizers is misleading. Which topic should be reviewed first?
\begin{enumerate}[label=\Alph*.]
\item Multi-head attention
\item Cross-entropy and perplexity
\item Self-attention math
\item QLoRA
\end{enumerate}
\noindent\textbf{Answer.} Cross-entropy and perplexity
\par\textit{Cross-entropy is negative log likelihood; perplexity exponentiates average token loss. Tradeoff: It is stable for model comparison on the same tokenization and data, not across arbitrary setups. Watch for: Comparing perplexity across vocabularies or tokenizers is misleading.}
\paragraph{MCQ-1235 [hard]} Which design note most accurately balances the benefits and costs of Cross-entropy and perplexity?
\begin{enumerate}[label=\Alph*.]
\item Diverse interactions improve capacity but add projections and KV memory.
\item Fine-tuning large models becomes cheaper, while kernels and quantization settings affect stability.
\item Every token can interact directly, while standard sequence cost grows quadratically.
\item It is stable for model comparison on the same tokenization and data, not across arbitrary setups.
\end{enumerate}
\noindent\textbf{Answer.} It is stable for model comparison on the same tokenization and data, not across arbitrary setups.
\par\textit{Cross-entropy is negative log likelihood; perplexity exponentiates average token loss. Tradeoff: It is stable for model comparison on the same tokenization and data, not across arbitrary setups. Watch for: Comparing perplexity across vocabularies or tokenizers is misleading.}
\paragraph{MCQ-1236 [medium]} Which explanation would be strongest in a system-design interview about Cross-entropy and perplexity?
\begin{enumerate}[label=\Alph*.]
\item Cross-entropy is negative log likelihood; perplexity exponentiates average token loss.
\item Attention computes scaled query-key scores, applies masking and softmax, then mixes value vectors.
\item QLoRA backpropagates through a frozen quantized base model into LoRA adapters using memory-saving quantization techniques.
\item Multiple heads project attention into separate subspaces before concatenation and output projection.
\end{enumerate}
\noindent\textbf{Answer.} Cross-entropy is negative log likelihood; perplexity exponentiates average token loss.
\par\textit{Cross-entropy is negative log likelihood; perplexity exponentiates average token loss. Tradeoff: It is stable for model comparison on the same tokenization and data, not across arbitrary setups. Watch for: Comparing perplexity across vocabularies or tokenizers is misleading.}
\paragraph{MCQ-1237 [hard]} Which production warning is most directly associated with Cross-entropy and perplexity?
\begin{enumerate}[label=\Alph*.]
\item Confusing training quantization with final serving quantization produces wrong expectations.
\item Comparing perplexity across vocabularies or tokenizers is misleading.
\item Omitting the square-root head-dimension scale saturates softmax gradients.
\item Assuming heads are independently interpretable can overstate mechanistic conclusions.
\end{enumerate}
\noindent\textbf{Answer.} Comparing perplexity across vocabularies or tokenizers is misleading.
\par\textit{Cross-entropy is negative log likelihood; perplexity exponentiates average token loss. Tradeoff: It is stable for model comparison on the same tokenization and data, not across arbitrary setups. Watch for: Comparing perplexity across vocabularies or tokenizers is misleading.}
\paragraph{FC-1238 [medium]} Define Cross-entropy and perplexity in interview-ready language.
\noindent\textbf{Answer.} Cross-entropy is negative log likelihood; perplexity exponentiates average token loss.
\par\textit{Cross-entropy is negative log likelihood; perplexity exponentiates average token loss.}
\paragraph{FC-1239 [hard]} State the main tradeoff and production pitfall for Cross-entropy and perplexity.
\noindent\textbf{Answer.} Tradeoff: It is stable for model comparison on the same tokenization and data, not across arbitrary setups. Pitfall: Comparing perplexity across vocabularies or tokenizers is misleading.
\par\textit{Tradeoff: It is stable for model comparison on the same tokenization and data, not across arbitrary setups. Pitfall: Comparing perplexity across vocabularies or tokenizers is misleading.}
\paragraph{LA-1240 [hard]} Explain Cross-entropy and perplexity from first principles, then show how you would evaluate and operate it in production.
\noindent\textbf{Answer.} Start with the mechanism: Cross-entropy is negative log likelihood; perplexity exponentiates average token loss. Make the tradeoff explicit: It is stable for model comparison on the same tokenization and data, not across arbitrary setups. Define workload-specific offline metrics and a latency/cost budget, test important slices, instrument the critical path, and create a rollback criterion. The production review must explicitly guard against this failure: Comparing perplexity across vocabularies or tokenizers is misleading.
\par\textit{A strong answer connects mechanism, measurable tradeoffs, failure isolation, observability, and rollback rather than listing product features.}
\paragraph{MCQ-1241 [medium]} Which statement best describes SFT?
\begin{enumerate}[label=\Alph*.]
\item A student learns from teacher probabilities, generated data, rationales, or intermediate representations.
\item LoRA trains low-rank update matrices while freezing base weights for parameter-efficient adaptation.
\item Supervised fine-tuning teaches instruction-response behavior from curated demonstrations.
\item Direct Preference Optimization fits preferred over rejected responses through a classification-style objective relative to a reference policy.
\end{enumerate}
\noindent\textbf{Answer.} Supervised fine-tuning teaches instruction-response behavior from curated demonstrations.
\par\textit{Supervised fine-tuning teaches instruction-response behavior from curated demonstrations. Tradeoff: It is direct and stable but can overfit style and inherit annotation errors. Watch for: Training on conflicting instructions without provenance makes behavior unstable.}
\paragraph{MCQ-1242 [hard]} What is the central engineering tradeoff of SFT?
\begin{enumerate}[label=\Alph*.]
\item Memory and storage fall, while rank and target-module choices limit expressiveness.
\item Serving cost falls but rare capabilities and calibration may be lost.
\item It is direct and stable but can overfit style and inherit annotation errors.
\item It simplifies preference training but still depends on pair quality and distribution coverage.
\end{enumerate}
\noindent\textbf{Answer.} It is direct and stable but can overfit style and inherit annotation errors.
\par\textit{Supervised fine-tuning teaches instruction-response behavior from curated demonstrations. Tradeoff: It is direct and stable but can overfit style and inherit annotation errors. Watch for: Training on conflicting instructions without provenance makes behavior unstable.}
\paragraph{MCQ-1243 [hard]} Which failure mode should an interviewer expect you to identify for SFT?
\begin{enumerate}[label=\Alph*.]
\item Training on conflicting instructions without provenance makes behavior unstable.
\item Distilling only easy teacher outputs creates a deceptively strong narrow model.
\item Treating implicit user choices as clean preferences introduces confounding.
\item Merging incompatible adapters without evaluation creates interference.
\end{enumerate}
\noindent\textbf{Answer.} Training on conflicting instructions without provenance makes behavior unstable.
\par\textit{Supervised fine-tuning teaches instruction-response behavior from curated demonstrations. Tradeoff: It is direct and stable but can overfit style and inherit annotation errors. Watch for: Training on conflicting instructions without provenance makes behavior unstable.}
\paragraph{MCQ-1244 [medium]} A design review is specifically concerned with this risk: Training on conflicting instructions without provenance makes behavior unstable. Which topic should be reviewed first?
\begin{enumerate}[label=\Alph*.]
\item LoRA
\item Knowledge distillation
\item SFT
\item DPO
\end{enumerate}
\noindent\textbf{Answer.} SFT
\par\textit{Supervised fine-tuning teaches instruction-response behavior from curated demonstrations. Tradeoff: It is direct and stable but can overfit style and inherit annotation errors. Watch for: Training on conflicting instructions without provenance makes behavior unstable.}
\paragraph{MCQ-1245 [hard]} Which design note most accurately balances the benefits and costs of SFT?
\begin{enumerate}[label=\Alph*.]
\item Memory and storage fall, while rank and target-module choices limit expressiveness.
\item Serving cost falls but rare capabilities and calibration may be lost.
\item It simplifies preference training but still depends on pair quality and distribution coverage.
\item It is direct and stable but can overfit style and inherit annotation errors.
\end{enumerate}
\noindent\textbf{Answer.} It is direct and stable but can overfit style and inherit annotation errors.
\par\textit{Supervised fine-tuning teaches instruction-response behavior from curated demonstrations. Tradeoff: It is direct and stable but can overfit style and inherit annotation errors. Watch for: Training on conflicting instructions without provenance makes behavior unstable.}
\paragraph{MCQ-1246 [medium]} Which explanation would be strongest in a system-design interview about SFT?
\begin{enumerate}[label=\Alph*.]
\item Direct Preference Optimization fits preferred over rejected responses through a classification-style objective relative to a reference policy.
\item A student learns from teacher probabilities, generated data, rationales, or intermediate representations.
\item LoRA trains low-rank update matrices while freezing base weights for parameter-efficient adaptation.
\item Supervised fine-tuning teaches instruction-response behavior from curated demonstrations.
\end{enumerate}
\noindent\textbf{Answer.} Supervised fine-tuning teaches instruction-response behavior from curated demonstrations.
\par\textit{Supervised fine-tuning teaches instruction-response behavior from curated demonstrations. Tradeoff: It is direct and stable but can overfit style and inherit annotation errors. Watch for: Training on conflicting instructions without provenance makes behavior unstable.}
\paragraph{MCQ-1247 [hard]} Which production warning is most directly associated with SFT?
\begin{enumerate}[label=\Alph*.]
\item Treating implicit user choices as clean preferences introduces confounding.
\item Training on conflicting instructions without provenance makes behavior unstable.
\item Merging incompatible adapters without evaluation creates interference.
\item Distilling only easy teacher outputs creates a deceptively strong narrow model.
\end{enumerate}
\noindent\textbf{Answer.} Training on conflicting instructions without provenance makes behavior unstable.
\par\textit{Supervised fine-tuning teaches instruction-response behavior from curated demonstrations. Tradeoff: It is direct and stable but can overfit style and inherit annotation errors. Watch for: Training on conflicting instructions without provenance makes behavior unstable.}
\paragraph{FC-1248 [medium]} Define SFT in interview-ready language.
\noindent\textbf{Answer.} Supervised fine-tuning teaches instruction-response behavior from curated demonstrations.
\par\textit{Supervised fine-tuning teaches instruction-response behavior from curated demonstrations.}
\paragraph{FC-1249 [hard]} State the main tradeoff and production pitfall for SFT.
\noindent\textbf{Answer.} Tradeoff: It is direct and stable but can overfit style and inherit annotation errors. Pitfall: Training on conflicting instructions without provenance makes behavior unstable.
\par\textit{Tradeoff: It is direct and stable but can overfit style and inherit annotation errors. Pitfall: Training on conflicting instructions without provenance makes behavior unstable.}
\paragraph{LA-1250 [hard]} Explain SFT from first principles, then show how you would evaluate and operate it in production.
\noindent\textbf{Answer.} Start with the mechanism: Supervised fine-tuning teaches instruction-response behavior from curated demonstrations. Make the tradeoff explicit: It is direct and stable but can overfit style and inherit annotation errors. Define workload-specific offline metrics and a latency/cost budget, test important slices, instrument the critical path, and create a rollback criterion. The production review must explicitly guard against this failure: Training on conflicting instructions without provenance makes behavior unstable.
\par\textit{A strong answer connects mechanism, measurable tradeoffs, failure isolation, observability, and rollback rather than listing product features.}
\paragraph{MCQ-1251 [medium]} Which statement best describes RLHF and PPO?
\begin{enumerate}[label=\Alph*.]
\item MoE routes each token to a subset of expert networks, increasing parameters without proportional per-token compute.
\item Supervised fine-tuning teaches instruction-response behavior from curated demonstrations.
\item Next-token training factorizes sequence likelihood from left to right under a causal mask.
\item A preference model supplies rewards while PPO constrains policy updates relative to a reference model.
\end{enumerate}
\noindent\textbf{Answer.} A preference model supplies rewards while PPO constrains policy updates relative to a reference model.
\par\textit{A preference model supplies rewards while PPO constrains policy updates relative to a reference model. Tradeoff: It can optimize human preferences but adds reward hacking and training complexity. Watch for: A biased reward model can produce confidently optimized failure modes.}
\paragraph{MCQ-1252 [hard]} What is the central engineering tradeoff of RLHF and PPO?
\begin{enumerate}[label=\Alph*.]
\item It is direct and stable but can overfit style and inherit annotation errors.
\item Capacity scales efficiently but routing balance, communication, and memory are hard.
\item It can optimize human preferences but adds reward hacking and training complexity.
\item It provides a simple scalable objective but does not directly optimize instruction usefulness or truth.
\end{enumerate}
\noindent\textbf{Answer.} It can optimize human preferences but adds reward hacking and training complexity.
\par\textit{A preference model supplies rewards while PPO constrains policy updates relative to a reference model. Tradeoff: It can optimize human preferences but adds reward hacking and training complexity. Watch for: A biased reward model can produce confidently optimized failure modes.}
\paragraph{MCQ-1253 [hard]} Which failure mode should an interviewer expect you to identify for RLHF and PPO?
\begin{enumerate}[label=\Alph*.]
\item Training on conflicting instructions without provenance makes behavior unstable.
\item A biased reward model can produce confidently optimized failure modes.
\item Expert collapse and uneven routing create stragglers and weak specialization.
\item Interpreting low perplexity as safe task performance is invalid.
\end{enumerate}
\noindent\textbf{Answer.} A biased reward model can produce confidently optimized failure modes.
\par\textit{A preference model supplies rewards while PPO constrains policy updates relative to a reference model. Tradeoff: It can optimize human preferences but adds reward hacking and training complexity. Watch for: A biased reward model can produce confidently optimized failure modes.}
\paragraph{MCQ-1254 [medium]} A design review is specifically concerned with this risk: A biased reward model can produce confidently optimized failure modes. Which topic should be reviewed first?
\begin{enumerate}[label=\Alph*.]
\item Causal language modeling
\item Mixture of Experts
\item SFT
\item RLHF and PPO
\end{enumerate}
\noindent\textbf{Answer.} RLHF and PPO
\par\textit{A preference model supplies rewards while PPO constrains policy updates relative to a reference model. Tradeoff: It can optimize human preferences but adds reward hacking and training complexity. Watch for: A biased reward model can produce confidently optimized failure modes.}
\paragraph{MCQ-1255 [hard]} Which design note most accurately balances the benefits and costs of RLHF and PPO?
\begin{enumerate}[label=\Alph*.]
\item It can optimize human preferences but adds reward hacking and training complexity.
\item Capacity scales efficiently but routing balance, communication, and memory are hard.
\item It is direct and stable but can overfit style and inherit annotation errors.
\item It provides a simple scalable objective but does not directly optimize instruction usefulness or truth.
\end{enumerate}
\noindent\textbf{Answer.} It can optimize human preferences but adds reward hacking and training complexity.
\par\textit{A preference model supplies rewards while PPO constrains policy updates relative to a reference model. Tradeoff: It can optimize human preferences but adds reward hacking and training complexity. Watch for: A biased reward model can produce confidently optimized failure modes.}
\paragraph{MCQ-1256 [medium]} Which explanation would be strongest in a system-design interview about RLHF and PPO?
\begin{enumerate}[label=\Alph*.]
\item Next-token training factorizes sequence likelihood from left to right under a causal mask.
\item A preference model supplies rewards while PPO constrains policy updates relative to a reference model.
\item Supervised fine-tuning teaches instruction-response behavior from curated demonstrations.
\item MoE routes each token to a subset of expert networks, increasing parameters without proportional per-token compute.
\end{enumerate}
\noindent\textbf{Answer.} A preference model supplies rewards while PPO constrains policy updates relative to a reference model.
\par\textit{A preference model supplies rewards while PPO constrains policy updates relative to a reference model. Tradeoff: It can optimize human preferences but adds reward hacking and training complexity. Watch for: A biased reward model can produce confidently optimized failure modes.}
\paragraph{MCQ-1257 [hard]} Which production warning is most directly associated with RLHF and PPO?
\begin{enumerate}[label=\Alph*.]
\item Interpreting low perplexity as safe task performance is invalid.
\item Expert collapse and uneven routing create stragglers and weak specialization.
\item A biased reward model can produce confidently optimized failure modes.
\item Training on conflicting instructions without provenance makes behavior unstable.
\end{enumerate}
\noindent\textbf{Answer.} A biased reward model can produce confidently optimized failure modes.
\par\textit{A preference model supplies rewards while PPO constrains policy updates relative to a reference model. Tradeoff: It can optimize human preferences but adds reward hacking and training complexity. Watch for: A biased reward model can produce confidently optimized failure modes.}
\paragraph{FC-1258 [medium]} Define RLHF and PPO in interview-ready language.
\noindent\textbf{Answer.} A preference model supplies rewards while PPO constrains policy updates relative to a reference model.
\par\textit{A preference model supplies rewards while PPO constrains policy updates relative to a reference model.}
\paragraph{FC-1259 [hard]} State the main tradeoff and production pitfall for RLHF and PPO.
\noindent\textbf{Answer.} Tradeoff: It can optimize human preferences but adds reward hacking and training complexity. Pitfall: A biased reward model can produce confidently optimized failure modes.
\par\textit{Tradeoff: It can optimize human preferences but adds reward hacking and training complexity. Pitfall: A biased reward model can produce confidently optimized failure modes.}
\paragraph{LA-1260 [hard]} Explain RLHF and PPO from first principles, then show how you would evaluate and operate it in production.
\noindent\textbf{Answer.} Start with the mechanism: A preference model supplies rewards while PPO constrains policy updates relative to a reference model. Make the tradeoff explicit: It can optimize human preferences but adds reward hacking and training complexity. Define workload-specific offline metrics and a latency/cost budget, test important slices, instrument the critical path, and create a rollback criterion. The production review must explicitly guard against this failure: A biased reward model can produce confidently optimized failure modes.
\par\textit{A strong answer connects mechanism, measurable tradeoffs, failure isolation, observability, and rollback rather than listing product features.}
\paragraph{MCQ-1261 [medium]} Which statement best describes DPO?
\begin{enumerate}[label=\Alph*.]
\item Position signals break permutation invariance using absolute, relative, rotary, or bias-based schemes.
\item Tokenizers map text or bytes into model vocabulary IDs using algorithms such as BPE, WordPiece, or unigram models.
\item Direct Preference Optimization fits preferred over rejected responses through a classification-style objective relative to a reference policy.
\item Attention computes scaled query-key scores, applies masking and softmax, then mixes value vectors.
\end{enumerate}
\noindent\textbf{Answer.} Direct Preference Optimization fits preferred over rejected responses through a classification-style objective relative to a reference policy.
\par\textit{Direct Preference Optimization fits preferred over rejected responses through a classification-style objective relative to a reference policy. Tradeoff: It simplifies preference training but still depends on pair quality and distribution coverage. Watch for: Treating implicit user choices as clean preferences introduces confounding.}
\paragraph{MCQ-1262 [hard]} What is the central engineering tradeoff of DPO?
\begin{enumerate}[label=\Alph*.]
\item It simplifies preference training but still depends on pair quality and distribution coverage.
\item Every token can interact directly, while standard sequence cost grows quadratically.
\item Larger vocabularies shorten sequences but increase embedding parameters and sparse coverage.
\item Relative methods extrapolate differently, while implementation and scaling choices affect long context.
\end{enumerate}
\noindent\textbf{Answer.} It simplifies preference training but still depends on pair quality and distribution coverage.
\par\textit{Direct Preference Optimization fits preferred over rejected responses through a classification-style objective relative to a reference policy. Tradeoff: It simplifies preference training but still depends on pair quality and distribution coverage. Watch for: Treating implicit user choices as clean preferences introduces confounding.}
\paragraph{MCQ-1263 [hard]} Which failure mode should an interviewer expect you to identify for DPO?
\begin{enumerate}[label=\Alph*.]
\item Treating implicit user choices as clean preferences introduces confounding.
\item Comparing context lengths across tokenizers as if tokens were equal misstates capacity and cost.
\item Omitting the square-root head-dimension scale saturates softmax gradients.
\item Extending context without validating the position scheme can collapse quality.
\end{enumerate}
\noindent\textbf{Answer.} Treating implicit user choices as clean preferences introduces confounding.
\par\textit{Direct Preference Optimization fits preferred over rejected responses through a classification-style objective relative to a reference policy. Tradeoff: It simplifies preference training but still depends on pair quality and distribution coverage. Watch for: Treating implicit user choices as clean preferences introduces confounding.}
\paragraph{MCQ-1264 [medium]} A design review is specifically concerned with this risk: Treating implicit user choices as clean preferences introduces confounding. Which topic should be reviewed first?
\begin{enumerate}[label=\Alph*.]
\item Self-attention math
\item Positional encoding
\item Tokenization
\item DPO
\end{enumerate}
\noindent\textbf{Answer.} DPO
\par\textit{Direct Preference Optimization fits preferred over rejected responses through a classification-style objective relative to a reference policy. Tradeoff: It simplifies preference training but still depends on pair quality and distribution coverage. Watch for: Treating implicit user choices as clean preferences introduces confounding.}
\paragraph{MCQ-1265 [hard]} Which design note most accurately balances the benefits and costs of DPO?
\begin{enumerate}[label=\Alph*.]
\item Every token can interact directly, while standard sequence cost grows quadratically.
\item Relative methods extrapolate differently, while implementation and scaling choices affect long context.
\item It simplifies preference training but still depends on pair quality and distribution coverage.
\item Larger vocabularies shorten sequences but increase embedding parameters and sparse coverage.
\end{enumerate}
\noindent\textbf{Answer.} It simplifies preference training but still depends on pair quality and distribution coverage.
\par\textit{Direct Preference Optimization fits preferred over rejected responses through a classification-style objective relative to a reference policy. Tradeoff: It simplifies preference training but still depends on pair quality and distribution coverage. Watch for: Treating implicit user choices as clean preferences introduces confounding.}
\paragraph{MCQ-1266 [medium]} Which explanation would be strongest in a system-design interview about DPO?
\begin{enumerate}[label=\Alph*.]
\item Position signals break permutation invariance using absolute, relative, rotary, or bias-based schemes.
\item Attention computes scaled query-key scores, applies masking and softmax, then mixes value vectors.
\item Tokenizers map text or bytes into model vocabulary IDs using algorithms such as BPE, WordPiece, or unigram models.
\item Direct Preference Optimization fits preferred over rejected responses through a classification-style objective relative to a reference policy.
\end{enumerate}
\noindent\textbf{Answer.} Direct Preference Optimization fits preferred over rejected responses through a classification-style objective relative to a reference policy.
\par\textit{Direct Preference Optimization fits preferred over rejected responses through a classification-style objective relative to a reference policy. Tradeoff: It simplifies preference training but still depends on pair quality and distribution coverage. Watch for: Treating implicit user choices as clean preferences introduces confounding.}
\paragraph{MCQ-1267 [hard]} Which production warning is most directly associated with DPO?
\begin{enumerate}[label=\Alph*.]
\item Treating implicit user choices as clean preferences introduces confounding.
\item Omitting the square-root head-dimension scale saturates softmax gradients.
\item Comparing context lengths across tokenizers as if tokens were equal misstates capacity and cost.
\item Extending context without validating the position scheme can collapse quality.
\end{enumerate}
\noindent\textbf{Answer.} Treating implicit user choices as clean preferences introduces confounding.
\par\textit{Direct Preference Optimization fits preferred over rejected responses through a classification-style objective relative to a reference policy. Tradeoff: It simplifies preference training but still depends on pair quality and distribution coverage. Watch for: Treating implicit user choices as clean preferences introduces confounding.}
\paragraph{FC-1268 [medium]} Define DPO in interview-ready language.
\noindent\textbf{Answer.} Direct Preference Optimization fits preferred over rejected responses through a classification-style objective relative to a reference policy.
\par\textit{Direct Preference Optimization fits preferred over rejected responses through a classification-style objective relative to a reference policy.}
\paragraph{FC-1269 [hard]} State the main tradeoff and production pitfall for DPO.
\noindent\textbf{Answer.} Tradeoff: It simplifies preference training but still depends on pair quality and distribution coverage. Pitfall: Treating implicit user choices as clean preferences introduces confounding.
\par\textit{Tradeoff: It simplifies preference training but still depends on pair quality and distribution coverage. Pitfall: Treating implicit user choices as clean preferences introduces confounding.}
\paragraph{LA-1270 [hard]} Explain DPO from first principles, then show how you would evaluate and operate it in production.
\noindent\textbf{Answer.} Start with the mechanism: Direct Preference Optimization fits preferred over rejected responses through a classification-style objective relative to a reference policy. Make the tradeoff explicit: It simplifies preference training but still depends on pair quality and distribution coverage. Define workload-specific offline metrics and a latency/cost budget, test important slices, instrument the critical path, and create a rollback criterion. The production review must explicitly guard against this failure: Treating implicit user choices as clean preferences introduces confounding.
\par\textit{A strong answer connects mechanism, measurable tradeoffs, failure isolation, observability, and rollback rather than listing product features.}
\paragraph{MCQ-1271 [medium]} Which statement best describes LoRA?
\begin{enumerate}[label=\Alph*.]
\item Supervised fine-tuning teaches instruction-response behavior from curated demonstrations.
\item LoRA trains low-rank update matrices while freezing base weights for parameter-efficient adaptation.
\item Attention computes scaled query-key scores, applies masking and softmax, then mixes value vectors.
\item Position signals break permutation invariance using absolute, relative, rotary, or bias-based schemes.
\end{enumerate}
\noindent\textbf{Answer.} LoRA trains low-rank update matrices while freezing base weights for parameter-efficient adaptation.
\par\textit{LoRA trains low-rank update matrices while freezing base weights for parameter-efficient adaptation. Tradeoff: Memory and storage fall, while rank and target-module choices limit expressiveness. Watch for: Merging incompatible adapters without evaluation creates interference.}
\paragraph{MCQ-1272 [hard]} What is the central engineering tradeoff of LoRA?
\begin{enumerate}[label=\Alph*.]
\item It is direct and stable but can overfit style and inherit annotation errors.
\item Every token can interact directly, while standard sequence cost grows quadratically.
\item Relative methods extrapolate differently, while implementation and scaling choices affect long context.
\item Memory and storage fall, while rank and target-module choices limit expressiveness.
\end{enumerate}
\noindent\textbf{Answer.} Memory and storage fall, while rank and target-module choices limit expressiveness.
\par\textit{LoRA trains low-rank update matrices while freezing base weights for parameter-efficient adaptation. Tradeoff: Memory and storage fall, while rank and target-module choices limit expressiveness. Watch for: Merging incompatible adapters without evaluation creates interference.}
\paragraph{MCQ-1273 [hard]} Which failure mode should an interviewer expect you to identify for LoRA?
\begin{enumerate}[label=\Alph*.]
\item Omitting the square-root head-dimension scale saturates softmax gradients.
\item Training on conflicting instructions without provenance makes behavior unstable.
\item Merging incompatible adapters without evaluation creates interference.
\item Extending context without validating the position scheme can collapse quality.
\end{enumerate}
\noindent\textbf{Answer.} Merging incompatible adapters without evaluation creates interference.
\par\textit{LoRA trains low-rank update matrices while freezing base weights for parameter-efficient adaptation. Tradeoff: Memory and storage fall, while rank and target-module choices limit expressiveness. Watch for: Merging incompatible adapters without evaluation creates interference.}
\paragraph{MCQ-1274 [medium]} A design review is specifically concerned with this risk: Merging incompatible adapters without evaluation creates interference. Which topic should be reviewed first?
\begin{enumerate}[label=\Alph*.]
\item Self-attention math
\item Positional encoding
\item SFT
\item LoRA
\end{enumerate}
\noindent\textbf{Answer.} LoRA
\par\textit{LoRA trains low-rank update matrices while freezing base weights for parameter-efficient adaptation. Tradeoff: Memory and storage fall, while rank and target-module choices limit expressiveness. Watch for: Merging incompatible adapters without evaluation creates interference.}
\paragraph{MCQ-1275 [hard]} Which design note most accurately balances the benefits and costs of LoRA?
\begin{enumerate}[label=\Alph*.]
\item Relative methods extrapolate differently, while implementation and scaling choices affect long context.
\item Every token can interact directly, while standard sequence cost grows quadratically.
\item It is direct and stable but can overfit style and inherit annotation errors.
\item Memory and storage fall, while rank and target-module choices limit expressiveness.
\end{enumerate}
\noindent\textbf{Answer.} Memory and storage fall, while rank and target-module choices limit expressiveness.
\par\textit{LoRA trains low-rank update matrices while freezing base weights for parameter-efficient adaptation. Tradeoff: Memory and storage fall, while rank and target-module choices limit expressiveness. Watch for: Merging incompatible adapters without evaluation creates interference.}
\paragraph{MCQ-1276 [medium]} Which explanation would be strongest in a system-design interview about LoRA?
\begin{enumerate}[label=\Alph*.]
\item Attention computes scaled query-key scores, applies masking and softmax, then mixes value vectors.
\item Position signals break permutation invariance using absolute, relative, rotary, or bias-based schemes.
\item Supervised fine-tuning teaches instruction-response behavior from curated demonstrations.
\item LoRA trains low-rank update matrices while freezing base weights for parameter-efficient adaptation.
\end{enumerate}
\noindent\textbf{Answer.} LoRA trains low-rank update matrices while freezing base weights for parameter-efficient adaptation.
\par\textit{LoRA trains low-rank update matrices while freezing base weights for parameter-efficient adaptation. Tradeoff: Memory and storage fall, while rank and target-module choices limit expressiveness. Watch for: Merging incompatible adapters without evaluation creates interference.}
\paragraph{MCQ-1277 [hard]} Which production warning is most directly associated with LoRA?
\begin{enumerate}[label=\Alph*.]
\item Merging incompatible adapters without evaluation creates interference.
\item Extending context without validating the position scheme can collapse quality.
\item Training on conflicting instructions without provenance makes behavior unstable.
\item Omitting the square-root head-dimension scale saturates softmax gradients.
\end{enumerate}
\noindent\textbf{Answer.} Merging incompatible adapters without evaluation creates interference.
\par\textit{LoRA trains low-rank update matrices while freezing base weights for parameter-efficient adaptation. Tradeoff: Memory and storage fall, while rank and target-module choices limit expressiveness. Watch for: Merging incompatible adapters without evaluation creates interference.}
\paragraph{FC-1278 [medium]} Define LoRA in interview-ready language.
\noindent\textbf{Answer.} LoRA trains low-rank update matrices while freezing base weights for parameter-efficient adaptation.
\par\textit{LoRA trains low-rank update matrices while freezing base weights for parameter-efficient adaptation.}
\paragraph{FC-1279 [hard]} State the main tradeoff and production pitfall for LoRA.
\noindent\textbf{Answer.} Tradeoff: Memory and storage fall, while rank and target-module choices limit expressiveness. Pitfall: Merging incompatible adapters without evaluation creates interference.
\par\textit{Tradeoff: Memory and storage fall, while rank and target-module choices limit expressiveness. Pitfall: Merging incompatible adapters without evaluation creates interference.}
\paragraph{LA-1280 [hard]} Explain LoRA from first principles, then show how you would evaluate and operate it in production.
\noindent\textbf{Answer.} Start with the mechanism: LoRA trains low-rank update matrices while freezing base weights for parameter-efficient adaptation. Make the tradeoff explicit: Memory and storage fall, while rank and target-module choices limit expressiveness. Define workload-specific offline metrics and a latency/cost budget, test important slices, instrument the critical path, and create a rollback criterion. The production review must explicitly guard against this failure: Merging incompatible adapters without evaluation creates interference.
\par\textit{A strong answer connects mechanism, measurable tradeoffs, failure isolation, observability, and rollback rather than listing product features.}
\paragraph{MCQ-1281 [medium]} Which statement best describes QLoRA?
\begin{enumerate}[label=\Alph*.]
\item Rotary position embeddings rotate query and key coordinates so dot products encode relative offsets.
\item QLoRA backpropagates through a frozen quantized base model into LoRA adapters using memory-saving quantization techniques.
\item Attention computes scaled query-key scores, applies masking and softmax, then mixes value vectors.
\item Cross-entropy is negative log likelihood; perplexity exponentiates average token loss.
\end{enumerate}
\noindent\textbf{Answer.} QLoRA backpropagates through a frozen quantized base model into LoRA adapters using memory-saving quantization techniques.
\par\textit{QLoRA backpropagates through a frozen quantized base model into LoRA adapters using memory-saving quantization techniques. Tradeoff: Fine-tuning large models becomes cheaper, while kernels and quantization settings affect stability. Watch for: Confusing training quantization with final serving quantization produces wrong expectations.}
\paragraph{MCQ-1282 [hard]} What is the central engineering tradeoff of QLoRA?
\begin{enumerate}[label=\Alph*.]
\item Every token can interact directly, while standard sequence cost grows quadratically.
\item It is stable for model comparison on the same tokenization and data, not across arbitrary setups.
\item RoPE integrates cleanly with attention but long-context scaling needs careful frequency treatment.
\item Fine-tuning large models becomes cheaper, while kernels and quantization settings affect stability.
\end{enumerate}
\noindent\textbf{Answer.} Fine-tuning large models becomes cheaper, while kernels and quantization settings affect stability.
\par\textit{QLoRA backpropagates through a frozen quantized base model into LoRA adapters using memory-saving quantization techniques. Tradeoff: Fine-tuning large models becomes cheaper, while kernels and quantization settings affect stability. Watch for: Confusing training quantization with final serving quantization produces wrong expectations.}
\paragraph{MCQ-1283 [hard]} Which failure mode should an interviewer expect you to identify for QLoRA?
\begin{enumerate}[label=\Alph*.]
\item Naively increasing the maximum length does not preserve trained frequency behavior.
\item Confusing training quantization with final serving quantization produces wrong expectations.
\item Omitting the square-root head-dimension scale saturates softmax gradients.
\item Comparing perplexity across vocabularies or tokenizers is misleading.
\end{enumerate}
\noindent\textbf{Answer.} Confusing training quantization with final serving quantization produces wrong expectations.
\par\textit{QLoRA backpropagates through a frozen quantized base model into LoRA adapters using memory-saving quantization techniques. Tradeoff: Fine-tuning large models becomes cheaper, while kernels and quantization settings affect stability. Watch for: Confusing training quantization with final serving quantization produces wrong expectations.}
\paragraph{MCQ-1284 [medium]} A design review is specifically concerned with this risk: Confusing training quantization with final serving quantization produces wrong expectations. Which topic should be reviewed first?
\begin{enumerate}[label=\Alph*.]
\item RoPE
\item QLoRA
\item Self-attention math
\item Cross-entropy and perplexity
\end{enumerate}
\noindent\textbf{Answer.} QLoRA
\par\textit{QLoRA backpropagates through a frozen quantized base model into LoRA adapters using memory-saving quantization techniques. Tradeoff: Fine-tuning large models becomes cheaper, while kernels and quantization settings affect stability. Watch for: Confusing training quantization with final serving quantization produces wrong expectations.}
\paragraph{MCQ-1285 [hard]} Which design note most accurately balances the benefits and costs of QLoRA?
\begin{enumerate}[label=\Alph*.]
\item Fine-tuning large models becomes cheaper, while kernels and quantization settings affect stability.
\item It is stable for model comparison on the same tokenization and data, not across arbitrary setups.
\item Every token can interact directly, while standard sequence cost grows quadratically.
\item RoPE integrates cleanly with attention but long-context scaling needs careful frequency treatment.
\end{enumerate}
\noindent\textbf{Answer.} Fine-tuning large models becomes cheaper, while kernels and quantization settings affect stability.
\par\textit{QLoRA backpropagates through a frozen quantized base model into LoRA adapters using memory-saving quantization techniques. Tradeoff: Fine-tuning large models becomes cheaper, while kernels and quantization settings affect stability. Watch for: Confusing training quantization with final serving quantization produces wrong expectations.}
\paragraph{MCQ-1286 [medium]} Which explanation would be strongest in a system-design interview about QLoRA?
\begin{enumerate}[label=\Alph*.]
\item Cross-entropy is negative log likelihood; perplexity exponentiates average token loss.
\item Rotary position embeddings rotate query and key coordinates so dot products encode relative offsets.
\item Attention computes scaled query-key scores, applies masking and softmax, then mixes value vectors.
\item QLoRA backpropagates through a frozen quantized base model into LoRA adapters using memory-saving quantization techniques.
\end{enumerate}
\noindent\textbf{Answer.} QLoRA backpropagates through a frozen quantized base model into LoRA adapters using memory-saving quantization techniques.
\par\textit{QLoRA backpropagates through a frozen quantized base model into LoRA adapters using memory-saving quantization techniques. Tradeoff: Fine-tuning large models becomes cheaper, while kernels and quantization settings affect stability. Watch for: Confusing training quantization with final serving quantization produces wrong expectations.}
\paragraph{MCQ-1287 [hard]} Which production warning is most directly associated with QLoRA?
\begin{enumerate}[label=\Alph*.]
\item Naively increasing the maximum length does not preserve trained frequency behavior.
\item Omitting the square-root head-dimension scale saturates softmax gradients.
\item Confusing training quantization with final serving quantization produces wrong expectations.
\item Comparing perplexity across vocabularies or tokenizers is misleading.
\end{enumerate}
\noindent\textbf{Answer.} Confusing training quantization with final serving quantization produces wrong expectations.
\par\textit{QLoRA backpropagates through a frozen quantized base model into LoRA adapters using memory-saving quantization techniques. Tradeoff: Fine-tuning large models becomes cheaper, while kernels and quantization settings affect stability. Watch for: Confusing training quantization with final serving quantization produces wrong expectations.}
\paragraph{FC-1288 [medium]} Define QLoRA in interview-ready language.
\noindent\textbf{Answer.} QLoRA backpropagates through a frozen quantized base model into LoRA adapters using memory-saving quantization techniques.
\par\textit{QLoRA backpropagates through a frozen quantized base model into LoRA adapters using memory-saving quantization techniques.}
\paragraph{FC-1289 [hard]} State the main tradeoff and production pitfall for QLoRA.
\noindent\textbf{Answer.} Tradeoff: Fine-tuning large models becomes cheaper, while kernels and quantization settings affect stability. Pitfall: Confusing training quantization with final serving quantization produces wrong expectations.
\par\textit{Tradeoff: Fine-tuning large models becomes cheaper, while kernels and quantization settings affect stability. Pitfall: Confusing training quantization with final serving quantization produces wrong expectations.}
\paragraph{LA-1290 [hard]} Explain QLoRA from first principles, then show how you would evaluate and operate it in production.
\noindent\textbf{Answer.} Start with the mechanism: QLoRA backpropagates through a frozen quantized base model into LoRA adapters using memory-saving quantization techniques. Make the tradeoff explicit: Fine-tuning large models becomes cheaper, while kernels and quantization settings affect stability. Define workload-specific offline metrics and a latency/cost budget, test important slices, instrument the critical path, and create a rollback criterion. The production review must explicitly guard against this failure: Confusing training quantization with final serving quantization produces wrong expectations.
\par\textit{A strong answer connects mechanism, measurable tradeoffs, failure isolation, observability, and rollback rather than listing product features.}
\paragraph{MCQ-1291 [medium]} Which statement best describes Knowledge distillation?
\begin{enumerate}[label=\Alph*.]
\item Multiple heads project attention into separate subspaces before concatenation and output projection.
\item QLoRA backpropagates through a frozen quantized base model into LoRA adapters using memory-saving quantization techniques.
\item LoRA trains low-rank update matrices while freezing base weights for parameter-efficient adaptation.
\item A student learns from teacher probabilities, generated data, rationales, or intermediate representations.
\end{enumerate}
\noindent\textbf{Answer.} A student learns from teacher probabilities, generated data, rationales, or intermediate representations.
\par\textit{A student learns from teacher probabilities, generated data, rationales, or intermediate representations. Tradeoff: Serving cost falls but rare capabilities and calibration may be lost. Watch for: Distilling only easy teacher outputs creates a deceptively strong narrow model.}
\paragraph{MCQ-1292 [hard]} What is the central engineering tradeoff of Knowledge distillation?
\begin{enumerate}[label=\Alph*.]
\item Fine-tuning large models becomes cheaper, while kernels and quantization settings affect stability.
\item Memory and storage fall, while rank and target-module choices limit expressiveness.
\item Serving cost falls but rare capabilities and calibration may be lost.
\item Diverse interactions improve capacity but add projections and KV memory.
\end{enumerate}
\noindent\textbf{Answer.} Serving cost falls but rare capabilities and calibration may be lost.
\par\textit{A student learns from teacher probabilities, generated data, rationales, or intermediate representations. Tradeoff: Serving cost falls but rare capabilities and calibration may be lost. Watch for: Distilling only easy teacher outputs creates a deceptively strong narrow model.}
\paragraph{MCQ-1293 [hard]} Which failure mode should an interviewer expect you to identify for Knowledge distillation?
\begin{enumerate}[label=\Alph*.]
\item Merging incompatible adapters without evaluation creates interference.
\item Distilling only easy teacher outputs creates a deceptively strong narrow model.
\item Assuming heads are independently interpretable can overstate mechanistic conclusions.
\item Confusing training quantization with final serving quantization produces wrong expectations.
\end{enumerate}
\noindent\textbf{Answer.} Distilling only easy teacher outputs creates a deceptively strong narrow model.
\par\textit{A student learns from teacher probabilities, generated data, rationales, or intermediate representations. Tradeoff: Serving cost falls but rare capabilities and calibration may be lost. Watch for: Distilling only easy teacher outputs creates a deceptively strong narrow model.}
\paragraph{MCQ-1294 [medium]} A design review is specifically concerned with this risk: Distilling only easy teacher outputs creates a deceptively strong narrow model. Which topic should be reviewed first?
\begin{enumerate}[label=\Alph*.]
\item QLoRA
\item LoRA
\item Knowledge distillation
\item Multi-head attention
\end{enumerate}
\noindent\textbf{Answer.} Knowledge distillation
\par\textit{A student learns from teacher probabilities, generated data, rationales, or intermediate representations. Tradeoff: Serving cost falls but rare capabilities and calibration may be lost. Watch for: Distilling only easy teacher outputs creates a deceptively strong narrow model.}
\paragraph{MCQ-1295 [hard]} Which design note most accurately balances the benefits and costs of Knowledge distillation?
\begin{enumerate}[label=\Alph*.]
\item Fine-tuning large models becomes cheaper, while kernels and quantization settings affect stability.
\item Serving cost falls but rare capabilities and calibration may be lost.
\item Diverse interactions improve capacity but add projections and KV memory.
\item Memory and storage fall, while rank and target-module choices limit expressiveness.
\end{enumerate}
\noindent\textbf{Answer.} Serving cost falls but rare capabilities and calibration may be lost.
\par\textit{A student learns from teacher probabilities, generated data, rationales, or intermediate representations. Tradeoff: Serving cost falls but rare capabilities and calibration may be lost. Watch for: Distilling only easy teacher outputs creates a deceptively strong narrow model.}
\paragraph{MCQ-1296 [medium]} Which explanation would be strongest in a system-design interview about Knowledge distillation?
\begin{enumerate}[label=\Alph*.]
\item A student learns from teacher probabilities, generated data, rationales, or intermediate representations.
\item LoRA trains low-rank update matrices while freezing base weights for parameter-efficient adaptation.
\item QLoRA backpropagates through a frozen quantized base model into LoRA adapters using memory-saving quantization techniques.
\item Multiple heads project attention into separate subspaces before concatenation and output projection.
\end{enumerate}
\noindent\textbf{Answer.} A student learns from teacher probabilities, generated data, rationales, or intermediate representations.
\par\textit{A student learns from teacher probabilities, generated data, rationales, or intermediate representations. Tradeoff: Serving cost falls but rare capabilities and calibration may be lost. Watch for: Distilling only easy teacher outputs creates a deceptively strong narrow model.}
\paragraph{MCQ-1297 [hard]} Which production warning is most directly associated with Knowledge distillation?
\begin{enumerate}[label=\Alph*.]
\item Assuming heads are independently interpretable can overstate mechanistic conclusions.
\item Merging incompatible adapters without evaluation creates interference.
\item Distilling only easy teacher outputs creates a deceptively strong narrow model.
\item Confusing training quantization with final serving quantization produces wrong expectations.
\end{enumerate}
\noindent\textbf{Answer.} Distilling only easy teacher outputs creates a deceptively strong narrow model.
\par\textit{A student learns from teacher probabilities, generated data, rationales, or intermediate representations. Tradeoff: Serving cost falls but rare capabilities and calibration may be lost. Watch for: Distilling only easy teacher outputs creates a deceptively strong narrow model.}
\paragraph{FC-1298 [medium]} Define Knowledge distillation in interview-ready language.
\noindent\textbf{Answer.} A student learns from teacher probabilities, generated data, rationales, or intermediate representations.
\par\textit{A student learns from teacher probabilities, generated data, rationales, or intermediate representations.}
\paragraph{FC-1299 [hard]} State the main tradeoff and production pitfall for Knowledge distillation.
\noindent\textbf{Answer.} Tradeoff: Serving cost falls but rare capabilities and calibration may be lost. Pitfall: Distilling only easy teacher outputs creates a deceptively strong narrow model.
\par\textit{Tradeoff: Serving cost falls but rare capabilities and calibration may be lost. Pitfall: Distilling only easy teacher outputs creates a deceptively strong narrow model.}
\paragraph{LA-1300 [hard]} Explain Knowledge distillation from first principles, then show how you would evaluate and operate it in production.
\noindent\textbf{Answer.} Start with the mechanism: A student learns from teacher probabilities, generated data, rationales, or intermediate representations. Make the tradeoff explicit: Serving cost falls but rare capabilities and calibration may be lost. Define workload-specific offline metrics and a latency/cost budget, test important slices, instrument the critical path, and create a rollback criterion. The production review must explicitly guard against this failure: Distilling only easy teacher outputs creates a deceptively strong narrow model.
\par\textit{A strong answer connects mechanism, measurable tradeoffs, failure isolation, observability, and rollback rather than listing product features.}
\subsection{Inference Optimization}
\paragraph{MCQ-1301 [medium]} Which statement best describes KV caching?
\begin{enumerate}[label=\Alph*.]
\item Requests join and leave a running decode batch at iteration boundaries rather than waiting for a fixed batch.
\item Pipeline parallelism assigns layer stages to devices and streams microbatches through them.
\item MoE experts are distributed across devices and tokens are routed with all-to-all communication.
\item Autoregressive decoding stores prior keys and values so each new token avoids recomputing attention over the entire prefix.
\end{enumerate}
\noindent\textbf{Answer.} Autoregressive decoding stores prior keys and values so each new token avoids recomputing attention over the entire prefix.
\par\textit{Autoregressive decoding stores prior keys and values so each new token avoids recomputing attention over the entire prefix. Tradeoff: Decode speed improves while cache memory grows with layers, sequence, batch, and KV heads. Watch for: Capacity planning only model weights causes out-of-memory failures at long context.}
\paragraph{MCQ-1302 [hard]} What is the central engineering tradeoff of KV caching?
\begin{enumerate}[label=\Alph*.]
\item It scales model depth but creates bubbles and per-request latency.
\item Parameter capacity scales, while communication and load balance determine performance.
\item Throughput and utilization rise, while scheduling fairness and latency isolation become harder.
\item Decode speed improves while cache memory grows with layers, sequence, batch, and KV heads.
\end{enumerate}
\noindent\textbf{Answer.} Decode speed improves while cache memory grows with layers, sequence, batch, and KV heads.
\par\textit{Autoregressive decoding stores prior keys and values so each new token avoids recomputing attention over the entire prefix. Tradeoff: Decode speed improves while cache memory grows with layers, sequence, batch, and KV heads. Watch for: Capacity planning only model weights causes out-of-memory failures at long context.}
\paragraph{MCQ-1303 [hard]} Which failure mode should an interviewer expect you to identify for KV caching?
\begin{enumerate}[label=\Alph*.]
\item Capacity planning only model weights causes out-of-memory failures at long context.
\item Too few microbatches leave stages idle.
\item Optimizing total tokens per second alone can starve short requests.
\item Average expert load hides hot experts that gate step time.
\end{enumerate}
\noindent\textbf{Answer.} Capacity planning only model weights causes out-of-memory failures at long context.
\par\textit{Autoregressive decoding stores prior keys and values so each new token avoids recomputing attention over the entire prefix. Tradeoff: Decode speed improves while cache memory grows with layers, sequence, batch, and KV heads. Watch for: Capacity planning only model weights causes out-of-memory failures at long context.}
\paragraph{MCQ-1304 [medium]} A design review is specifically concerned with this risk: Capacity planning only model weights causes out-of-memory failures at long context. Which topic should be reviewed first?
\begin{enumerate}[label=\Alph*.]
\item Expert parallelism
\item Continuous batching
\item KV caching
\item Pipeline parallelism
\end{enumerate}
\noindent\textbf{Answer.} KV caching
\par\textit{Autoregressive decoding stores prior keys and values so each new token avoids recomputing attention over the entire prefix. Tradeoff: Decode speed improves while cache memory grows with layers, sequence, batch, and KV heads. Watch for: Capacity planning only model weights causes out-of-memory failures at long context.}
\paragraph{MCQ-1305 [hard]} Which design note most accurately balances the benefits and costs of KV caching?
\begin{enumerate}[label=\Alph*.]
\item It scales model depth but creates bubbles and per-request latency.
\item Throughput and utilization rise, while scheduling fairness and latency isolation become harder.
\item Parameter capacity scales, while communication and load balance determine performance.
\item Decode speed improves while cache memory grows with layers, sequence, batch, and KV heads.
\end{enumerate}
\noindent\textbf{Answer.} Decode speed improves while cache memory grows with layers, sequence, batch, and KV heads.
\par\textit{Autoregressive decoding stores prior keys and values so each new token avoids recomputing attention over the entire prefix. Tradeoff: Decode speed improves while cache memory grows with layers, sequence, batch, and KV heads. Watch for: Capacity planning only model weights causes out-of-memory failures at long context.}
\paragraph{MCQ-1306 [medium]} Which explanation would be strongest in a system-design interview about KV caching?
\begin{enumerate}[label=\Alph*.]
\item Autoregressive decoding stores prior keys and values so each new token avoids recomputing attention over the entire prefix.
\item MoE experts are distributed across devices and tokens are routed with all-to-all communication.
\item Requests join and leave a running decode batch at iteration boundaries rather than waiting for a fixed batch.
\item Pipeline parallelism assigns layer stages to devices and streams microbatches through them.
\end{enumerate}
\noindent\textbf{Answer.} Autoregressive decoding stores prior keys and values so each new token avoids recomputing attention over the entire prefix.
\par\textit{Autoregressive decoding stores prior keys and values so each new token avoids recomputing attention over the entire prefix. Tradeoff: Decode speed improves while cache memory grows with layers, sequence, batch, and KV heads. Watch for: Capacity planning only model weights causes out-of-memory failures at long context.}
\paragraph{MCQ-1307 [hard]} Which production warning is most directly associated with KV caching?
\begin{enumerate}[label=\Alph*.]
\item Capacity planning only model weights causes out-of-memory failures at long context.
\item Too few microbatches leave stages idle.
\item Optimizing total tokens per second alone can starve short requests.
\item Average expert load hides hot experts that gate step time.
\end{enumerate}
\noindent\textbf{Answer.} Capacity planning only model weights causes out-of-memory failures at long context.
\par\textit{Autoregressive decoding stores prior keys and values so each new token avoids recomputing attention over the entire prefix. Tradeoff: Decode speed improves while cache memory grows with layers, sequence, batch, and KV heads. Watch for: Capacity planning only model weights causes out-of-memory failures at long context.}
\paragraph{FC-1308 [medium]} Define KV caching in interview-ready language.
\noindent\textbf{Answer.} Autoregressive decoding stores prior keys and values so each new token avoids recomputing attention over the entire prefix.
\par\textit{Autoregressive decoding stores prior keys and values so each new token avoids recomputing attention over the entire prefix.}
\paragraph{FC-1309 [hard]} State the main tradeoff and production pitfall for KV caching.
\noindent\textbf{Answer.} Tradeoff: Decode speed improves while cache memory grows with layers, sequence, batch, and KV heads. Pitfall: Capacity planning only model weights causes out-of-memory failures at long context.
\par\textit{Tradeoff: Decode speed improves while cache memory grows with layers, sequence, batch, and KV heads. Pitfall: Capacity planning only model weights causes out-of-memory failures at long context.}
\paragraph{LA-1310 [hard]} Explain KV caching from first principles, then show how you would evaluate and operate it in production.
\noindent\textbf{Answer.} Start with the mechanism: Autoregressive decoding stores prior keys and values so each new token avoids recomputing attention over the entire prefix. Make the tradeoff explicit: Decode speed improves while cache memory grows with layers, sequence, batch, and KV heads. Define workload-specific offline metrics and a latency/cost budget, test important slices, instrument the critical path, and create a rollback criterion. The production review must explicitly guard against this failure: Capacity planning only model weights causes out-of-memory failures at long context.
\par\textit{A strong answer connects mechanism, measurable tradeoffs, failure isolation, observability, and rollback rather than listing product features.}
\paragraph{MCQ-1311 [medium]} Which statement best describes PagedAttention?
\begin{enumerate}[label=\Alph*.]
\item FlashAttention tiles exact attention to reduce high-bandwidth-memory traffic and materialization of the score matrix.
\item Tensor parallelism shards matrix operations across devices and communicates partial results within layers.
\item PagedAttention manages KV cache in non-contiguous blocks using paging-inspired allocation to reduce fragmentation and enable sharing.
\item Replicas serve independent requests and scale throughput with routing and load balancing.
\end{enumerate}
\noindent\textbf{Answer.} PagedAttention manages KV cache in non-contiguous blocks using paging-inspired allocation to reduce fragmentation and enable sharing.
\par\textit{PagedAttention manages KV cache in non-contiguous blocks using paging-inspired allocation to reduce fragmentation and enable sharing. Tradeoff: Utilization and batching improve, while block tables and kernels add runtime complexity. Watch for: Oversized blocks waste tail space; tiny blocks add metadata overhead.}
\paragraph{MCQ-1312 [hard]} What is the central engineering tradeoff of PagedAttention?
\begin{enumerate}[label=\Alph*.]
\item It is operationally simple but duplicates model weights on every replica.
\item It enables large models and lower per-request latency but requires high-bandwidth interconnects.
\item Utilization and batching improve, while block tables and kernels add runtime complexity.
\item Wall-clock speed and memory improve without approximate attention, but kernels depend on shapes and hardware.
\end{enumerate}
\noindent\textbf{Answer.} Utilization and batching improve, while block tables and kernels add runtime complexity.
\par\textit{PagedAttention manages KV cache in non-contiguous blocks using paging-inspired allocation to reduce fragmentation and enable sharing. Tradeoff: Utilization and batching improve, while block tables and kernels add runtime complexity. Watch for: Oversized blocks waste tail space; tiny blocks add metadata overhead.}
\paragraph{MCQ-1313 [hard]} Which failure mode should an interviewer expect you to identify for PagedAttention?
\begin{enumerate}[label=\Alph*.]
\item Random routing ignores prefix-cache locality and uneven sequence lengths.
\item Oversized blocks waste tail space; tiny blocks add metadata overhead.
\item Calling it subquadratic compute confuses IO optimization with arithmetic complexity.
\item Spanning slow network links can make communication dominate compute.
\end{enumerate}
\noindent\textbf{Answer.} Oversized blocks waste tail space; tiny blocks add metadata overhead.
\par\textit{PagedAttention manages KV cache in non-contiguous blocks using paging-inspired allocation to reduce fragmentation and enable sharing. Tradeoff: Utilization and batching improve, while block tables and kernels add runtime complexity. Watch for: Oversized blocks waste tail space; tiny blocks add metadata overhead.}
\paragraph{MCQ-1314 [medium]} A design review is specifically concerned with this risk: Oversized blocks waste tail space; tiny blocks add metadata overhead. Which topic should be reviewed first?
\begin{enumerate}[label=\Alph*.]
\item FlashAttention
\item Data parallel inference
\item PagedAttention
\item Tensor parallelism
\end{enumerate}
\noindent\textbf{Answer.} PagedAttention
\par\textit{PagedAttention manages KV cache in non-contiguous blocks using paging-inspired allocation to reduce fragmentation and enable sharing. Tradeoff: Utilization and batching improve, while block tables and kernels add runtime complexity. Watch for: Oversized blocks waste tail space; tiny blocks add metadata overhead.}
\paragraph{MCQ-1315 [hard]} Which design note most accurately balances the benefits and costs of PagedAttention?
\begin{enumerate}[label=\Alph*.]
\item Wall-clock speed and memory improve without approximate attention, but kernels depend on shapes and hardware.
\item It is operationally simple but duplicates model weights on every replica.
\item Utilization and batching improve, while block tables and kernels add runtime complexity.
\item It enables large models and lower per-request latency but requires high-bandwidth interconnects.
\end{enumerate}
\noindent\textbf{Answer.} Utilization and batching improve, while block tables and kernels add runtime complexity.
\par\textit{PagedAttention manages KV cache in non-contiguous blocks using paging-inspired allocation to reduce fragmentation and enable sharing. Tradeoff: Utilization and batching improve, while block tables and kernels add runtime complexity. Watch for: Oversized blocks waste tail space; tiny blocks add metadata overhead.}
\paragraph{MCQ-1316 [medium]} Which explanation would be strongest in a system-design interview about PagedAttention?
\begin{enumerate}[label=\Alph*.]
\item PagedAttention manages KV cache in non-contiguous blocks using paging-inspired allocation to reduce fragmentation and enable sharing.
\item Tensor parallelism shards matrix operations across devices and communicates partial results within layers.
\item Replicas serve independent requests and scale throughput with routing and load balancing.
\item FlashAttention tiles exact attention to reduce high-bandwidth-memory traffic and materialization of the score matrix.
\end{enumerate}
\noindent\textbf{Answer.} PagedAttention manages KV cache in non-contiguous blocks using paging-inspired allocation to reduce fragmentation and enable sharing.
\par\textit{PagedAttention manages KV cache in non-contiguous blocks using paging-inspired allocation to reduce fragmentation and enable sharing. Tradeoff: Utilization and batching improve, while block tables and kernels add runtime complexity. Watch for: Oversized blocks waste tail space; tiny blocks add metadata overhead.}
\paragraph{MCQ-1317 [hard]} Which production warning is most directly associated with PagedAttention?
\begin{enumerate}[label=\Alph*.]
\item Spanning slow network links can make communication dominate compute.
\item Calling it subquadratic compute confuses IO optimization with arithmetic complexity.
\item Random routing ignores prefix-cache locality and uneven sequence lengths.
\item Oversized blocks waste tail space; tiny blocks add metadata overhead.
\end{enumerate}
\noindent\textbf{Answer.} Oversized blocks waste tail space; tiny blocks add metadata overhead.
\par\textit{PagedAttention manages KV cache in non-contiguous blocks using paging-inspired allocation to reduce fragmentation and enable sharing. Tradeoff: Utilization and batching improve, while block tables and kernels add runtime complexity. Watch for: Oversized blocks waste tail space; tiny blocks add metadata overhead.}
\paragraph{FC-1318 [medium]} Define PagedAttention in interview-ready language.
\noindent\textbf{Answer.} PagedAttention manages KV cache in non-contiguous blocks using paging-inspired allocation to reduce fragmentation and enable sharing.
\par\textit{PagedAttention manages KV cache in non-contiguous blocks using paging-inspired allocation to reduce fragmentation and enable sharing.}
\paragraph{FC-1319 [hard]} State the main tradeoff and production pitfall for PagedAttention.
\noindent\textbf{Answer.} Tradeoff: Utilization and batching improve, while block tables and kernels add runtime complexity. Pitfall: Oversized blocks waste tail space; tiny blocks add metadata overhead.
\par\textit{Tradeoff: Utilization and batching improve, while block tables and kernels add runtime complexity. Pitfall: Oversized blocks waste tail space; tiny blocks add metadata overhead.}
\paragraph{LA-1320 [hard]} Explain PagedAttention from first principles, then show how you would evaluate and operate it in production.
\noindent\textbf{Answer.} Start with the mechanism: PagedAttention manages KV cache in non-contiguous blocks using paging-inspired allocation to reduce fragmentation and enable sharing. Make the tradeoff explicit: Utilization and batching improve, while block tables and kernels add runtime complexity. Define workload-specific offline metrics and a latency/cost budget, test important slices, instrument the critical path, and create a rollback criterion. The production review must explicitly guard against this failure: Oversized blocks waste tail space; tiny blocks add metadata overhead.
\par\textit{A strong answer connects mechanism, measurable tradeoffs, failure isolation, observability, and rollback rather than listing product features.}
\paragraph{MCQ-1321 [medium]} Which statement best describes Continuous batching?
\begin{enumerate}[label=\Alph*.]
\item Requests join and leave a running decode batch at iteration boundaries rather than waiting for a fixed batch.
\item PagedAttention manages KV cache in non-contiguous blocks using paging-inspired allocation to reduce fragmentation and enable sharing.
\item Autoregressive decoding stores prior keys and values so each new token avoids recomputing attention over the entire prefix.
\item FlashAttention tiles exact attention to reduce high-bandwidth-memory traffic and materialization of the score matrix.
\end{enumerate}
\noindent\textbf{Answer.} Requests join and leave a running decode batch at iteration boundaries rather than waiting for a fixed batch.
\par\textit{Requests join and leave a running decode batch at iteration boundaries rather than waiting for a fixed batch. Tradeoff: Throughput and utilization rise, while scheduling fairness and latency isolation become harder. Watch for: Optimizing total tokens per second alone can starve short requests.}
\paragraph{MCQ-1322 [hard]} What is the central engineering tradeoff of Continuous batching?
\begin{enumerate}[label=\Alph*.]
\item Utilization and batching improve, while block tables and kernels add runtime complexity.
\item Throughput and utilization rise, while scheduling fairness and latency isolation become harder.
\item Decode speed improves while cache memory grows with layers, sequence, batch, and KV heads.
\item Wall-clock speed and memory improve without approximate attention, but kernels depend on shapes and hardware.
\end{enumerate}
\noindent\textbf{Answer.} Throughput and utilization rise, while scheduling fairness and latency isolation become harder.
\par\textit{Requests join and leave a running decode batch at iteration boundaries rather than waiting for a fixed batch. Tradeoff: Throughput and utilization rise, while scheduling fairness and latency isolation become harder. Watch for: Optimizing total tokens per second alone can starve short requests.}
\paragraph{MCQ-1323 [hard]} Which failure mode should an interviewer expect you to identify for Continuous batching?
\begin{enumerate}[label=\Alph*.]
\item Oversized blocks waste tail space; tiny blocks add metadata overhead.
\item Calling it subquadratic compute confuses IO optimization with arithmetic complexity.
\item Capacity planning only model weights causes out-of-memory failures at long context.
\item Optimizing total tokens per second alone can starve short requests.
\end{enumerate}
\noindent\textbf{Answer.} Optimizing total tokens per second alone can starve short requests.
\par\textit{Requests join and leave a running decode batch at iteration boundaries rather than waiting for a fixed batch. Tradeoff: Throughput and utilization rise, while scheduling fairness and latency isolation become harder. Watch for: Optimizing total tokens per second alone can starve short requests.}
\paragraph{MCQ-1324 [medium]} A design review is specifically concerned with this risk: Optimizing total tokens per second alone can starve short requests. Which topic should be reviewed first?
\begin{enumerate}[label=\Alph*.]
\item PagedAttention
\item KV caching
\item Continuous batching
\item FlashAttention
\end{enumerate}
\noindent\textbf{Answer.} Continuous batching
\par\textit{Requests join and leave a running decode batch at iteration boundaries rather than waiting for a fixed batch. Tradeoff: Throughput and utilization rise, while scheduling fairness and latency isolation become harder. Watch for: Optimizing total tokens per second alone can starve short requests.}
\paragraph{MCQ-1325 [hard]} Which design note most accurately balances the benefits and costs of Continuous batching?
\begin{enumerate}[label=\Alph*.]
\item Wall-clock speed and memory improve without approximate attention, but kernels depend on shapes and hardware.
\item Throughput and utilization rise, while scheduling fairness and latency isolation become harder.
\item Decode speed improves while cache memory grows with layers, sequence, batch, and KV heads.
\item Utilization and batching improve, while block tables and kernels add runtime complexity.
\end{enumerate}
\noindent\textbf{Answer.} Throughput and utilization rise, while scheduling fairness and latency isolation become harder.
\par\textit{Requests join and leave a running decode batch at iteration boundaries rather than waiting for a fixed batch. Tradeoff: Throughput and utilization rise, while scheduling fairness and latency isolation become harder. Watch for: Optimizing total tokens per second alone can starve short requests.}
\paragraph{MCQ-1326 [medium]} Which explanation would be strongest in a system-design interview about Continuous batching?
\begin{enumerate}[label=\Alph*.]
\item Requests join and leave a running decode batch at iteration boundaries rather than waiting for a fixed batch.
\item Autoregressive decoding stores prior keys and values so each new token avoids recomputing attention over the entire prefix.
\item PagedAttention manages KV cache in non-contiguous blocks using paging-inspired allocation to reduce fragmentation and enable sharing.
\item FlashAttention tiles exact attention to reduce high-bandwidth-memory traffic and materialization of the score matrix.
\end{enumerate}
\noindent\textbf{Answer.} Requests join and leave a running decode batch at iteration boundaries rather than waiting for a fixed batch.
\par\textit{Requests join and leave a running decode batch at iteration boundaries rather than waiting for a fixed batch. Tradeoff: Throughput and utilization rise, while scheduling fairness and latency isolation become harder. Watch for: Optimizing total tokens per second alone can starve short requests.}
\paragraph{MCQ-1327 [hard]} Which production warning is most directly associated with Continuous batching?
\begin{enumerate}[label=\Alph*.]
\item Optimizing total tokens per second alone can starve short requests.
\item Oversized blocks waste tail space; tiny blocks add metadata overhead.
\item Calling it subquadratic compute confuses IO optimization with arithmetic complexity.
\item Capacity planning only model weights causes out-of-memory failures at long context.
\end{enumerate}
\noindent\textbf{Answer.} Optimizing total tokens per second alone can starve short requests.
\par\textit{Requests join and leave a running decode batch at iteration boundaries rather than waiting for a fixed batch. Tradeoff: Throughput and utilization rise, while scheduling fairness and latency isolation become harder. Watch for: Optimizing total tokens per second alone can starve short requests.}
\paragraph{FC-1328 [medium]} Define Continuous batching in interview-ready language.
\noindent\textbf{Answer.} Requests join and leave a running decode batch at iteration boundaries rather than waiting for a fixed batch.
\par\textit{Requests join and leave a running decode batch at iteration boundaries rather than waiting for a fixed batch.}
\paragraph{FC-1329 [hard]} State the main tradeoff and production pitfall for Continuous batching.
\noindent\textbf{Answer.} Tradeoff: Throughput and utilization rise, while scheduling fairness and latency isolation become harder. Pitfall: Optimizing total tokens per second alone can starve short requests.
\par\textit{Tradeoff: Throughput and utilization rise, while scheduling fairness and latency isolation become harder. Pitfall: Optimizing total tokens per second alone can starve short requests.}
\paragraph{LA-1330 [hard]} Explain Continuous batching from first principles, then show how you would evaluate and operate it in production.
\noindent\textbf{Answer.} Start with the mechanism: Requests join and leave a running decode batch at iteration boundaries rather than waiting for a fixed batch. Make the tradeoff explicit: Throughput and utilization rise, while scheduling fairness and latency isolation become harder. Define workload-specific offline metrics and a latency/cost budget, test important slices, instrument the critical path, and create a rollback criterion. The production review must explicitly guard against this failure: Optimizing total tokens per second alone can starve short requests.
\par\textit{A strong answer connects mechanism, measurable tradeoffs, failure isolation, observability, and rollback rather than listing product features.}
\paragraph{MCQ-1331 [medium]} Which statement best describes Prefix caching?
\begin{enumerate}[label=\Alph*.]
\item Grouped-query and multi-query attention share key-value heads across query heads to shrink KV cache and memory bandwidth.
\item Tensor parallelism shards matrix operations across devices and communicates partial results within layers.
\item Shared prompt prefixes reuse precomputed KV blocks across requests with matching tokens and model state.
\item A cheaper draft proposes tokens and a target model verifies them while preserving the target distribution.
\end{enumerate}
\noindent\textbf{Answer.} Shared prompt prefixes reuse precomputed KV blocks across requests with matching tokens and model state.
\par\textit{Shared prompt prefixes reuse precomputed KV blocks across requests with matching tokens and model state. Tradeoff: Repeated system prompts become cheaper, but hit rate, invalidation, and tenant isolation matter. Watch for: Cache keys that omit adapters or model revision return incorrect state.}
\paragraph{MCQ-1332 [hard]} What is the central engineering tradeoff of Prefix caching?
\begin{enumerate}[label=\Alph*.]
\item Decode efficiency improves, with possible quality loss relative to full multi-head KV.
\item Repeated system prompts become cheaper, but hit rate, invalidation, and tenant isolation matter.
\item It enables large models and lower per-request latency but requires high-bandwidth interconnects.
\item Latency falls when acceptance is high, but drafting and verification overhead can dominate.
\end{enumerate}
\noindent\textbf{Answer.} Repeated system prompts become cheaper, but hit rate, invalidation, and tenant isolation matter.
\par\textit{Shared prompt prefixes reuse precomputed KV blocks across requests with matching tokens and model state. Tradeoff: Repeated system prompts become cheaper, but hit rate, invalidation, and tenant isolation matter. Watch for: Cache keys that omit adapters or model revision return incorrect state.}
\paragraph{MCQ-1333 [hard]} Which failure mode should an interviewer expect you to identify for Prefix caching?
\begin{enumerate}[label=\Alph*.]
\item Converting weights naively after training does not reproduce a model trained for shared KV heads.
\item Spanning slow network links can make communication dominate compute.
\item Cache keys that omit adapters or model revision return incorrect state.
\item Measuring draft speed without acceptance rate gives meaningless projections.
\end{enumerate}
\noindent\textbf{Answer.} Cache keys that omit adapters or model revision return incorrect state.
\par\textit{Shared prompt prefixes reuse precomputed KV blocks across requests with matching tokens and model state. Tradeoff: Repeated system prompts become cheaper, but hit rate, invalidation, and tenant isolation matter. Watch for: Cache keys that omit adapters or model revision return incorrect state.}
\paragraph{MCQ-1334 [medium]} A design review is specifically concerned with this risk: Cache keys that omit adapters or model revision return incorrect state. Which topic should be reviewed first?
\begin{enumerate}[label=\Alph*.]
\item Prefix caching
\item Tensor parallelism
\item Speculative decoding
\item GQA and MQA
\end{enumerate}
\noindent\textbf{Answer.} Prefix caching
\par\textit{Shared prompt prefixes reuse precomputed KV blocks across requests with matching tokens and model state. Tradeoff: Repeated system prompts become cheaper, but hit rate, invalidation, and tenant isolation matter. Watch for: Cache keys that omit adapters or model revision return incorrect state.}
\paragraph{MCQ-1335 [hard]} Which design note most accurately balances the benefits and costs of Prefix caching?
\begin{enumerate}[label=\Alph*.]
\item It enables large models and lower per-request latency but requires high-bandwidth interconnects.
\item Decode efficiency improves, with possible quality loss relative to full multi-head KV.
\item Repeated system prompts become cheaper, but hit rate, invalidation, and tenant isolation matter.
\item Latency falls when acceptance is high, but drafting and verification overhead can dominate.
\end{enumerate}
\noindent\textbf{Answer.} Repeated system prompts become cheaper, but hit rate, invalidation, and tenant isolation matter.
\par\textit{Shared prompt prefixes reuse precomputed KV blocks across requests with matching tokens and model state. Tradeoff: Repeated system prompts become cheaper, but hit rate, invalidation, and tenant isolation matter. Watch for: Cache keys that omit adapters or model revision return incorrect state.}
\paragraph{MCQ-1336 [medium]} Which explanation would be strongest in a system-design interview about Prefix caching?
\begin{enumerate}[label=\Alph*.]
\item Shared prompt prefixes reuse precomputed KV blocks across requests with matching tokens and model state.
\item Tensor parallelism shards matrix operations across devices and communicates partial results within layers.
\item A cheaper draft proposes tokens and a target model verifies them while preserving the target distribution.
\item Grouped-query and multi-query attention share key-value heads across query heads to shrink KV cache and memory bandwidth.
\end{enumerate}
\noindent\textbf{Answer.} Shared prompt prefixes reuse precomputed KV blocks across requests with matching tokens and model state.
\par\textit{Shared prompt prefixes reuse precomputed KV blocks across requests with matching tokens and model state. Tradeoff: Repeated system prompts become cheaper, but hit rate, invalidation, and tenant isolation matter. Watch for: Cache keys that omit adapters or model revision return incorrect state.}
\paragraph{MCQ-1337 [hard]} Which production warning is most directly associated with Prefix caching?
\begin{enumerate}[label=\Alph*.]
\item Spanning slow network links can make communication dominate compute.
\item Cache keys that omit adapters or model revision return incorrect state.
\item Converting weights naively after training does not reproduce a model trained for shared KV heads.
\item Measuring draft speed without acceptance rate gives meaningless projections.
\end{enumerate}
\noindent\textbf{Answer.} Cache keys that omit adapters or model revision return incorrect state.
\par\textit{Shared prompt prefixes reuse precomputed KV blocks across requests with matching tokens and model state. Tradeoff: Repeated system prompts become cheaper, but hit rate, invalidation, and tenant isolation matter. Watch for: Cache keys that omit adapters or model revision return incorrect state.}
\paragraph{FC-1338 [medium]} Define Prefix caching in interview-ready language.
\noindent\textbf{Answer.} Shared prompt prefixes reuse precomputed KV blocks across requests with matching tokens and model state.
\par\textit{Shared prompt prefixes reuse precomputed KV blocks across requests with matching tokens and model state.}
\paragraph{FC-1339 [hard]} State the main tradeoff and production pitfall for Prefix caching.
\noindent\textbf{Answer.} Tradeoff: Repeated system prompts become cheaper, but hit rate, invalidation, and tenant isolation matter. Pitfall: Cache keys that omit adapters or model revision return incorrect state.
\par\textit{Tradeoff: Repeated system prompts become cheaper, but hit rate, invalidation, and tenant isolation matter. Pitfall: Cache keys that omit adapters or model revision return incorrect state.}
\paragraph{LA-1340 [hard]} Explain Prefix caching from first principles, then show how you would evaluate and operate it in production.
\noindent\textbf{Answer.} Start with the mechanism: Shared prompt prefixes reuse precomputed KV blocks across requests with matching tokens and model state. Make the tradeoff explicit: Repeated system prompts become cheaper, but hit rate, invalidation, and tenant isolation matter. Define workload-specific offline metrics and a latency/cost budget, test important slices, instrument the critical path, and create a rollback criterion. The production review must explicitly guard against this failure: Cache keys that omit adapters or model revision return incorrect state.
\par\textit{A strong answer connects mechanism, measurable tradeoffs, failure isolation, observability, and rollback rather than listing product features.}
\paragraph{MCQ-1341 [medium]} Which statement best describes FlashAttention?
\begin{enumerate}[label=\Alph*.]
\item PagedAttention manages KV cache in non-contiguous blocks using paging-inspired allocation to reduce fragmentation and enable sharing.
\item FlashAttention tiles exact attention to reduce high-bandwidth-memory traffic and materialization of the score matrix.
\item MoE experts are distributed across devices and tokens are routed with all-to-all communication.
\item Requests join and leave a running decode batch at iteration boundaries rather than waiting for a fixed batch.
\end{enumerate}
\noindent\textbf{Answer.} FlashAttention tiles exact attention to reduce high-bandwidth-memory traffic and materialization of the score matrix.
\par\textit{FlashAttention tiles exact attention to reduce high-bandwidth-memory traffic and materialization of the score matrix. Tradeoff: Wall-clock speed and memory improve without approximate attention, but kernels depend on shapes and hardware. Watch for: Calling it subquadratic compute confuses IO optimization with arithmetic complexity.}
\paragraph{MCQ-1342 [hard]} What is the central engineering tradeoff of FlashAttention?
\begin{enumerate}[label=\Alph*.]
\item Throughput and utilization rise, while scheduling fairness and latency isolation become harder.
\item Wall-clock speed and memory improve without approximate attention, but kernels depend on shapes and hardware.
\item Parameter capacity scales, while communication and load balance determine performance.
\item Utilization and batching improve, while block tables and kernels add runtime complexity.
\end{enumerate}
\noindent\textbf{Answer.} Wall-clock speed and memory improve without approximate attention, but kernels depend on shapes and hardware.
\par\textit{FlashAttention tiles exact attention to reduce high-bandwidth-memory traffic and materialization of the score matrix. Tradeoff: Wall-clock speed and memory improve without approximate attention, but kernels depend on shapes and hardware. Watch for: Calling it subquadratic compute confuses IO optimization with arithmetic complexity.}
\paragraph{MCQ-1343 [hard]} Which failure mode should an interviewer expect you to identify for FlashAttention?
\begin{enumerate}[label=\Alph*.]
\item Oversized blocks waste tail space; tiny blocks add metadata overhead.
\item Average expert load hides hot experts that gate step time.
\item Optimizing total tokens per second alone can starve short requests.
\item Calling it subquadratic compute confuses IO optimization with arithmetic complexity.
\end{enumerate}
\noindent\textbf{Answer.} Calling it subquadratic compute confuses IO optimization with arithmetic complexity.
\par\textit{FlashAttention tiles exact attention to reduce high-bandwidth-memory traffic and materialization of the score matrix. Tradeoff: Wall-clock speed and memory improve without approximate attention, but kernels depend on shapes and hardware. Watch for: Calling it subquadratic compute confuses IO optimization with arithmetic complexity.}
\paragraph{MCQ-1344 [medium]} A design review is specifically concerned with this risk: Calling it subquadratic compute confuses IO optimization with arithmetic complexity. Which topic should be reviewed first?
\begin{enumerate}[label=\Alph*.]
\item Expert parallelism
\item FlashAttention
\item PagedAttention
\item Continuous batching
\end{enumerate}
\noindent\textbf{Answer.} FlashAttention
\par\textit{FlashAttention tiles exact attention to reduce high-bandwidth-memory traffic and materialization of the score matrix. Tradeoff: Wall-clock speed and memory improve without approximate attention, but kernels depend on shapes and hardware. Watch for: Calling it subquadratic compute confuses IO optimization with arithmetic complexity.}
\paragraph{MCQ-1345 [hard]} Which design note most accurately balances the benefits and costs of FlashAttention?
\begin{enumerate}[label=\Alph*.]
\item Utilization and batching improve, while block tables and kernels add runtime complexity.
\item Wall-clock speed and memory improve without approximate attention, but kernels depend on shapes and hardware.
\item Parameter capacity scales, while communication and load balance determine performance.
\item Throughput and utilization rise, while scheduling fairness and latency isolation become harder.
\end{enumerate}
\noindent\textbf{Answer.} Wall-clock speed and memory improve without approximate attention, but kernels depend on shapes and hardware.
\par\textit{FlashAttention tiles exact attention to reduce high-bandwidth-memory traffic and materialization of the score matrix. Tradeoff: Wall-clock speed and memory improve without approximate attention, but kernels depend on shapes and hardware. Watch for: Calling it subquadratic compute confuses IO optimization with arithmetic complexity.}
\paragraph{MCQ-1346 [medium]} Which explanation would be strongest in a system-design interview about FlashAttention?
\begin{enumerate}[label=\Alph*.]
\item PagedAttention manages KV cache in non-contiguous blocks using paging-inspired allocation to reduce fragmentation and enable sharing.
\item Requests join and leave a running decode batch at iteration boundaries rather than waiting for a fixed batch.
\item MoE experts are distributed across devices and tokens are routed with all-to-all communication.
\item FlashAttention tiles exact attention to reduce high-bandwidth-memory traffic and materialization of the score matrix.
\end{enumerate}
\noindent\textbf{Answer.} FlashAttention tiles exact attention to reduce high-bandwidth-memory traffic and materialization of the score matrix.
\par\textit{FlashAttention tiles exact attention to reduce high-bandwidth-memory traffic and materialization of the score matrix. Tradeoff: Wall-clock speed and memory improve without approximate attention, but kernels depend on shapes and hardware. Watch for: Calling it subquadratic compute confuses IO optimization with arithmetic complexity.}
\paragraph{MCQ-1347 [hard]} Which production warning is most directly associated with FlashAttention?
\begin{enumerate}[label=\Alph*.]
\item Optimizing total tokens per second alone can starve short requests.
\item Oversized blocks waste tail space; tiny blocks add metadata overhead.
\item Calling it subquadratic compute confuses IO optimization with arithmetic complexity.
\item Average expert load hides hot experts that gate step time.
\end{enumerate}
\noindent\textbf{Answer.} Calling it subquadratic compute confuses IO optimization with arithmetic complexity.
\par\textit{FlashAttention tiles exact attention to reduce high-bandwidth-memory traffic and materialization of the score matrix. Tradeoff: Wall-clock speed and memory improve without approximate attention, but kernels depend on shapes and hardware. Watch for: Calling it subquadratic compute confuses IO optimization with arithmetic complexity.}
\paragraph{FC-1348 [medium]} Define FlashAttention in interview-ready language.
\noindent\textbf{Answer.} FlashAttention tiles exact attention to reduce high-bandwidth-memory traffic and materialization of the score matrix.
\par\textit{FlashAttention tiles exact attention to reduce high-bandwidth-memory traffic and materialization of the score matrix.}
\paragraph{FC-1349 [hard]} State the main tradeoff and production pitfall for FlashAttention.
\noindent\textbf{Answer.} Tradeoff: Wall-clock speed and memory improve without approximate attention, but kernels depend on shapes and hardware. Pitfall: Calling it subquadratic compute confuses IO optimization with arithmetic complexity.
\par\textit{Tradeoff: Wall-clock speed and memory improve without approximate attention, but kernels depend on shapes and hardware. Pitfall: Calling it subquadratic compute confuses IO optimization with arithmetic complexity.}
\paragraph{LA-1350 [hard]} Explain FlashAttention from first principles, then show how you would evaluate and operate it in production.
\noindent\textbf{Answer.} Start with the mechanism: FlashAttention tiles exact attention to reduce high-bandwidth-memory traffic and materialization of the score matrix. Make the tradeoff explicit: Wall-clock speed and memory improve without approximate attention, but kernels depend on shapes and hardware. Define workload-specific offline metrics and a latency/cost budget, test important slices, instrument the critical path, and create a rollback criterion. The production review must explicitly guard against this failure: Calling it subquadratic compute confuses IO optimization with arithmetic complexity.
\par\textit{A strong answer connects mechanism, measurable tradeoffs, failure isolation, observability, and rollback rather than listing product features.}
\paragraph{MCQ-1351 [medium]} Which statement best describes GQA and MQA?
\begin{enumerate}[label=\Alph*.]
\item PagedAttention manages KV cache in non-contiguous blocks using paging-inspired allocation to reduce fragmentation and enable sharing.
\item Requests join and leave a running decode batch at iteration boundaries rather than waiting for a fixed batch.
\item Servers emit tokens incrementally with backpressure, disconnect handling, and final usage accounting.
\item Grouped-query and multi-query attention share key-value heads across query heads to shrink KV cache and memory bandwidth.
\end{enumerate}
\noindent\textbf{Answer.} Grouped-query and multi-query attention share key-value heads across query heads to shrink KV cache and memory bandwidth.
\par\textit{Grouped-query and multi-query attention share key-value heads across query heads to shrink KV cache and memory bandwidth. Tradeoff: Decode efficiency improves, with possible quality loss relative to full multi-head KV. Watch for: Converting weights naively after training does not reproduce a model trained for shared KV heads.}
\paragraph{MCQ-1352 [hard]} What is the central engineering tradeoff of GQA and MQA?
\begin{enumerate}[label=\Alph*.]
\item Throughput and utilization rise, while scheduling fairness and latency isolation become harder.
\item Perceived latency improves, but partial outputs complicate moderation and retries.
\item Decode efficiency improves, with possible quality loss relative to full multi-head KV.
\item Utilization and batching improve, while block tables and kernels add runtime complexity.
\end{enumerate}
\noindent\textbf{Answer.} Decode efficiency improves, with possible quality loss relative to full multi-head KV.
\par\textit{Grouped-query and multi-query attention share key-value heads across query heads to shrink KV cache and memory bandwidth. Tradeoff: Decode efficiency improves, with possible quality loss relative to full multi-head KV. Watch for: Converting weights naively after training does not reproduce a model trained for shared KV heads.}
\paragraph{MCQ-1353 [hard]} Which failure mode should an interviewer expect you to identify for GQA and MQA?
\begin{enumerate}[label=\Alph*.]
\item Oversized blocks waste tail space; tiny blocks add metadata overhead.
\item Converting weights naively after training does not reproduce a model trained for shared KV heads.
\item Retrying after visible partial output can duplicate text or side effects.
\item Optimizing total tokens per second alone can starve short requests.
\end{enumerate}
\noindent\textbf{Answer.} Converting weights naively after training does not reproduce a model trained for shared KV heads.
\par\textit{Grouped-query and multi-query attention share key-value heads across query heads to shrink KV cache and memory bandwidth. Tradeoff: Decode efficiency improves, with possible quality loss relative to full multi-head KV. Watch for: Converting weights naively after training does not reproduce a model trained for shared KV heads.}
\paragraph{MCQ-1354 [medium]} A design review is specifically concerned with this risk: Converting weights naively after training does not reproduce a model trained for shared KV heads. Which topic should be reviewed first?
\begin{enumerate}[label=\Alph*.]
\item PagedAttention
\item Streaming generation
\item GQA and MQA
\item Continuous batching
\end{enumerate}
\noindent\textbf{Answer.} GQA and MQA
\par\textit{Grouped-query and multi-query attention share key-value heads across query heads to shrink KV cache and memory bandwidth. Tradeoff: Decode efficiency improves, with possible quality loss relative to full multi-head KV. Watch for: Converting weights naively after training does not reproduce a model trained for shared KV heads.}
\paragraph{MCQ-1355 [hard]} Which design note most accurately balances the benefits and costs of GQA and MQA?
\begin{enumerate}[label=\Alph*.]
\item Throughput and utilization rise, while scheduling fairness and latency isolation become harder.
\item Perceived latency improves, but partial outputs complicate moderation and retries.
\item Utilization and batching improve, while block tables and kernels add runtime complexity.
\item Decode efficiency improves, with possible quality loss relative to full multi-head KV.
\end{enumerate}
\noindent\textbf{Answer.} Decode efficiency improves, with possible quality loss relative to full multi-head KV.
\par\textit{Grouped-query and multi-query attention share key-value heads across query heads to shrink KV cache and memory bandwidth. Tradeoff: Decode efficiency improves, with possible quality loss relative to full multi-head KV. Watch for: Converting weights naively after training does not reproduce a model trained for shared KV heads.}
\paragraph{MCQ-1356 [medium]} Which explanation would be strongest in a system-design interview about GQA and MQA?
\begin{enumerate}[label=\Alph*.]
\item Requests join and leave a running decode batch at iteration boundaries rather than waiting for a fixed batch.
\item PagedAttention manages KV cache in non-contiguous blocks using paging-inspired allocation to reduce fragmentation and enable sharing.
\item Grouped-query and multi-query attention share key-value heads across query heads to shrink KV cache and memory bandwidth.
\item Servers emit tokens incrementally with backpressure, disconnect handling, and final usage accounting.
\end{enumerate}
\noindent\textbf{Answer.} Grouped-query and multi-query attention share key-value heads across query heads to shrink KV cache and memory bandwidth.
\par\textit{Grouped-query and multi-query attention share key-value heads across query heads to shrink KV cache and memory bandwidth. Tradeoff: Decode efficiency improves, with possible quality loss relative to full multi-head KV. Watch for: Converting weights naively after training does not reproduce a model trained for shared KV heads.}
\paragraph{MCQ-1357 [hard]} Which production warning is most directly associated with GQA and MQA?
\begin{enumerate}[label=\Alph*.]
\item Optimizing total tokens per second alone can starve short requests.
\item Retrying after visible partial output can duplicate text or side effects.
\item Oversized blocks waste tail space; tiny blocks add metadata overhead.
\item Converting weights naively after training does not reproduce a model trained for shared KV heads.
\end{enumerate}
\noindent\textbf{Answer.} Converting weights naively after training does not reproduce a model trained for shared KV heads.
\par\textit{Grouped-query and multi-query attention share key-value heads across query heads to shrink KV cache and memory bandwidth. Tradeoff: Decode efficiency improves, with possible quality loss relative to full multi-head KV. Watch for: Converting weights naively after training does not reproduce a model trained for shared KV heads.}
\paragraph{FC-1358 [medium]} Define GQA and MQA in interview-ready language.
\noindent\textbf{Answer.} Grouped-query and multi-query attention share key-value heads across query heads to shrink KV cache and memory bandwidth.
\par\textit{Grouped-query and multi-query attention share key-value heads across query heads to shrink KV cache and memory bandwidth.}
\paragraph{FC-1359 [hard]} State the main tradeoff and production pitfall for GQA and MQA.
\noindent\textbf{Answer.} Tradeoff: Decode efficiency improves, with possible quality loss relative to full multi-head KV. Pitfall: Converting weights naively after training does not reproduce a model trained for shared KV heads.
\par\textit{Tradeoff: Decode efficiency improves, with possible quality loss relative to full multi-head KV. Pitfall: Converting weights naively after training does not reproduce a model trained for shared KV heads.}
\paragraph{LA-1360 [hard]} Explain GQA and MQA from first principles, then show how you would evaluate and operate it in production.
\noindent\textbf{Answer.} Start with the mechanism: Grouped-query and multi-query attention share key-value heads across query heads to shrink KV cache and memory bandwidth. Make the tradeoff explicit: Decode efficiency improves, with possible quality loss relative to full multi-head KV. Define workload-specific offline metrics and a latency/cost budget, test important slices, instrument the critical path, and create a rollback criterion. The production review must explicitly guard against this failure: Converting weights naively after training does not reproduce a model trained for shared KV heads.
\par\textit{A strong answer connects mechanism, measurable tradeoffs, failure isolation, observability, and rollback rather than listing product features.}
\paragraph{MCQ-1361 [medium]} Which statement best describes Speculative decoding?
\begin{enumerate}[label=\Alph*.]
\item A cheaper draft proposes tokens and a target model verifies them while preserving the target distribution.
\item Autoregressive decoding stores prior keys and values so each new token avoids recomputing attention over the entire prefix.
\item Requests join and leave a running decode batch at iteration boundaries rather than waiting for a fixed batch.
\item Replicas serve independent requests and scale throughput with routing and load balancing.
\end{enumerate}
\noindent\textbf{Answer.} A cheaper draft proposes tokens and a target model verifies them while preserving the target distribution.
